% =====================================================================
%  第 6 章  随机微分方程与扩散过程
% =====================================================================
\chapter{随机微分方程与扩散过程}
\label{ch:6}

第~5~章开头提出的问题, 现在可以回答了. 一个受白噪声扰动的动力系统
\[
\frac{\dd X_t}{\dd t}=b(t,X_t)+\sigma(t,X_t)\,\text{``}\frac{\dd B_t}{\dd t}\text{''}
\]
应当理解为积分方程 $X_t=X_0+\int_0^tb(s,X_s)\dd s+\int_0^t\sigma(s,X_s)\dd B_s$, 其中第二个积分是 It\^o 积分. 这样的方程称为\textbf{随机微分方程} (SDE). 它的解是一类特别重要的马尔可夫过程, 称为\textbf{扩散过程}: 局部地看, 它在 $(t,t+\dd t]$ 内的增量近似服从均值为 $b\dd t$、方差为 $\sigma^2\dd t$ 的正态分布.

本章是全书的综合. 存在唯一性定理的证明结合了 Doob 不等式 (第~3~章) 与 It\^o 等距 (第~5~章), 马尔可夫性的推导依赖代入引理与单调类定理 (第~2,~4~章). 生成元与 Kolmogorov 方程推广了布朗运动与热方程的联系 (第~4~章), 而 It\^o 公式与可选停时定理又为 Feynman-Kac 公式及 Dirichlet 问题提供了工具. 期权定价除了使用 Girsanov 定理与鞅表示定理, 还要用到反射原理和首达时的分布.

本章的结构如下. 6.1 节给出定义并求解线性方程; 6.2 节证明 Lipschitz 条件下强解的存在唯一性, 并讨论弱解; 6.3 节研究解的马尔可夫性与生成元, 推导 Kolmogorov 方程并求解一维扩散的出口问题; 6.4 节建立随机微分方程与偏微分方程的联系, 并以 Black-Scholes 理论、障碍期权与美式期权为应用; 6.5 节介绍数值方法. 为了使记号简洁, 理论部分只讨论一维方程 ($X$ 与 $B$ 都是一维的); 多维的情形只需把绝对值换成范数, 证明逐字相同. 本章中信息流 $\{\mathcal{F}_t\}$ 满足通常条件, $B$ 是 $\{\mathcal{F}_t\}$-布朗运动.

% =====================================================================
\section{随机微分方程}
\label{sec:6-1}

% ---------------------------------------------------------------------
\subsection{强解与弱解}

设 $b,\sigma:[0,T]\times\R\to\R$ 是 Borel 函数, 分别称为\textbf{漂移系数}与\textbf{扩散系数}. 考虑方程
\begin{equation}\label{eq:sde}
\dd X_t=b(t,X_t)\dd t+\sigma(t,X_t)\dd B_t,\qquad X_0=\xi.
\end{equation}

\begin{definition}[强解]\label{def:strong-solution}
设 $\xi$ 是 $\mathcal{F}_0$-可测的随机变量. 称轨道连续的适应过程 $X=\{X_t\}_{t\in[0,T]}$ 为方程 \eqref{eq:sde} 的\textbf{强解}, 如果 $\int_0^T\big(\abs{b(s,X_s)}+\sigma(s,X_s)^2\big)\dd s<\infty$ $\as$, 且几乎必然地, 对一切 $t\in[0,T]$,
\begin{equation}\label{eq:sde-integral}
X_t=\xi+\int_0^tb(s,X_s)\dd s+\int_0^t\sigma(s,X_s)\dd B_s.
\end{equation}
若方程的任意两个强解 (对同一个 $B$ 与 $\xi$) 都不可区分, 称方程有\textbf{轨道唯一性}.
\end{definition}

条件 $\int_0^T(\abs{b}+\sigma^2)\dd s<\infty$ 保证 \eqref{eq:sde-integral} 右端有意义: $s\mapsto\sigma(s,X_s)$ 循序可测 (Borel 函数与循序可测过程的复合) 且属于 $\mathscr{L}^2_{\mathrm{loc}}$. 所以强解是一个 It\^o 过程 (定义~\ref{def:ito-process}), 系数为 $b(t,X_t)$ 与 $\sigma(t,X_t)$. ``强'' 是指解在\emph{给定的}概率空间上、相对于\emph{给定的}布朗运动构造的, 它是 ``输入'' $(\xi,B)$ 的函数. 有时我们只关心解的分布, 此时可以放宽要求.

\begin{definition}[弱解]\label{def:weak-solution}
称方程 \eqref{eq:sde} (初始分布为 $\mu$) 有\textbf{弱解}, 如果存在一个概率空间 $(\Omega,\mathcal{F},\P)$、其上满足通常条件的信息流 $\{\mathcal{F}_t\}$、一个 $\{\mathcal{F}_t\}$-布朗运动 $W$ 与一个轨道连续的适应过程 $X$, 使得 $X_0\sim\mu$, 且 \eqref{eq:sde-integral} 在 $B$ 换成 $W$ 后成立. 若任意两个弱解 (可能在不同的概率空间上) 的 $X$ 有相同的有限维分布, 称方程有\textbf{分布唯一性}.
\end{definition}

强解当然是弱解. 二者的区别在于, 弱解中的布朗运动 $W$ 是和解 $X$ 一起被构造出来的, $X$ 不必是 $W$ 的函数. 6.2.2 小节将给出一个有弱解而没有强解的例子.

% ---------------------------------------------------------------------
\subsection{线性方程}

当 $b$ 与 $\sigma$ 关于 $x$ 是线性函数时, 方程可以显式求解. 方法与常微分方程的 ``常数变易法'' 相同, 只是要用 It\^o 公式代替普通的链式法则.

\begin{proposition}[一维线性方程]\label{prop:linear-sde}
设 $a_1,a_2,\sigma_1,\sigma_2$ 是常数, $x\in\R$. 方程
\begin{equation}\label{eq:linear-sde}
\dd X_t=(a_1X_t+a_2)\dd t+(\sigma_1X_t+\sigma_2)\dd B_t,\qquad X_0=x
\end{equation}
的解为
\begin{align*}
X_t&=\Phi_t\Big(x+(a_2-\sigma_1\sigma_2)\int_0^t\Phi_s^{-1}\dd s+\sigma_2\int_0^t\Phi_s^{-1}\dd B_s\Big),\\
\Phi_t&=\exp\Big(\big(a_1-\tfrac{1}{2}\sigma_1^2\big)t+\sigma_1B_t\Big).
\end{align*}
\end{proposition}

\begin{proof}
由 It\^o 公式 (例~\ref{ex:gbm-ou} (1)), $\Phi$ 满足齐次方程 $\dd\Phi_t=a_1\Phi_t\dd t+\sigma_1\Phi_t\dd B_t$, $\Phi_0=1$. 对 $f(x)=x^{-1}$ 应用 It\^o 公式 ($f'=-x^{-2}$, $f''=2x^{-3}$, $\dd\ip{\Phi}_t=\sigma_1^2\Phi_t^2\dd t$),
\[
\dd\Phi_t^{-1}=-\Phi_t^{-2}\dd\Phi_t+\Phi_t^{-3}\dd\ip{\Phi}_t=\Phi_t^{-1}\big((\sigma_1^2-a_1)\dd t-\sigma_1\dd B_t\big).
\]
记 $Y_t$ 为括号中的过程, 则 $\dd Y_t=(a_2-\sigma_1\sigma_2)\Phi_t^{-1}\dd t+\sigma_2\Phi_t^{-1}\dd B_t$, $\ip{\Phi,Y}$ 的密度为 $\sigma_1\Phi_t\cdot\sigma_2\Phi_t^{-1}=\sigma_1\sigma_2$. 由乘积法则 (推论~\ref{cor:product-rule}),
\begin{align*}
\dd(\Phi_tY_t)&=Y_t\dd\Phi_t+\Phi_t\dd Y_t+\dd\ip{\Phi,Y}_t\\
&=a_1\Phi_tY_t\dd t+\sigma_1\Phi_tY_t\dd B_t+(a_2-\sigma_1\sigma_2)\dd t+\sigma_2\dd B_t+\sigma_1\sigma_2\dd t\\
&=(a_1X_t+a_2)\dd t+(\sigma_1X_t+\sigma_2)\dd B_t,
\end{align*}
其中 $X_t=\Phi_tY_t$, $X_0=x$. (解的唯一性由下一节的定理~\ref{thm:sde-existence} 保证.)
\end{proof}

由于 $\Phi$ 是可以显式写出的正过程, 命题~\ref{prop:linear-sde} 给出的解是 $B$ 的显式泛函. 几个重要的特例如下.

\begin{example}[可解方程]\label{ex:linear-examples}
(1) \emph{几何布朗运动.} $a_2=\sigma_2=0$, $a_1=\mu$, $\sigma_1=\sigma$: $X_t=x\Phi_t=x\exp((\mu-\frac{1}{2}\sigma^2)t+\sigma B_t)$, 这就是 4.5.5 小节的几何布朗运动.

(2) \emph{Ornstein-Uhlenbeck 过程 (Vasicek 模型).} 设 $\kappa,\sigma>0$, $\theta\in\R$, 方程
\[
\dd X_t=\kappa(\theta-X_t)\dd t+\sigma\dd B_t
\]
描述一个被拉向长期均值 $\theta$ 的过程, 回复速度为 $\kappa$. 这里 $\sigma_1=0$, $\Phi_t=e^{-\kappa t}$, 所以
\[
X_t=\theta+(x-\theta)e^{-\kappa t}+\sigma\int_0^te^{-\kappa(t-s)}\dd B_s.
\]
由 Wiener 积分 (命题~\ref{prop:wiener-integral}), $X_t\sim N\big(\theta+(x-\theta)e^{-\kappa t},\frac{\sigma^2}{2\kappa}(1-e^{-2\kappa t})\big)$. 当 $t\to\infty$ 时, 分布收敛到 $N(\theta,\frac{\sigma^2}{2\kappa})$, 与初值无关 (图~\ref{fig:ou-paths}). 在金融中, 这是短期利率的 Vasicek 模型.

(3) \emph{布朗桥.} 方程 $\dd X_t=-\frac{X_t}{1-t}\dd t+\dd B_t$ ($0\le t<1$), $X_0=0$ 的系数依赖于时间. 同样的常数变易法 (以 $\Phi_t=1-t$ 代替) 给出
\[
X_t=(1-t)\int_0^t\frac{\dd B_s}{1-s},\qquad0\le t<1.
\]
由命题~\ref{prop:wiener-integral}, 对 $s\le t<1$, $\Cov(X_s,X_t)=(1-s)(1-t)\int_0^s\frac{\dd u}{(1-u)^2}=(1-s)(1-t)\big(\frac{1}{1-s}-1\big)=s(1-t)$, 这正是布朗桥的协方差 (定义~\ref{def:bridge}). 漂移 $-\frac{X_t}{1-t}$ 在 $t\to1$ 时越来越强, 把过程 ``钉'' 回零点.
\end{example}

\begin{figure}[htbp]
\centering
\begin{tikzpicture}
\begin{axis}[SAaxis, axis x line=bottom, axis y line=left, width=12cm, height=5.6cm, xmin=0, xmax=6.1, ymin=-2.2, ymax=3.4,
  xtick={0,1,...,6}, ytick={-2,-1,0,1,2,3}, xlabel={$t$}]
  \addplot[draw=none,fill=SAYellow!25] coordinates {(0.000,3.000) (0.050,3.290) (0.100,3.317) (0.150,3.302) (0.200,3.268) (0.250,3.223) (0.300,3.172) (0.350,3.117) (0.400,3.060) (0.450,3.002) (0.500,2.944) (0.550,2.886) (0.600,2.829) (0.650,2.772) (0.700,2.717) (0.750,2.664) (0.800,2.611) (0.850,2.561) (0.900,2.512) (0.950,2.464) (1.000,2.419) (1.050,2.375) (1.100,2.332) (1.150,2.291) (1.200,2.252) (1.250,2.214) (1.300,2.178) (1.350,2.144) (1.400,2.110) (1.450,2.078) (1.500,2.048) (1.550,2.019) (1.600,1.991) (1.650,1.964) (1.700,1.938) (1.750,1.914) (1.800,1.891) (1.850,1.868) (1.900,1.847) (1.950,1.827) (2.000,1.807) (2.050,1.789) (2.100,1.771) (2.150,1.754) (2.200,1.738) (2.250,1.723) (2.300,1.708) (2.350,1.694) (2.400,1.681) (2.450,1.668) (2.500,1.656) (2.550,1.644) (2.600,1.633) (2.650,1.623) (2.700,1.613) (2.750,1.603) (2.800,1.594) (2.850,1.585) (2.900,1.577) (2.950,1.569) (3.000,1.562) (3.050,1.555) (3.100,1.548) (3.150,1.541) (3.200,1.535) (3.250,1.529) (3.300,1.524) (3.350,1.519) (3.400,1.514) (3.450,1.509) (3.500,1.504) (3.550,1.500) (3.600,1.496) (3.650,1.492) (3.700,1.488) (3.750,1.484) (3.800,1.481) (3.850,1.478) (3.900,1.475) (3.950,1.472) (4.000,1.469) (4.050,1.466) (4.100,1.464) (4.150,1.461) (4.200,1.459) (4.250,1.457) (4.300,1.455) (4.350,1.453) (4.400,1.451) (4.450,1.449) (4.500,1.447) (4.550,1.446) (4.600,1.444) (4.650,1.443) (4.700,1.441) (4.750,1.440) (4.800,1.439) (4.850,1.438) (4.900,1.437) (4.950,1.435) (5.000,1.434) (5.050,1.433) (5.100,1.432) (5.150,1.432) (5.200,1.431) (5.250,1.430) (5.300,1.429) (5.350,1.428) (5.400,1.428) (5.450,1.427) (5.500,1.426) (5.550,1.426) (5.600,1.425) (5.650,1.425) (5.700,1.424) (5.750,1.424) (5.800,1.423) (5.850,1.423) (5.900,1.422) (5.950,1.422) (6.000,1.422) (6.000,-1.407) (5.950,-1.406) (5.900,-1.406) (5.850,-1.406) (5.800,-1.405) (5.750,-1.405) (5.700,-1.404) (5.650,-1.404) (5.600,-1.403) (5.550,-1.403) (5.500,-1.402) (5.450,-1.401) (5.400,-1.401) (5.350,-1.400) (5.300,-1.399) (5.250,-1.398) (5.200,-1.398) (5.150,-1.397) (5.100,-1.396) (5.050,-1.395) (5.000,-1.394) (4.950,-1.393) (4.900,-1.392) (4.850,-1.391) (4.800,-1.389) (4.750,-1.388) (4.700,-1.387) (4.650,-1.385) (4.600,-1.384) (4.550,-1.382) (4.500,-1.381) (4.450,-1.379) (4.400,-1.377) (4.350,-1.375) (4.300,-1.373) (4.250,-1.371) (4.200,-1.369) (4.150,-1.367) (4.100,-1.364) (4.050,-1.362) (4.000,-1.359) (3.950,-1.356) (3.900,-1.353) (3.850,-1.350) (3.800,-1.347) (3.750,-1.343) (3.700,-1.340) (3.650,-1.336) (3.600,-1.332) (3.550,-1.327) (3.500,-1.323) (3.450,-1.318) (3.400,-1.313) (3.350,-1.308) (3.300,-1.303) (3.250,-1.297) (3.200,-1.291) (3.150,-1.284) (3.100,-1.278) (3.050,-1.271) (3.000,-1.263) (2.950,-1.255) (2.900,-1.247) (2.850,-1.238) (2.800,-1.229) (2.750,-1.220) (2.700,-1.209) (2.650,-1.199) (2.600,-1.187) (2.550,-1.176) (2.500,-1.163) (2.450,-1.150) (2.400,-1.136) (2.350,-1.122) (2.300,-1.106) (2.250,-1.090) (2.200,-1.073) (2.150,-1.055) (2.100,-1.036) (2.050,-1.016) (2.000,-0.995) (1.950,-0.973) (1.900,-0.950) (1.850,-0.925) (1.800,-0.899) (1.750,-0.871) (1.700,-0.842) (1.650,-0.812) (1.600,-0.779) (1.550,-0.745) (1.500,-0.709) (1.450,-0.671) (1.400,-0.631) (1.350,-0.588) (1.300,-0.543) (1.250,-0.495) (1.200,-0.445) (1.150,-0.392) (1.100,-0.335) (1.050,-0.275) (1.000,-0.211) (0.950,-0.144) (0.900,-0.072) (0.850,0.004) (0.800,0.085) (0.750,0.171) (0.700,0.262) (0.650,0.360) (0.600,0.464) (0.550,0.576) (0.500,0.695) (0.450,0.823) (0.400,0.962) (0.350,1.111) (0.300,1.273) (0.250,1.449) (0.200,1.644) (0.150,1.862) (0.100,2.112) (0.050,2.417) (0.000,3.000)} -- cycle;
  \addplot[SADark!20] coordinates {(0,0) (6,0)};
  \addplot[SABlue,thin] coordinates {(0.000,3.000) (0.010,3.173) (0.020,2.887) (0.030,2.900) (0.040,2.815) (0.050,2.742) (0.060,2.693) (0.070,2.465) (0.080,2.418) (0.090,2.308) (0.100,2.615) (0.110,2.612) (0.120,2.551) (0.130,2.497) (0.140,2.406) (0.150,2.277) (0.160,2.215) (0.170,2.241) (0.180,2.195) (0.190,2.269) (0.200,2.226) (0.210,2.207) (0.220,2.338) (0.230,2.369) (0.240,2.296) (0.250,2.255) (0.260,2.286) (0.270,2.456) (0.280,2.404) (0.290,2.356) (0.300,2.433) (0.310,2.320) (0.320,2.268) (0.330,2.333) (0.340,2.368) (0.350,2.353) (0.360,2.397) (0.370,2.091) (0.380,2.172) (0.390,2.055) (0.400,1.869) (0.410,1.877) (0.420,1.929) (0.430,1.865) (0.440,1.739) (0.450,1.725) (0.460,1.702) (0.470,1.825) (0.480,1.881) (0.490,1.882) (0.500,1.974) (0.510,1.934) (0.520,1.822) (0.530,1.862) (0.540,1.902) (0.550,1.862) (0.560,1.765) (0.570,1.770) (0.580,1.505) (0.590,1.558) (0.600,1.592) (0.610,1.413) (0.620,1.405) (0.630,1.295) (0.640,1.357) (0.650,1.141) (0.660,1.039) (0.670,1.099) (0.680,1.203) (0.690,0.977) (0.700,0.918) (0.710,0.941) (0.720,0.871) (0.730,1.021) (0.740,0.892) (0.750,0.918) (0.760,0.805) (0.770,0.937) (0.780,0.925) (0.790,0.879) (0.800,0.699) (0.810,0.860) (0.820,0.926) (0.830,0.992) (0.840,1.095) (0.850,1.119) (0.860,1.044) (0.870,0.954) (0.880,0.865) (0.890,0.993) (0.900,0.838) (0.910,0.770) (0.920,0.730) (0.930,0.746) (0.940,0.795) (0.950,0.663) (0.960,0.484) (0.970,0.479) (0.980,0.595) (0.990,0.665) (1.000,0.679) (1.010,0.641) (1.020,0.664) (1.030,0.633) (1.040,0.708) (1.050,0.622) (1.060,0.629) (1.070,0.612) (1.080,0.660) (1.090,0.676) (1.100,0.923) (1.110,1.063) (1.120,1.201) (1.130,0.986) (1.140,0.942) (1.150,0.872) (1.160,0.917) (1.170,0.681) (1.180,0.791) (1.190,0.889) (1.200,0.751) (1.210,0.646) (1.220,0.560) (1.230,0.559) (1.240,0.617) (1.250,0.814) (1.260,0.787) (1.270,0.855) (1.280,0.862) (1.290,1.029) (1.300,1.092) (1.310,1.218) (1.320,1.098) (1.330,1.068) (1.340,0.977) (1.350,1.117) (1.360,1.171) (1.370,1.129) (1.380,1.073) (1.390,1.110) (1.400,1.029) (1.410,0.927) (1.420,0.965) (1.430,1.201) (1.440,1.164) (1.450,1.097) (1.460,0.970) (1.470,0.827) (1.480,0.871) (1.490,0.947) (1.500,0.939) (1.510,0.963) (1.520,0.965) (1.530,0.969) (1.540,0.807) (1.550,0.754) (1.560,0.757) (1.570,0.672) (1.580,0.618) (1.590,0.530) (1.600,0.492) (1.610,0.572) (1.620,0.526) (1.630,0.505) (1.640,0.581) (1.650,0.639) (1.660,0.802) (1.670,0.586) (1.680,0.665) (1.690,0.610) (1.700,0.618) (1.710,0.695) (1.720,0.796) (1.730,0.891) (1.740,0.898) (1.750,1.049) (1.760,1.011) (1.770,0.986) (1.780,1.056) (1.790,0.991) (1.800,1.196) (1.810,1.285) (1.820,1.489) (1.830,1.472) (1.840,1.419) (1.850,1.421) (1.860,1.480) (1.870,1.407) (1.880,1.431) (1.890,1.415) (1.900,1.562) (1.910,1.480) (1.920,1.570) (1.930,1.490) (1.940,1.380) (1.950,1.295) (1.960,1.164) (1.970,1.167) (1.980,1.258) (1.990,1.262) (2.000,1.311) (2.010,1.461) (2.020,1.473) (2.030,1.478) (2.040,1.436) (2.050,1.261) (2.060,1.324) (2.070,1.136) (2.080,1.190) (2.090,1.177) (2.100,1.277) (2.110,1.258) (2.120,1.163) (2.130,1.187) (2.140,1.120) (2.150,1.091) (2.160,1.084) (2.170,1.060) (2.180,1.030) (2.190,0.937) (2.200,0.909) (2.210,0.687) (2.220,0.665) (2.230,0.539) (2.240,0.645) (2.250,0.765) (2.260,0.563) (2.270,0.572) (2.280,0.554) (2.290,0.444) (2.300,0.493) (2.310,0.442) (2.320,0.262) (2.330,0.232) (2.340,0.215) (2.350,0.224) (2.360,0.099) (2.370,0.159) (2.380,0.231) (2.390,0.115) (2.400,0.048) (2.410,0.040) (2.420,-0.017) (2.430,0.156) (2.440,0.176) (2.450,0.073) (2.460,-0.006) (2.470,-0.011) (2.480,0.217) (2.490,0.197) (2.500,0.208) (2.510,0.257) (2.520,0.250) (2.530,0.475) (2.540,0.417) (2.550,0.487) (2.560,0.498) (2.570,0.555) (2.580,0.433) (2.590,0.604) (2.600,0.586) (2.610,0.584) (2.620,0.485) (2.630,0.407) (2.640,0.459) (2.650,0.553) (2.660,0.622) (2.670,0.737) (2.680,0.800) (2.690,0.795) (2.700,0.871) (2.710,0.921) (2.720,0.902) (2.730,0.965) (2.740,1.083) (2.750,1.096) (2.760,1.050) (2.770,1.111) (2.780,1.289) (2.790,1.255) (2.800,1.233) (2.810,1.207) (2.820,1.317) (2.830,1.121) (2.840,1.146) (2.850,1.254) (2.860,1.160) (2.870,1.297) (2.880,1.337) (2.890,1.269) (2.900,1.277) (2.910,1.112) (2.920,1.044) (2.930,1.216) (2.940,1.115) (2.950,1.287) (2.960,1.266) (2.970,1.353) (2.980,1.343) (2.990,1.107) (3.000,1.039) (3.010,1.048) (3.020,0.923) (3.030,0.779) (3.040,0.808) (3.050,0.735) (3.060,0.559) (3.070,0.486) (3.080,0.567) (3.090,0.514) (3.100,0.602) (3.110,0.748) (3.120,0.743) (3.130,0.624) (3.140,0.488) (3.150,0.480) (3.160,0.563) (3.170,0.585) (3.180,0.633) (3.190,0.583) (3.200,0.607) (3.210,0.554) (3.220,0.516) (3.230,0.382) (3.240,0.222) (3.250,0.172) (3.260,0.060) (3.270,-0.036) (3.280,0.031) (3.290,0.020) (3.300,0.310) (3.310,0.399) (3.320,0.332) (3.330,0.406) (3.340,0.438) (3.350,0.365) (3.360,0.471) (3.370,0.389) (3.380,0.081) (3.390,0.159) (3.400,0.096) (3.410,0.243) (3.420,0.176) (3.430,0.057) (3.440,0.210) (3.450,0.101) (3.460,0.120) (3.470,0.263) (3.480,0.272) (3.490,0.262) (3.500,0.269) (3.510,0.350) (3.520,0.266) (3.530,0.324) (3.540,0.378) (3.550,0.607) (3.560,0.567) (3.570,0.471) (3.580,0.539) (3.590,0.494) (3.600,0.518) (3.610,0.526) (3.620,0.637) (3.630,0.565) (3.640,0.551) (3.650,0.291) (3.660,0.200) (3.670,0.351) (3.680,0.279) (3.690,0.207) (3.700,0.084) (3.710,0.140) (3.720,0.157) (3.730,0.037) (3.740,0.103) (3.750,0.128) (3.760,0.202) (3.770,0.136) (3.780,0.307) (3.790,0.345) (3.800,0.330) (3.810,0.275) (3.820,0.248) (3.830,0.331) (3.840,0.301) (3.850,0.405) (3.860,0.459) (3.870,0.588) (3.880,0.664) (3.890,0.651) (3.900,0.454) (3.910,0.479) (3.920,0.553) (3.930,0.564) (3.940,0.614) (3.950,0.561) (3.960,0.517) (3.970,0.418) (3.980,0.550) (3.990,0.560) (4.000,0.505) (4.010,0.368) (4.020,0.349) (4.030,0.489) (4.040,0.530) (4.050,0.453) (4.060,0.590) (4.070,0.727) (4.080,0.817) (4.090,0.887) (4.100,0.994) (4.110,1.038) (4.120,1.105) (4.130,1.016) (4.140,1.155) (4.150,1.167) (4.160,1.356) (4.170,1.140) (4.180,1.198) (4.190,1.284) (4.200,1.237) (4.210,1.270) (4.220,1.332) (4.230,1.376) (4.240,1.309) (4.250,1.282) (4.260,1.370) (4.270,1.484) (4.280,1.479) (4.290,1.472) (4.300,1.477) (4.310,1.566) (4.320,1.445) (4.330,1.298) (4.340,1.298) (4.350,1.175) (4.360,1.105) (4.370,1.102) (4.380,1.139) (4.390,1.021) (4.400,1.091) (4.410,0.904) (4.420,0.949) (4.430,0.782) (4.440,0.648) (4.450,0.705) (4.460,0.658) (4.470,0.572) (4.480,0.522) (4.490,0.563) (4.500,0.552) (4.510,0.544) (4.520,0.643) (4.530,0.566) (4.540,0.513) (4.550,0.456) (4.560,0.568) (4.570,0.661) (4.580,0.782) (4.590,0.771) (4.600,0.748) (4.610,0.825) (4.620,0.778) (4.630,0.853) (4.640,0.916) (4.650,0.838) (4.660,0.821) (4.670,0.801) (4.680,0.966) (4.690,1.213) (4.700,1.191) (4.710,1.181) (4.720,0.942) (4.730,0.954) (4.740,0.944) (4.750,0.828) (4.760,0.898) (4.770,0.733) (4.780,0.704) (4.790,0.711) (4.800,0.580) (4.810,0.659) (4.820,0.758) (4.830,0.612) (4.840,0.466) (4.850,0.521) (4.860,0.513) (4.870,0.387) (4.880,0.354) (4.890,0.340) (4.900,0.322) (4.910,0.401) (4.920,0.436) (4.930,0.607) (4.940,0.738) (4.950,0.798) (4.960,0.897) (4.970,1.020) (4.980,1.080) (4.990,0.948) (5.000,1.024) (5.010,1.021) (5.020,1.172) (5.030,1.169) (5.040,1.210) (5.050,1.231) (5.060,1.386) (5.070,1.390) (5.080,1.349) (5.090,1.366) (5.100,1.309) (5.110,1.312) (5.120,1.298) (5.130,1.312) (5.140,1.491) (5.150,1.435) (5.160,1.425) (5.170,1.391) (5.180,1.444) (5.190,1.312) (5.200,1.230) (5.210,1.214) (5.220,1.256) (5.230,1.239) (5.240,1.177) (5.250,1.079) (5.260,1.119) (5.270,1.251) (5.280,1.271) (5.290,1.277) (5.300,1.309) (5.310,1.202) (5.320,1.159) (5.330,1.121) (5.340,1.139) (5.350,1.124) (5.360,0.992) (5.370,1.094) (5.380,1.252) (5.390,1.213) (5.400,1.241) (5.410,1.255) (5.420,1.321) (5.430,1.160) (5.440,1.309) (5.450,1.250) (5.460,1.363) (5.470,1.357) (5.480,1.305) (5.490,1.278) (5.500,1.355) (5.510,1.255) (5.520,1.206) (5.530,1.146) (5.540,0.963) (5.550,1.083) (5.560,1.130) (5.570,1.117) (5.580,1.221) (5.590,1.114) (5.600,1.096) (5.610,1.256) (5.620,1.341) (5.630,1.107) (5.640,0.764) (5.650,0.653) (5.660,0.752) (5.670,0.698) (5.680,0.669) (5.690,0.723) (5.700,0.628) (5.710,0.643) (5.720,0.715) (5.730,0.755) (5.740,0.786) (5.750,0.719) (5.760,0.763) (5.770,0.713) (5.780,0.835) (5.790,1.083) (5.800,0.998) (5.810,0.772) (5.820,0.722) (5.830,0.752) (5.840,0.849) (5.850,0.948) (5.860,0.861) (5.870,0.773) (5.880,0.892) (5.890,0.774) (5.900,0.875) (5.910,0.844) (5.920,0.912) (5.930,0.973) (5.940,1.040) (5.950,1.112) (5.960,1.002) (5.970,1.022) (5.980,1.013) (5.990,1.034) (6.000,1.146)};
  \addplot[SARed,thin] coordinates {(0.000,3.000) (0.010,2.943) (0.020,2.843) (0.030,2.742) (0.040,2.684) (0.050,2.732) (0.060,2.624) (0.070,2.713) (0.080,2.715) (0.090,2.666) (0.100,2.507) (0.110,2.506) (0.120,2.550) (0.130,2.349) (0.140,2.444) (0.150,2.584) (0.160,2.592) (0.170,2.615) (0.180,2.829) (0.190,2.923) (0.200,2.805) (0.210,2.775) (0.220,2.732) (0.230,2.690) (0.240,2.428) (0.250,2.479) (0.260,2.391) (0.270,2.384) (0.280,2.346) (0.290,2.410) (0.300,2.356) (0.310,2.254) (0.320,2.284) (0.330,2.401) (0.340,2.346) (0.350,2.459) (0.360,2.460) (0.370,2.383) (0.380,2.511) (0.390,2.404) (0.400,2.481) (0.410,2.407) (0.420,2.428) (0.430,2.281) (0.440,2.175) (0.450,2.259) (0.460,2.255) (0.470,2.237) (0.480,2.171) (0.490,2.194) (0.500,2.033) (0.510,1.901) (0.520,1.866) (0.530,1.765) (0.540,1.678) (0.550,1.654) (0.560,1.624) (0.570,1.585) (0.580,1.735) (0.590,1.819) (0.600,1.923) (0.610,1.901) (0.620,1.832) (0.630,1.669) (0.640,1.618) (0.650,1.394) (0.660,1.322) (0.670,1.385) (0.680,1.533) (0.690,1.613) (0.700,1.512) (0.710,1.546) (0.720,1.452) (0.730,1.519) (0.740,1.493) (0.750,1.563) (0.760,1.804) (0.770,1.878) (0.780,2.091) (0.790,2.088) (0.800,2.034) (0.810,2.147) (0.820,2.297) (0.830,2.219) (0.840,2.205) (0.850,2.040) (0.860,2.000) (0.870,2.068) (0.880,2.039) (0.890,1.962) (0.900,1.803) (0.910,1.821) (0.920,1.882) (0.930,1.923) (0.940,1.929) (0.950,1.783) (0.960,1.941) (0.970,1.916) (0.980,1.883) (0.990,1.937) (1.000,1.702) (1.010,1.679) (1.020,1.649) (1.030,1.669) (1.040,1.787) (1.050,1.816) (1.060,1.615) (1.070,1.566) (1.080,1.571) (1.090,1.570) (1.100,1.426) (1.110,1.247) (1.120,1.289) (1.130,1.310) (1.140,1.295) (1.150,1.111) (1.160,1.276) (1.170,0.975) (1.180,0.929) (1.190,1.128) (1.200,1.101) (1.210,1.113) (1.220,0.985) (1.230,1.110) (1.240,0.970) (1.250,0.986) (1.260,0.990) (1.270,0.875) (1.280,0.955) (1.290,0.804) (1.300,0.728) (1.310,0.718) (1.320,0.670) (1.330,0.710) (1.340,0.795) (1.350,0.758) (1.360,0.696) (1.370,0.638) (1.380,0.595) (1.390,0.670) (1.400,0.641) (1.410,0.676) (1.420,0.660) (1.430,0.569) (1.440,0.649) (1.450,0.540) (1.460,0.535) (1.470,0.678) (1.480,0.562) (1.490,0.551) (1.500,0.388) (1.510,0.333) (1.520,0.181) (1.530,0.143) (1.540,0.089) (1.550,0.101) (1.560,0.257) (1.570,0.283) (1.580,0.346) (1.590,0.337) (1.600,0.309) (1.610,0.223) (1.620,0.042) (1.630,0.048) (1.640,0.210) (1.650,0.255) (1.660,0.200) (1.670,0.362) (1.680,0.315) (1.690,0.250) (1.700,0.127) (1.710,0.030) (1.720,0.005) (1.730,0.204) (1.740,0.182) (1.750,0.077) (1.760,0.016) (1.770,-0.006) (1.780,-0.176) (1.790,-0.162) (1.800,-0.091) (1.810,-0.058) (1.820,0.003) (1.830,-0.126) (1.840,-0.294) (1.850,-0.163) (1.860,-0.052) (1.870,-0.205) (1.880,-0.182) (1.890,-0.178) (1.900,-0.122) (1.910,-0.100) (1.920,-0.149) (1.930,-0.128) (1.940,-0.111) (1.950,-0.115) (1.960,-0.298) (1.970,-0.593) (1.980,-0.666) (1.990,-0.822) (2.000,-0.789) (2.010,-0.710) (2.020,-0.676) (2.030,-0.706) (2.040,-0.692) (2.050,-0.569) (2.060,-0.636) (2.070,-0.538) (2.080,-0.462) (2.090,-0.432) (2.100,-0.283) (2.110,-0.351) (2.120,-0.254) (2.130,-0.344) (2.140,-0.417) (2.150,-0.518) (2.160,-0.469) (2.170,-0.558) (2.180,-0.547) (2.190,-0.654) (2.200,-0.542) (2.210,-0.354) (2.220,-0.328) (2.230,-0.161) (2.240,-0.104) (2.250,0.113) (2.260,0.122) (2.270,-0.035) (2.280,0.113) (2.290,0.100) (2.300,-0.004) (2.310,-0.070) (2.320,-0.231) (2.330,-0.245) (2.340,-0.349) (2.350,-0.198) (2.360,-0.130) (2.370,-0.294) (2.380,-0.441) (2.390,-0.493) (2.400,-0.680) (2.410,-0.684) (2.420,-0.682) (2.430,-0.752) (2.440,-0.799) (2.450,-0.741) (2.460,-0.592) (2.470,-0.438) (2.480,-0.430) (2.490,-0.340) (2.500,-0.410) (2.510,-0.365) (2.520,-0.457) (2.530,-0.380) (2.540,-0.500) (2.550,-0.591) (2.560,-0.586) (2.570,-0.524) (2.580,-0.662) (2.590,-0.709) (2.600,-0.688) (2.610,-0.659) (2.620,-0.736) (2.630,-0.719) (2.640,-0.738) (2.650,-0.591) (2.660,-0.530) (2.670,-0.656) (2.680,-0.869) (2.690,-0.850) (2.700,-0.724) (2.710,-0.700) (2.720,-0.605) (2.730,-0.647) (2.740,-0.776) (2.750,-0.694) (2.760,-0.678) (2.770,-0.603) (2.780,-0.734) (2.790,-0.709) (2.800,-0.764) (2.810,-0.676) (2.820,-0.727) (2.830,-0.780) (2.840,-0.831) (2.850,-0.697) (2.860,-0.687) (2.870,-0.756) (2.880,-0.850) (2.890,-0.946) (2.900,-0.791) (2.910,-0.963) (2.920,-1.038) (2.930,-0.965) (2.940,-0.838) (2.950,-0.844) (2.960,-0.758) (2.970,-0.674) (2.980,-0.570) (2.990,-0.559) (3.000,-0.389) (3.010,-0.364) (3.020,-0.206) (3.030,-0.312) (3.040,-0.213) (3.050,-0.334) (3.060,-0.477) (3.070,-0.420) (3.080,-0.496) (3.090,-0.443) (3.100,-0.542) (3.110,-0.374) (3.120,-0.365) (3.130,-0.471) (3.140,-0.445) (3.150,-0.348) (3.160,-0.417) (3.170,-0.357) (3.180,-0.145) (3.190,-0.053) (3.200,-0.048) (3.210,-0.099) (3.220,-0.036) (3.230,-0.181) (3.240,-0.476) (3.250,-0.533) (3.260,-0.630) (3.270,-0.444) (3.280,-0.421) (3.290,-0.509) (3.300,-0.614) (3.310,-0.552) (3.320,-0.623) (3.330,-0.669) (3.340,-0.665) (3.350,-0.644) (3.360,-0.621) (3.370,-0.597) (3.380,-0.455) (3.390,-0.275) (3.400,-0.430) (3.410,-0.162) (3.420,-0.141) (3.430,-0.009) (3.440,-0.075) (3.450,-0.035) (3.460,-0.006) (3.470,0.013) (3.480,0.105) (3.490,0.174) (3.500,0.080) (3.510,0.074) (3.520,0.154) (3.530,0.228) (3.540,0.200) (3.550,0.177) (3.560,0.168) (3.570,0.064) (3.580,-0.058) (3.590,0.107) (3.600,0.074) (3.610,0.090) (3.620,0.022) (3.630,0.159) (3.640,0.339) (3.650,0.321) (3.660,0.377) (3.670,0.305) (3.680,0.485) (3.690,0.580) (3.700,0.498) (3.710,0.625) (3.720,0.662) (3.730,0.704) (3.740,0.740) (3.750,0.674) (3.760,0.621) (3.770,0.684) (3.780,0.679) (3.790,0.585) (3.800,0.789) (3.810,0.767) (3.820,0.711) (3.830,0.642) (3.840,0.711) (3.850,0.808) (3.860,0.831) (3.870,0.916) (3.880,0.765) (3.890,0.727) (3.900,0.680) (3.910,0.680) (3.920,0.729) (3.930,0.525) (3.940,0.837) (3.950,0.749) (3.960,0.782) (3.970,0.863) (3.980,0.945) (3.990,0.966) (4.000,0.673) (4.010,0.564) (4.020,0.639) (4.030,0.546) (4.040,0.441) (4.050,0.207) (4.060,0.331) (4.070,0.383) (4.080,0.295) (4.090,0.283) (4.100,0.274) (4.110,0.164) (4.120,0.139) (4.130,0.166) (4.140,0.222) (4.150,0.349) (4.160,0.351) (4.170,0.304) (4.180,0.344) (4.190,0.280) (4.200,0.155) (4.210,0.270) (4.220,0.231) (4.230,0.243) (4.240,0.307) (4.250,0.196) (4.260,0.260) (4.270,0.304) (4.280,0.293) (4.290,0.332) (4.300,0.255) (4.310,0.035) (4.320,0.008) (4.330,-0.081) (4.340,0.008) (4.350,0.036) (4.360,-0.015) (4.370,-0.008) (4.380,0.128) (4.390,0.242) (4.400,0.278) (4.410,0.320) (4.420,0.291) (4.430,0.276) (4.440,0.334) (4.450,0.321) (4.460,0.349) (4.470,0.236) (4.480,0.069) (4.490,-0.099) (4.500,-0.140) (4.510,-0.005) (4.520,0.115) (4.530,0.016) (4.540,-0.138) (4.550,-0.056) (4.560,-0.007) (4.570,-0.142) (4.580,-0.034) (4.590,-0.015) (4.600,0.119) (4.610,0.153) (4.620,0.242) (4.630,0.273) (4.640,0.177) (4.650,0.154) (4.660,0.086) (4.670,0.169) (4.680,0.070) (4.690,-0.082) (4.700,-0.012) (4.710,0.104) (4.720,0.248) (4.730,0.356) (4.740,0.204) (4.750,0.133) (4.760,0.079) (4.770,0.077) (4.780,-0.137) (4.790,-0.135) (4.800,-0.105) (4.810,-0.198) (4.820,-0.183) (4.830,-0.061) (4.840,-0.055) (4.850,0.052) (4.860,0.154) (4.870,0.194) (4.880,0.231) (4.890,0.193) (4.900,0.316) (4.910,0.267) (4.920,0.196) (4.930,0.265) (4.940,0.192) (4.950,0.233) (4.960,0.337) (4.970,0.343) (4.980,0.351) (4.990,0.441) (5.000,0.512) (5.010,0.355) (5.020,0.223) (5.030,0.354) (5.040,0.391) (5.050,0.181) (5.060,0.047) (5.070,0.050) (5.080,-0.046) (5.090,0.003) (5.100,0.012) (5.110,0.062) (5.120,-0.053) (5.130,0.012) (5.140,0.208) (5.150,0.129) (5.160,0.260) (5.170,0.358) (5.180,0.420) (5.190,0.412) (5.200,0.294) (5.210,0.206) (5.220,0.282) (5.230,0.379) (5.240,0.499) (5.250,0.592) (5.260,0.678) (5.270,0.631) (5.280,0.506) (5.290,0.524) (5.300,0.571) (5.310,0.576) (5.320,0.422) (5.330,0.368) (5.340,0.220) (5.350,0.235) (5.360,0.130) (5.370,0.115) (5.380,0.148) (5.390,0.163) (5.400,0.300) (5.410,0.189) (5.420,-0.027) (5.430,0.070) (5.440,0.211) (5.450,0.354) (5.460,0.382) (5.470,0.434) (5.480,0.445) (5.490,0.314) (5.500,0.257) (5.510,0.307) (5.520,0.250) (5.530,0.440) (5.540,0.285) (5.550,0.242) (5.560,0.353) (5.570,0.292) (5.580,0.180) (5.590,0.046) (5.600,-0.024) (5.610,0.129) (5.620,0.172) (5.630,0.225) (5.640,0.174) (5.650,0.259) (5.660,0.271) (5.670,0.180) (5.680,0.136) (5.690,0.158) (5.700,0.276) (5.710,0.271) (5.720,0.321) (5.730,0.522) (5.740,0.441) (5.750,0.316) (5.760,0.415) (5.770,0.457) (5.780,0.446) (5.790,0.472) (5.800,0.438) (5.810,0.488) (5.820,0.453) (5.830,0.532) (5.840,0.413) (5.850,0.332) (5.860,0.343) (5.870,0.297) (5.880,0.366) (5.890,0.362) (5.900,0.308) (5.910,0.330) (5.920,0.277) (5.930,0.123) (5.940,0.013) (5.950,-0.063) (5.960,-0.106) (5.970,-0.091) (5.980,0.004) (5.990,0.063) (6.000,0.244)};
  \addplot[SAGreenText,thin] coordinates {(0.000,3.000) (0.010,3.101) (0.020,3.056) (0.030,3.194) (0.040,3.284) (0.050,3.210) (0.060,3.214) (0.070,3.295) (0.080,3.291) (0.090,3.364) (0.100,3.374) (0.110,3.415) (0.120,3.354) (0.130,3.363) (0.140,3.298) (0.150,3.337) (0.160,3.140) (0.170,3.147) (0.180,3.245) (0.190,3.275) (0.200,3.245) (0.210,3.271) (0.220,3.159) (0.230,3.049) (0.240,2.988) (0.250,2.948) (0.260,3.020) (0.270,2.905) (0.280,2.981) (0.290,2.799) (0.300,2.698) (0.310,2.672) (0.320,2.765) (0.330,2.681) (0.340,2.446) (0.350,2.468) (0.360,2.445) (0.370,2.318) (0.380,2.195) (0.390,2.109) (0.400,2.102) (0.410,2.152) (0.420,1.988) (0.430,2.201) (0.440,2.180) (0.450,2.125) (0.460,2.176) (0.470,2.170) (0.480,2.385) (0.490,2.475) (0.500,2.418) (0.510,2.360) (0.520,2.291) (0.530,2.239) (0.540,2.288) (0.550,2.205) (0.560,2.348) (0.570,2.321) (0.580,2.288) (0.590,2.364) (0.600,2.148) (0.610,2.038) (0.620,2.081) (0.630,2.066) (0.640,2.048) (0.650,2.072) (0.660,1.996) (0.670,1.963) (0.680,1.812) (0.690,1.612) (0.700,1.726) (0.710,1.620) (0.720,1.537) (0.730,1.664) (0.740,1.757) (0.750,1.692) (0.760,1.626) (0.770,1.606) (0.780,1.499) (0.790,1.542) (0.800,1.545) (0.810,1.547) (0.820,1.610) (0.830,1.670) (0.840,1.840) (0.850,1.977) (0.860,1.791) (0.870,1.866) (0.880,1.740) (0.890,1.584) (0.900,1.643) (0.910,1.561) (0.920,1.518) (0.930,1.613) (0.940,1.535) (0.950,1.434) (0.960,1.456) (0.970,1.567) (0.980,1.608) (0.990,1.441) (1.000,1.366) (1.010,1.341) (1.020,1.524) (1.030,1.601) (1.040,1.674) (1.050,1.781) (1.060,1.736) (1.070,1.770) (1.080,1.805) (1.090,1.694) (1.100,1.679) (1.110,1.723) (1.120,1.602) (1.130,1.670) (1.140,1.733) (1.150,1.822) (1.160,1.731) (1.170,1.728) (1.180,1.859) (1.190,1.924) (1.200,1.830) (1.210,1.838) (1.220,1.965) (1.230,1.941) (1.240,1.964) (1.250,2.203) (1.260,2.181) (1.270,2.153) (1.280,2.035) (1.290,1.953) (1.300,2.085) (1.310,2.053) (1.320,1.918) (1.330,1.810) (1.340,1.796) (1.350,1.857) (1.360,1.768) (1.370,1.825) (1.380,1.761) (1.390,1.782) (1.400,1.818) (1.410,1.863) (1.420,1.804) (1.430,1.772) (1.440,1.807) (1.450,1.662) (1.460,1.665) (1.470,1.701) (1.480,1.789) (1.490,1.769) (1.500,1.734) (1.510,1.732) (1.520,1.643) (1.530,1.607) (1.540,1.632) (1.550,1.645) (1.560,1.525) (1.570,1.478) (1.580,1.253) (1.590,1.320) (1.600,1.254) (1.610,1.265) (1.620,1.143) (1.630,1.123) (1.640,1.155) (1.650,1.241) (1.660,1.092) (1.670,1.164) (1.680,1.135) (1.690,1.045) (1.700,1.081) (1.710,1.039) (1.720,0.998) (1.730,0.917) (1.740,0.928) (1.750,0.921) (1.760,0.893) (1.770,0.787) (1.780,0.726) (1.790,0.757) (1.800,0.875) (1.810,0.862) (1.820,0.992) (1.830,0.979) (1.840,0.859) (1.850,0.911) (1.860,0.804) (1.870,0.834) (1.880,0.930) (1.890,0.888) (1.900,0.869) (1.910,1.104) (1.920,1.114) (1.930,1.128) (1.940,1.192) (1.950,1.169) (1.960,1.146) (1.970,1.171) (1.980,0.920) (1.990,1.049) (2.000,0.931) (2.010,0.813) (2.020,0.856) (2.030,0.849) (2.040,0.714) (2.050,0.726) (2.060,0.604) (2.070,0.527) (2.080,0.371) (2.090,0.283) (2.100,0.318) (2.110,0.325) (2.120,0.351) (2.130,0.437) (2.140,0.427) (2.150,0.355) (2.160,0.322) (2.170,0.481) (2.180,0.478) (2.190,0.503) (2.200,0.487) (2.210,0.504) (2.220,0.540) (2.230,0.532) (2.240,0.510) (2.250,0.425) (2.260,0.335) (2.270,0.371) (2.280,0.439) (2.290,0.525) (2.300,0.523) (2.310,0.434) (2.320,0.491) (2.330,0.545) (2.340,0.389) (2.350,0.463) (2.360,0.514) (2.370,0.673) (2.380,0.552) (2.390,0.418) (2.400,0.490) (2.410,0.305) (2.420,0.228) (2.430,0.116) (2.440,0.012) (2.450,-0.179) (2.460,-0.202) (2.470,-0.105) (2.480,-0.061) (2.490,-0.018) (2.500,-0.055) (2.510,0.088) (2.520,0.138) (2.530,0.240) (2.540,0.082) (2.550,0.115) (2.560,-0.001) (2.570,0.033) (2.580,0.051) (2.590,-0.062) (2.600,-0.045) (2.610,0.070) (2.620,0.118) (2.630,0.175) (2.640,0.088) (2.650,0.019) (2.660,-0.092) (2.670,-0.078) (2.680,0.040) (2.690,0.083) (2.700,0.082) (2.710,-0.030) (2.720,0.014) (2.730,-0.010) (2.740,-0.198) (2.750,-0.368) (2.760,-0.398) (2.770,-0.482) (2.780,-0.314) (2.790,-0.307) (2.800,-0.407) (2.810,-0.451) (2.820,-0.288) (2.830,-0.379) (2.840,-0.299) (2.850,-0.252) (2.860,-0.104) (2.870,-0.050) (2.880,-0.089) (2.890,-0.036) (2.900,-0.028) (2.910,0.069) (2.920,0.063) (2.930,0.259) (2.940,0.247) (2.950,0.288) (2.960,0.221) (2.970,0.232) (2.980,0.234) (2.990,0.249) (3.000,0.337) (3.010,0.281) (3.020,0.196) (3.030,0.305) (3.040,0.275) (3.050,0.207) (3.060,0.320) (3.070,0.396) (3.080,0.131) (3.090,0.094) (3.100,0.172) (3.110,0.255) (3.120,0.203) (3.130,0.357) (3.140,0.277) (3.150,0.253) (3.160,0.265) (3.170,0.264) (3.180,0.343) (3.190,0.443) (3.200,0.436) (3.210,0.572) (3.220,0.501) (3.230,0.477) (3.240,0.645) (3.250,0.651) (3.260,0.753) (3.270,0.705) (3.280,0.674) (3.290,0.660) (3.300,0.673) (3.310,0.636) (3.320,0.634) (3.330,0.697) (3.340,0.812) (3.350,0.818) (3.360,0.877) (3.370,0.962) (3.380,0.849) (3.390,0.688) (3.400,0.844) (3.410,0.970) (3.420,0.969) (3.430,0.998) (3.440,0.898) (3.450,0.755) (3.460,0.782) (3.470,0.631) (3.480,0.718) (3.490,0.647) (3.500,0.673) (3.510,0.861) (3.520,1.042) (3.530,1.008) (3.540,1.107) (3.550,0.997) (3.560,0.932) (3.570,0.963) (3.580,1.093) (3.590,1.029) (3.600,1.075) (3.610,1.077) (3.620,0.944) (3.630,0.924) (3.640,1.067) (3.650,1.050) (3.660,1.092) (3.670,1.166) (3.680,1.084) (3.690,1.198) (3.700,1.338) (3.710,1.485) (3.720,1.603) (3.730,1.544) (3.740,1.473) (3.750,1.383) (3.760,1.282) (3.770,1.193) (3.780,1.281) (3.790,1.416) (3.800,1.341) (3.810,1.410) (3.820,1.273) (3.830,1.230) (3.840,1.121) (3.850,0.950) (3.860,0.948) (3.870,0.948) (3.880,0.900) (3.890,0.841) (3.900,0.736) (3.910,0.717) (3.920,0.652) (3.930,0.671) (3.940,0.546) (3.950,0.564) (3.960,0.611) (3.970,0.729) (3.980,0.828) (3.990,0.839) (4.000,0.764) (4.010,0.576) (4.020,0.570) (4.030,0.706) (4.040,0.865) (4.050,0.604) (4.060,0.524) (4.070,0.485) (4.080,0.377) (4.090,0.393) (4.100,0.394) (4.110,0.339) (4.120,0.333) (4.130,0.172) (4.140,0.074) (4.150,-0.015) (4.160,0.075) (4.170,-0.087) (4.180,-0.106) (4.190,-0.027) (4.200,-0.032) (4.210,0.050) (4.220,-0.015) (4.230,0.029) (4.240,0.119) (4.250,0.341) (4.260,0.410) (4.270,0.549) (4.280,0.545) (4.290,0.551) (4.300,0.702) (4.310,0.754) (4.320,0.889) (4.330,0.952) (4.340,1.154) (4.350,1.381) (4.360,1.286) (4.370,1.459) (4.380,1.490) (4.390,1.524) (4.400,1.499) (4.410,1.557) (4.420,1.433) (4.430,1.349) (4.440,1.261) (4.450,1.280) (4.460,1.361) (4.470,1.423) (4.480,1.448) (4.490,1.377) (4.500,1.191) (4.510,1.222) (4.520,1.227) (4.530,1.254) (4.540,1.386) (4.550,1.357) (4.560,1.326) (4.570,1.086) (4.580,1.133) (4.590,1.120) (4.600,1.038) (4.610,1.069) (4.620,0.958) (4.630,0.897) (4.640,1.087) (4.650,1.153) (4.660,1.237) (4.670,1.302) (4.680,1.356) (4.690,1.364) (4.700,1.581) (4.710,1.561) (4.720,1.621) (4.730,1.586) (4.740,1.537) (4.750,1.282) (4.760,1.080) (4.770,1.226) (4.780,1.364) (4.790,1.309) (4.800,1.295) (4.810,1.345) (4.820,1.235) (4.830,1.401) (4.840,1.285) (4.850,1.203) (4.860,0.917) (4.870,1.061) (4.880,1.058) (4.890,1.210) (4.900,1.100) (4.910,1.124) (4.920,1.122) (4.930,1.132) (4.940,1.329) (4.950,1.543) (4.960,1.666) (4.970,1.611) (4.980,1.685) (4.990,1.804) (5.000,1.842) (5.010,1.882) (5.020,1.913) (5.030,1.805) (5.040,1.813) (5.050,1.955) (5.060,1.838) (5.070,1.829) (5.080,1.895) (5.090,1.741) (5.100,1.711) (5.110,1.527) (5.120,1.452) (5.130,1.298) (5.140,1.197) (5.150,1.137) (5.160,1.099) (5.170,1.068) (5.180,0.819) (5.190,0.969) (5.200,0.933) (5.210,0.815) (5.220,0.744) (5.230,0.768) (5.240,0.660) (5.250,0.613) (5.260,0.560) (5.270,0.608) (5.280,0.527) (5.290,0.567) (5.300,0.596) (5.310,0.576) (5.320,0.693) (5.330,0.689) (5.340,0.565) (5.350,0.589) (5.360,0.595) (5.370,0.585) (5.380,0.602) (5.390,0.830) (5.400,0.845) (5.410,0.882) (5.420,0.849) (5.430,0.971) (5.440,0.956) (5.450,0.866) (5.460,0.736) (5.470,0.781) (5.480,0.847) (5.490,0.681) (5.500,0.602) (5.510,0.660) (5.520,0.756) (5.530,0.617) (5.540,0.570) (5.550,0.647) (5.560,0.587) (5.570,0.536) (5.580,0.477) (5.590,0.342) (5.600,0.343) (5.610,0.370) (5.620,0.227) (5.630,0.234) (5.640,0.342) (5.650,0.359) (5.660,0.450) (5.670,0.390) (5.680,0.522) (5.690,0.439) (5.700,0.622) (5.710,0.571) (5.720,0.493) (5.730,0.399) (5.740,0.354) (5.750,0.301) (5.760,0.229) (5.770,0.194) (5.780,0.258) (5.790,0.343) (5.800,0.306) (5.810,0.425) (5.820,0.322) (5.830,0.216) (5.840,0.101) (5.850,0.159) (5.860,0.003) (5.870,-0.136) (5.880,-0.008) (5.890,0.162) (5.900,0.294) (5.910,0.142) (5.920,0.141) (5.930,0.222) (5.940,0.229) (5.950,0.065) (5.960,0.030) (5.970,-0.140) (5.980,-0.262) (5.990,-0.046) (6.000,0.021)};
  \addplot[SADark,thick,dashed,domain=0:6,samples=80] {3*exp(-x)};
\end{axis}
\end{tikzpicture}
\caption{Ornstein-Uhlenbeck 过程 $\dd X_t=-X_t\dd t+\dd B_t$ ($\kappa=\sigma=1$, $\theta=0$) 从 $X_0=3$ 出发的三条轨道. 虚线为均值 $3e^{-t}$, 阴影为均值 $\pm2$ 倍标准差的范围. 均值指数衰减, 方差增长到平稳值 $\frac12$}
\label{fig:ou-paths}
\end{figure}

% ---------------------------------------------------------------------
\subsection{可约化为布朗运动的方程}

若扩散系数 $\sigma(x)>0$ 不依赖于时间且足够光滑, 适当的变量替换可以把扩散系数化为 $1$. 设 $F$ 满足 $F'(x)=\frac{1}{\sigma(x)}$. 由 It\^o 公式, $Y_t=F(X_t)$ 满足
\begin{equation}\label{eq:lamperti}
\dd Y_t=F'(X_t)\dd X_t+\frac{1}{2}F''(X_t)\sigma(X_t)^2\dd t=\Big(\frac{b(X_t)}{\sigma(X_t)}-\frac{1}{2}\sigma'(X_t)\Big)\dd t+\dd B_t,
\end{equation}
因为 $F''=-\frac{\sigma'}{\sigma^2}$. 这称为 \textbf{Lamperti 变换}. 若 \eqref{eq:lamperti} 中的漂移为零, 则 $Y$ 是 (平移的) 布朗运动, 原方程的解为 $X_t=F^{-1}(F(x)+B_t)$.

\begin{example}\label{ex:lamperti}
方程 $\dd X_t=\frac{1}{2}X_t\dd t+\sqrt{1+X_t^2}\dd B_t$: 取 $F(x)=\operatorname{arsinh}x$, 则 $F'=\frac{1}{\sqrt{1+x^2}}$, 漂移 $\frac{b}{\sigma}-\frac{1}{2}\sigma'=\frac{x/2}{\sqrt{1+x^2}}-\frac{1}{2}\cdot\frac{x}{\sqrt{1+x^2}}=0$. 所以 $X_t=\sinh\big(\operatorname{arsinh}x+B_t\big)$. 可以用 It\^o 公式直接验证.
\end{example}

% =====================================================================
\section{存在性与唯一性}
\label{sec:6-2}

常微分方程 $\dot{x}=b(x)$ 的经典理论告诉我们: 若 $b$ 满足 Lipschitz 条件, 则解存在唯一, 证明方法是 Picard 迭代. 随机微分方程的情形完全类似, 只是距离要用 $\E\sup_t\abs{\cdot}^2$ 来度量, 随机积分项要用 Doob 不等式与 It\^o 等距来估计.

% ---------------------------------------------------------------------
\subsection{Lipschitz 条件下的强解}

本节假设存在常数 $K$, 使得对一切 $t\in[0,T]$ 与 $x,y\in\R$,
\begin{align}
&\abs{b(t,x)-b(t,y)}+\abs{\sigma(t,x)-\sigma(t,y)}\le K\abs{x-y},\tag{H1}\label{eq:H1}\\
&\abs{b(t,x)}+\abs{\sigma(t,x)}\le K(1+\abs{x}).\tag{H2}\label{eq:H2}
\end{align}
\eqref{eq:H1} 是关于 $x$ 的一致 Lipschitz 条件, \eqref{eq:H2} 是线性增长条件 (若 $b(t,0)$ 与 $\sigma(t,0)$ 有界, 它由 \eqref{eq:H1} 推出). 记 $\mathscr{C}^2_T$ 为满足 $\E\sup_{t\le T}\abs{X_t}^2<\infty$ 的连续适应过程全体.

\begin{lemma}[Gronwall 不等式]\label{lem:gronwall}
设 $u$ 是 $[0,T]$ 上的非负有界 Borel 函数, $a,c\ge0$, 且 $u(t)\le a+c\int_0^tu(s)\dd s$ 对一切 $t\in[0,T]$ 成立. 则 $u(t)\le ae^{ct}$.
\end{lemma}

\begin{proof}
对 $n$ 归纳证明
\[
u(t)\le a\sum_{k=0}^{n-1}\frac{(ct)^k}{k!}+c^n\int_0^t\frac{(t-s)^{n-1}}{(n-1)!}u(s)\dd s.
\]
$n=1$ 即假设. 设它对 $n$ 成立, 把 $u(s)\le a+c\int_0^su(r)\dd r$ 代入最后一项, 并由 Fubini 定理交换积分次序,
\begin{align*}
c^n\int_0^t\frac{(t-s)^{n-1}}{(n-1)!}u(s)\dd s&\le a\frac{(ct)^n}{n!}+c^{n+1}\int_0^tu(r)\int_r^t\frac{(t-s)^{n-1}}{(n-1)!}\dd s\dd r\\
&=a\frac{(ct)^n}{n!}+c^{n+1}\int_0^t\frac{(t-r)^n}{n!}u(r)\dd r.
\end{align*}
归纳完成. 最后一项不超过 $\sup u\cdot\frac{(ct)^n}{n!}\to0$, 令 $n\to\infty$ 即得.
\end{proof}

对连续适应过程 $X$ 与 $\mathcal{F}_0$-可测的 $\xi$, 记
\[
\Psi(\xi,X)_t=\xi+\int_0^tb(s,X_s)\dd s+\int_0^t\sigma(s,X_s)\dd B_s,\qquad t\in[0,T].
\]
方程 \eqref{eq:sde} 的强解就是映射 $X\mapsto\Psi(\xi,X)$ 的不动点. 下面的估计是全部证明的核心.

\begin{lemma}[基本估计]\label{lem:sde-estimate}
设 \eqref{eq:H1} 成立, $X,Y$ 是循序可测过程, $\E\int_0^T(X_s^2+Y_s^2)\dd s<\infty$, $\xi,\eta\in L^2(\mathcal{F}_0)$, 且 $b(\cdot,X),\sigma(\cdot,X)$ 等满足 $\E\int_0^T(b^2+\sigma^2)\dd s<\infty$ (在 \eqref{eq:H2} 下自动成立). 则对 $t\le T$,
\[
\E\Big[\sup_{s\le t}\abs{\Psi(\xi,X)_s-\Psi(\eta,Y)_s}^2\Big]\le3\E\abs{\xi-\eta}^2+L\int_0^t\E\abs{X_s-Y_s}^2\dd s,
\]
其中 $L=3K^2(T+4)$.
\end{lemma}

\begin{proof}
记 $\Delta b_s=b(s,X_s)-b(s,Y_s)$, $\Delta\sigma_s=\sigma(s,X_s)-\sigma(s,Y_s)$. 由 $(a_1+a_2+a_3)^2\le3(a_1^2+a_2^2+a_3^2)$,
\[
\sup_{s\le t}\abs{\Psi(\xi,X)_s-\Psi(\eta,Y)_s}^2\le3\abs{\xi-\eta}^2+3\sup_{s\le t}\Big\lvert\int_0^s\Delta b_u\dd u\Big\rvert^2+3\sup_{s\le t}\Big\lvert\int_0^s\Delta\sigma_u\dd B_u\Big\rvert^2.
\]
对第二项, 由 Cauchy-Schwarz 不等式与 \eqref{eq:H1}, $\big(\int_0^t\abs{\Delta b_u}\dd u\big)^2\le t\int_0^t\abs{\Delta b_u}^2\dd u\le TK^2\int_0^t\abs{X_u-Y_u}^2\dd u$. 对第三项, $\Delta\sigma\in\mathscr{L}^2_T$, 由 Doob 不等式~\ref{thm:ct-doob} 与 It\^o 等距,
\[
\E\sup_{s\le t}\Big\lvert\int_0^s\Delta\sigma_u\dd B_u\Big\rvert^2\le4\E\int_0^t\abs{\Delta\sigma_u}^2\dd u\le4K^2\int_0^t\E\abs{X_u-Y_u}^2\dd u.
\]
取期望并相加即得.
\end{proof}

\begin{theorem}[强解的存在唯一性]\label{thm:sde-existence}
设 \eqref{eq:H1}\eqref{eq:H2} 成立, $\xi\in L^2(\mathcal{F}_0)$. 则方程 \eqref{eq:sde} 有强解 $X\in\mathscr{C}^2_T$, 并且有轨道唯一性; 同时存在只依赖于 $K,T$ 的常数 $C$, 使
\begin{equation}\label{eq:sde-moment}
\E\Big[\sup_{t\le T}\abs{X_t}^2\Big]\le C\big(1+\E\xi^2\big).
\end{equation}
\end{theorem}

\begin{proof}
\emph{唯一性.} 设 $X$, $Y$ 都是强解 (不必属于 $\mathscr{C}^2_T$). 对 $N\ge1$, 令 $\tau_N=\inf\{t:\abs{X_t}\vee\abs{Y_t}\ge N\}\wedge T$. 由停止性质 (定理~\ref{thm:ito-local}),
\[
X_{t\wedge\tau_N}=\xi+\int_0^t\1_{\{s\le\tau_N\}}b(s,X_s)\dd s+\int_0^t\1_{\{s\le\tau_N\}}\sigma(s,X_s)\dd B_s,
\]
对 $Y$ 同理. 在 $\{s\le\tau_N\}$ 上 $X_s=X_{s\wedge\tau_N}$, 所以 $\1_{\{s\le\tau_N\}}\abs{b(s,X_s)-b(s,Y_s)}\le K\abs{X_{s\wedge\tau_N}-Y_{s\wedge\tau_N}}$, 对 $\sigma$ 同理; 并且在 $(0,\tau_N]$ 上 $\abs{X},\abs{Y}\le N$, 所以被积过程有界. 与引理~\ref{lem:sde-estimate} 的证明相同 (此处初值相同), 函数 $e(t)=\E\sup_{s\le t}\abs{X_{s\wedge\tau_N}-Y_{s\wedge\tau_N}}^2$ 满足 $e(t)\le L\int_0^te(s)\dd s$. $e\le4N^2$ 有界 (若 $\abs{\xi}\ge N$, 则 $\tau_N=0$, 差为零; 否则停止后的两个过程都被 $N$ 控制), 由 Gronwall 不等式 $e\equiv0$. 所以 $X=Y$ 在 $[0,\tau_N]$ 上成立. 由轨道连续性, $\tau_N\uparrow T$ 且最终等于 $T$, 所以 $X$ 与 $Y$ 不可区分.

\emph{存在性.} 令 $X^0_t\equiv\xi$, $X^{k+1}=\Psi(\xi,X^k)$. 若 $X^k\in\mathscr{C}^2_T$, 则由 \eqref{eq:H2}, $b(\cdot,X^k)^2+\sigma(\cdot,X^k)^2\le2K^2(1+\abs{X^k}^2)$ 可积, 与引理~\ref{lem:sde-estimate} 的证明相同可得 $\E\sup_t\abs{X^{k+1}_t}^2<\infty$, 所以每个 $X^k\in\mathscr{C}^2_T$. 记 $D_k(t)=\E\sup_{s\le t}\abs{X^{k+1}_s-X^k_s}^2$. 由引理~\ref{lem:sde-estimate} (初值相同),
\[
D_k(t)\le L\int_0^t\E\abs{X^k_s-X^{k-1}_s}^2\dd s\le L\int_0^tD_{k-1}(s)\dd s.
\]
归纳得 $D_k(t)\le D_0(T)\frac{(Lt)^k}{k!}$. 记 $\norm{Z}=\big(\E\sup_{t\le T}\abs{Z_t}^2\big)^{1/2}$, 则 $\sum_k\norm{X^{k+1}-X^k}\le\sum_k\big(D_0(T)\frac{(LT)^k}{k!}\big)^{1/2}<\infty$. 于是 $\E\sum_k\sup_t\abs{X^{k+1}_t-X^k_t}<\infty$, 几乎必然地 $X^k$ 在 $[0,T]$ 上一致收敛, 极限 $X$ 连续、适应 (在零概率集上令 $X\equiv0$). 由 Fatou 引理与 Minkowski 不等式, $\norm{X-X^k}\le\sum_{j\ge k}\norm{X^{j+1}-X^j}\to0$, 特别地 $X\in\mathscr{C}^2_T$. 再由引理~\ref{lem:sde-estimate},
\[
\norm{\Psi(\xi,X)-X^{k+1}}^2=\norm{\Psi(\xi,X)-\Psi(\xi,X^k)}^2\le LT\norm{X-X^k}^2\to0,
\]
而 $X^{k+1}\to X$, 所以 $\norm{\Psi(\xi,X)-X}=0$: $X$ 与 $\Psi(\xi,X)$ 都连续, 二者不可区分, 即 $X$ 是强解.

\emph{矩估计.} 记 $m(t)=\E\sup_{s\le t}\abs{X_s}^2<\infty$. 与引理~\ref{lem:sde-estimate} 的证明相同, 用 \eqref{eq:H2} 代替 \eqref{eq:H1} ($\abs{b}^2,\abs{\sigma}^2\le2K^2(1+\abs{x}^2)$),
\begin{align*}
m(t)&\le3\E\xi^2+6K^2(T+4)\int_0^t\big(1+m(s)\big)\dd s\\
&\le3\E\xi^2+6K^2(T+4)T+6K^2(T+4)\int_0^tm(s)\dd s.
\end{align*}
由 Gronwall 不等式, $m(T)\le\big(3\E\xi^2+6K^2(T+4)T\big)e^{6K^2(T+4)T}$, 即得 \eqref{eq:sde-moment}.
\end{proof}

同样的论证给出解对初值的连续依赖.

\begin{proposition}[对初值的连续依赖]\label{prop:sde-continuity}
设 \eqref{eq:H1}\eqref{eq:H2} 成立, $X^\xi$, $X^\eta$ 分别是初值为 $\xi,\eta\in L^2(\mathcal{F}_0)$ 的解. 则
\[
\E\Big[\sup_{t\le T}\abs{X^\xi_t-X^\eta_t}^2\Big]\le3e^{LT}\,\E\abs{\xi-\eta}^2.
\]
\end{proposition}

\begin{proof}
由引理~\ref{lem:sde-estimate}, $u(t)=\E\sup_{s\le t}\abs{X^\xi_s-X^\eta_s}^2$ 满足 $u(t)\le3\E\abs{\xi-\eta}^2+L\int_0^tu(s)\dd s$, 再用 Gronwall 不等式.
\end{proof}

\begin{remark}
(1) \emph{初值.} 条件 $\xi\in L^2$ 可以去掉: 对一般的 $\mathcal{F}_0$-可测的 $\xi$, 在 $\{\abs{\xi}\le n\}\in\mathcal{F}_0$ 上用初值 $\xi\1_{\{\abs{\xi}\le n\}}$ 的解, 由局部性 (引理~\ref{lem:locality}) 与唯一性, 这些解可以拼接起来.

(2) \emph{Lipschitz 条件的作用.} 若只有局部 Lipschitz 条件, 解在 ``爆炸时刻'' 之前存在唯一. 例如 $\dot{x}=x^2$, $x_0=1$ 的解 $\frac{1}{1-t}$ 在 $t=1$ 爆炸, 线性增长条件 \eqref{eq:H2} 正是用来排除这种情况的. 若连局部 Lipschitz 条件也不满足, 唯一性可能失效: $\dot{x}=3x^{2/3}$, $x_0=0$ 有解 $x\equiv0$ 与 $x=t^3$.

(3) \emph{Yamada-Watanabe 条件.} 对一维方程, Lipschitz 条件对 $\sigma$ 可以放宽: 若 $b$ 满足 Lipschitz 条件, 而 $\abs{\sigma(t,x)-\sigma(t,y)}\le K\abs{x-y}^{1/2}$, 则轨道唯一性仍成立 (见~\cite{karatzas} 命题 5.2.13). 这适用于金融中的 Cox-Ingersoll-Ross 利率模型 $\dd X_t=\kappa(\theta-X_t)\dd t+\sigma\sqrt{X_t}\dd B_t$, 其扩散系数 $\sigma\sqrt{x}$ 在 $0$ 附近不是 Lipschitz 的.

(4) \emph{解是噪声的泛函.} 由构造, 每个 Picard 迭代 $X^k$ 是 $(\xi,B)$ 的 (同一个) 可测泛函, 所以解 $X$ 也是. 特别地, $X$ 的分布由 $\xi$ 的分布与 $(b,\sigma)$ 唯一决定, 与概率空间的选取无关.
\end{remark}

% ---------------------------------------------------------------------
\subsection{弱解}

弱解的概念允许我们处理系数不满足 Lipschitz 条件的方程. Girsanov 定理提供了构造弱解的一般方法: 先取一个布朗运动作为 ``解'', 再改变测度, 使它获得需要的漂移.

\begin{proposition}[用 Girsanov 定理构造弱解]\label{prop:weak-girsanov}
设 $b:[0,T]\times\R\to\R$ 是有界 Borel 函数, $x\in\R$. 则方程
\[
\dd X_t=b(t,X_t)\dd t+\dd W_t,\qquad X_0=x
\]
有弱解.
\end{proposition}

\begin{proof}
在某个概率空间上取标准布朗运动 $B$ 及其扩充信息流, 令 $X_t=x+B_t$. 过程 $\theta_t=b(t,X_t)$ 有界且循序可测, 由命题~\ref{prop:novikov} (1), $Z_t=\exp(\int_0^t\theta\dd B-\frac{1}{2}\int_0^t\theta^2\dd s)$ 在 $[0,T]$ 上是鞅. 令 $\widetilde{\P}=Z_T\cdot\P$. 由 Girsanov 定理~\ref{thm:girsanov}, $W_t=B_t-\int_0^tb(s,X_s)\dd s$ 是 $\widetilde{\P}$ 下的布朗运动, 而
\[
X_t=x+B_t=x+\int_0^tb(s,X_s)\dd s+W_t.
\]
所以 $(\Omega,\mathcal{F},\widetilde{\P},\{\mathcal{F}_t\},W,X)$ 是一个弱解 ($\widetilde{\P}$ 与 $\P$ 等价, 信息流在 $\widetilde{\P}$ 下仍满足通常条件).
\end{proof}

这里 $b$ 可以是不连续的, 例如 $b(x)=-\sgn(x)$ (总是把过程拉向原点的 ``常数力''). 可以证明, 这样构造的弱解在分布意义下唯一 (见~\cite{karatzas} 命题 5.3.10); 事实上, Zvonkin 与 Veretennikov 证明了此时还有强解. 下面的例子说明, 强解确实可能不存在.

\begin{example}[Tanaka 方程]\label{ex:tanaka-sde}
考虑方程
\[
\dd X_t=\sgn(X_t)\dd W_t,\qquad X_0=0,
\]
其中 $\sgn(x)=1$ ($x>0$), $-1$ ($x\le0$), 与例~\ref{ex:levy-applications} 相同.

(1) \emph{弱解存在.} 设 $X$ 是标准布朗运动, 令 $W_t=\int_0^t\sgn(X_s)\dd X_s$. 由 L\'evy 刻画 (例~\ref{ex:levy-applications} (1)), $W$ 是布朗运动, 且 $\int_0^t\sgn(X_s)\dd W_s=\int_0^t\sgn(X_s)^2\dd X_s=X_t$. 所以 $(X,W)$ 是弱解.

(2) \emph{分布唯一.} 若 $(X,W)$ 是任一弱解, 则 $X=\int\sgn(X)\dd W$ 满足 $\ip{X}_t=\int_0^t\sgn(X_s)^2\dd s=t$, 由 L\'evy 刻画~\ref{thm:levy-char}, $X$ 是布朗运动. 所以一切弱解有相同的分布.

(3) \emph{轨道唯一性不成立.} 设 $(X,W)$ 是弱解. 由于 $\sgn(-x)=-\sgn(x)$ 对 $x\ne0$ 成立, 而 $X$ 是布朗运动, $\int_0^T\1_{\{X_s=0\}}\dd s=0$ $\as$, 所以 $\sgn(-X_s)=-\sgn(X_s)$ 对几乎每个 $s$ 成立, 从而 $\int_0^t\sgn(-X_s)\dd W_s=-X_t$: $-X$ 也是 (同一个 $W$ 驱动的) 解.

(4) \emph{强解不存在.} 设 $(X,W)$ 是弱解. 由 Tanaka 公式~\ref{thm:tanaka} (应用于布朗运动 $X$),
\[
W_t=\int_0^t\sgn(X_s)\dd X_s=\abs{X_t}-L_t,\qquad L_t=\lim_{\varepsilon\downarrow0}\frac{1}{2\varepsilon}\int_0^t\1_{\{\abs{X_s}<\varepsilon\}}\dd s.
\]
右端只依赖于 $\abs{X}$, 所以 $\mathcal{F}^W_t\subseteq\mathcal{F}^{\abs{X}}_t$ (扩充后). 若 $X$ 是强解, 即 $X_t$ 是 $\mathcal{F}^W_t$-可测的, 则 $X_t$ 是 $\mathcal{F}^{\abs{X}}_t$-可测的, 即存在可测泛函 $G$ 使 $X_t=G(\abs{X_s}:s\le t)$ $\as$ 但 $-X$ 也是布朗运动, $(X,\abs{X})$ 与 $(-X,\abs{X})$ 同分布, 所以 $-X_t=G(\abs{X_s}:s\le t)$ 也几乎必然成立, 从而 $X_t=0$ $\as$, 对 $t>0$ 矛盾.
\end{example}

\begin{intuition}
强解是 ``可以由噪声计算出来'' 的解: 知道了驱动的布朗运动 $W$ 的整条轨道, 就知道了 $X$. 对 Tanaka 方程, $W$ 只记录了 $\abs{X}$ 的信息, 每当 $X$ 回到零点, 它下一段 ``向上走还是向下走'' 是 $W$ 无法决定的额外随机性. 弱解则只要求 $(X,W)$ 的联合分布正确, 允许 $X$ 携带 $W$ 之外的随机性. Yamada-Watanabe 定理把两种概念联系起来: 弱解存在 + 轨道唯一性 $\Rightarrow$ 强解存在且分布唯一 (见~\cite{karatzas} 推论 5.3.23).
\end{intuition}

% =====================================================================
\section{扩散过程: 马尔可夫性与生成元}
\label{sec:6-3}

本节假设系数不依赖于时间: $b=b(x)$, $\sigma=\sigma(x)$, 且满足 \eqref{eq:H1}\eqref{eq:H2} (于是 $b,\sigma$ 连续). 对 $s\ge0$ 与 $\mathcal{F}_s$-可测的 $\eta\in L^2$, 记 $X^{s,\eta}=\{X^{s,\eta}_t\}_{t\ge s}$ 为从时刻 $s$、位置 $\eta$ 出发的解:
\[
X^{s,\eta}_t=\eta+\int_s^tb(X^{s,\eta}_u)\dd u+\int_s^t\sigma(X^{s,\eta}_u)\dd B_u,\qquad t\ge s.
\]
定理~\ref{thm:sde-existence} 与命题~\ref{prop:sde-continuity} 对初始时刻 $s$ 同样成立 (证明只需平移时间). 记 $X^x=X^{0,x}$, 并用 $\P_x$, $\E_x$ 表示与 $X^x$ 有关的概率与期望. 时齐方程的解称为 (It\^o) \textbf{扩散过程}.

% ---------------------------------------------------------------------
\subsection{马尔可夫性}

\begin{lemma}\label{lem:sde-shift}
设 $s\ge0$, $y,z\in\R$, 记 $W_u=B_{s+u}-B_s$.
\begin{enumerate}
\item $\{X^{s,y}_{s+u}\}_{u\ge0}$ 关于 $\sigma(W_u:u\ge0)$ (扩充后) 可测, 从而与 $\mathcal{F}_s$ 独立;
\item $\{X^{s,y}_{s+u}\}_{u\ge0}$ 与 $\{X^y_u\}_{u\ge0}$ 有相同的有限维分布;
\item 对有界连续函数 $f$ 与 $h\ge0$, $y\mapsto P_hf(y):=\E f(X^y_h)$ 是连续函数.
\end{enumerate}
\end{lemma}

\begin{proof}
(1)(2) 由定理~\ref{thm:bm-markov}, $W$ 是与 $\mathcal{F}_s$ 独立的布朗运动. $Y_u=X^{s,y}_{s+u}$ 满足以 $W$ 为驱动的方程 $Y_u=y+\int_0^ub(Y_r)\dd r+\int_0^u\sigma(Y_r)\dd W_r$ (随机积分的时间平移可由 Riemann 和逼近验证, 命题~\ref{prop:riemann-approx}). 对这个方程作 Picard 迭代: 每个迭代都是连续过程, 由命题~\ref{prop:riemann-approx}, 关于 $W$ 的随机积分是 $W$ 与被积过程的 Riemann 和的依概率极限; 对 $k$ 归纳可知, 每个迭代都关于 $\sigma(W)$ 可测, 而且它作为 $W$ 的泛函, 与以 $B$ 为驱动、从 $y$ 出发的迭代是 ``同一个'' 泛函 (Riemann 和的表达式相同). 由于 $W$ 与 $B$ 同分布, 二者的迭代 (连同驱动的布朗运动) 有相同的有限维分布; 取极限 (几乎必然收敛保持同分布) 即得 (1)(2).

(3) 由命题~\ref{prop:sde-continuity}, $z\to y$ 时 $X^z_h\to X^y_h$ 于 $L^2$, 从而依概率收敛. 由于 $f$ 有界连续, $f(X^z_h)\to f(X^y_h)$ 依概率成立, 由有界收敛定理 $P_hf(z)\to P_hf(y)$.
\end{proof}

\begin{theorem}[扩散过程的马尔可夫性]\label{thm:sde-markov}
设 $X=X^{0,\xi}$, $\xi\in L^2(\mathcal{F}_0)$. 则对 $s,h\ge0$ 与有界 Borel 函数 $f$,
\[
\E\big[f(X_{s+h})\,\big|\,\mathcal{F}_s\big]=P_hf(X_s),\qquad P_hf(y)=\E f(X^y_h).
\]
所以 $X$ 是时齐马尔可夫过程, 转移概率为 $P_h(y,A)=\P(X^y_h\in A)$.
\end{theorem}

\begin{proof}
\emph{第一步.} 由唯一性, $X_{s+\cdot}$ 就是从时刻 $s$、位置 $X_s$ 出发的解: $X_t=X^{s,X_s}_t$ ($t\ge s$), 因为二者满足同一个方程.

\emph{第二步: 初值取可数多个值.} 设 $\eta=\sum_jy_j\1_{A_j}$, 其中 $\{A_j\}\subseteq\mathcal{F}_s$ 是 $\Omega$ 的可数分割. 令 $Y=\sum_j\1_{A_j}X^{s,y_j}$. 在 $A_j$ 上, $Y=X^{s,y_j}$, 由局部性引理~\ref{lem:locality}, 在 $A_j$ 上 $\int_s^t\sigma(Y_u)\dd B_u=\int_s^t\sigma(X^{s,y_j}_u)\dd B_u$; 所以 $Y$ 满足初值为 $\eta$ 的方程, 由唯一性 $Y=X^{s,\eta}$. 于是, 由提出已知因子、引理~\ref{lem:sde-shift} (1)(2),
\begin{align*}
\E\big[f(X^{s,\eta}_{s+h})\,\big|\,\mathcal{F}_s\big]&=\sum_j\1_{A_j}\E\big[f(X^{s,y_j}_{s+h})\,\big|\,\mathcal{F}_s\big]=\sum_j\1_{A_j}\E f(X^{s,y_j}_{s+h})\\
&=\sum_j\1_{A_j}P_hf(y_j)=P_hf(\eta).
\end{align*}

\emph{第三步: 一般初值, $f$ 有界连续.} 令 $\eta_n=2^{-n}\lfloor2^nX_s\rfloor$, 则 $\abs{\eta_n-X_s}\le2^{-n}$. 由命题~\ref{prop:sde-continuity} (初始时刻为 $s$), $X^{s,\eta_n}_{s+h}\to X^{s,X_s}_{s+h}=X_{s+h}$ 于 $L^2$. 由有界收敛定理, 第二步等式的左端在 $L^1$ 中收敛到 $\E[f(X_{s+h})\,|\,\mathcal{F}_s]$; 由引理~\ref{lem:sde-shift} (3), 右端 $P_hf(\eta_n)\to P_hf(X_s)$.

\emph{第四步: $f$ 有界 Borel.} 满足结论 (且 $P_hf$ 为 Borel 函数) 的有界 Borel 函数 $f$ 全体是包含常数的线性空间, 由单调收敛定理与条件单调收敛定理对有界单调极限封闭, 由第三步包含有界连续函数. 由引理~\ref{lem:cont-to-borel} ($n=1$), 它包含全体有界 Borel 函数.
\end{proof}

与布朗运动相同 (定理~\ref{thm:bm-markov} 与定理~\ref{thm:strong-markov}), 扩散过程对整条未来轨道的泛函也有马尔可夫性, 并且有强马尔可夫性: 对停时 $\tau$, 在 $\{\tau<\infty\}$ 上 $\E[F(X_{\tau+\cdot})\,|\,\mathcal{F}_\tau]=\E_{X_\tau}[F]$ (见~\cite{karatzas} 定理 5.4.20). 引理~\ref{lem:sde-shift} (3) 说明 $P_h$ 把有界连续函数映为有界连续函数, 这称为 \textbf{Feller 性}.

% ---------------------------------------------------------------------
\subsection{生成元与 Dynkin 公式}

马尔可夫过程的转移概率满足 Chapman-Kolmogorov 方程 $P_{t+s}=P_tP_s$ (定理~\ref{thm:chapman-kolmogorov}), 所以整个半群 $\{P_t\}$ 由它在 $t=0$ 附近的行为决定, 就像 $e^{tA}$ 由 $A$ 决定一样.

\begin{definition}[生成元]\label{def:generator}
扩散过程的\textbf{无穷小生成元}定义为
\[
\mathcal{A}f(x)=\lim_{t\downarrow0}\frac{\E_xf(X_t)-f(x)}{t},
\]
定义域为使极限对一切 $x$ 存在的函数 $f$. 对二阶可微的 $f$, 记
\begin{equation}\label{eq:L-operator}
Lf(x)=b(x)f'(x)+\frac{1}{2}\sigma(x)^2f''(x).
\end{equation}
\end{definition}

记 $C^2_c$ 为二阶连续可微且支撑紧的函数全体.

\begin{proposition}[Dynkin 公式]\label{prop:dynkin}
设 $f\in C^2_c$, $\tau$ 是停时, $\E_x\tau<\infty$. 则
\[
\E_xf(X_\tau)=f(x)+\E_x\int_0^\tau Lf(X_s)\dd s.
\]
\end{proposition}

\begin{proof}
由 It\^o 公式, $f(X_t)=f(x)+\int_0^tLf(X_s)\dd s+\int_0^tf'(X_s)\sigma(X_s)\dd B_s$. 由于 $f'$ 支撑紧而 $\sigma$ 连续, $f'\sigma$ 有界, 所以随机积分是鞅 (定理~\ref{thm:ito-properties}). 对有界停时 $\tau\wedge n$ 应用可选停时定理 (推论~\ref{cor:ct-ost-bounded}),
\[
\E_xf(X_{\tau\wedge n})=f(x)+\E_x\int_0^{\tau\wedge n}Lf(X_s)\dd s.
\]
令 $n\to\infty$: 左端由有界收敛定理收敛 ($\tau<\infty$ $\as$); 右端被积函数 $Lf$ 有界, $\abs{\int_0^{\tau\wedge n}Lf\dd s}\le\norm{Lf}_\infty\tau\in L^1$, 由控制收敛定理收敛.
\end{proof}

\begin{corollary}[生成元的形式]\label{cor:generator}
对 $f\in C^2_c$, $\mathcal{A}f=Lf$.
\end{corollary}

\begin{proof}
在 Dynkin 公式中取 $\tau=t$, 由 Fubini 定理,
\[
\frac{\E_xf(X_t)-f(x)}{t}=\frac{1}{t}\int_0^t\E_xLf(X_s)\dd s.
\]
$Lf$ 有界连续, $X$ 的轨道连续, 由有界收敛定理, $s\mapsto\E_xLf(X_s)$ 在 $s=0$ 处连续, 其值为 $Lf(x)$. 所以右端趋于 $Lf(x)$.
\end{proof}

\begin{intuition}
生成元把扩散过程归结为两个 ``局部特征'': 在状态 $x$, 过程在短时间 $\Delta t$ 内的增量近似满足
\[
\E[\Delta X\,|\,X_t=x]\approx b(x)\Delta t,\qquad\Var(\Delta X\,|\,X_t=x)\approx\sigma(x)^2\Delta t.
\]
$Lf$ 中的一阶项来自平均位移, 二阶项来自方差 (这正是 It\^o 公式的修正项). 对布朗运动, $L=\frac{1}{2}\frac{\dd^2}{\dd x^2}$ (第~4~章的直观); 对 OU 过程 $\dd X=-\kappa X\dd t+\sigma\dd B$, $L=-\kappa x\frac{\dd}{\dd x}+\frac{\sigma^2}{2}\frac{\dd^2}{\dd x^2}$; 对几何布朗运动, $L=\mu x\frac{\dd}{\dd x}+\frac{\sigma^2x^2}{2}\frac{\dd^2}{\dd x^2}$. 多维情形 $\dd\rvec{X}=\rvec{b}\dd t+\mat{\Sigma}\dd\rvec{B}$, 生成元为 $L=\sum_ib_i\partial_i+\frac{1}{2}\sum_{i,j}a_{ij}\partial_{ij}$, 其中 $\mat{A}=(a_{ij})=\mat{\Sigma}\mat{\Sigma}^\T$ 是 ``扩散矩阵''.
\end{intuition}

% ---------------------------------------------------------------------
\subsection{Kolmogorov 方程与平稳分布}

第~4~章证明了布朗运动的转移密度满足热方程 (定理~\ref{thm:heat-equation}). 对一般的扩散过程, 热方程中的 $\frac{1}{2}\partial_{xx}$ 换成生成元 $L$ 或其伴随算子.

\begin{theorem}[Kolmogorov 向后方程]\label{thm:backward}
设 $f$ 有界连续, $u\in C^{1,2}((0,\infty)\times\R)\cap C([0,\infty)\times\R)$ 有界, 且对每个 $T$, $\partial_xu$ 在 $(0,T]\times\R$ 上有界. 若
\[
\partial_tu=Lu\quad\text{于 }(0,\infty)\times\R,\qquad u(0,x)=f(x),
\]
则 $u(t,x)=\E_xf(X_t)$.
\end{theorem}

\begin{proof}
固定 $t>0$ 与 $x$, 令 $v(s,y)=u(t-s,y)$ ($0\le s<t$), 则 $\partial_sv+Lv=0$. 对 $0<\varepsilon<t$, 在 $[0,t-\varepsilon]$ 上对 $v(s,X_s)$ ($X=X^x$) 应用 It\^o 公式,
\[
v(t-\varepsilon,X_{t-\varepsilon})=v(0,x)+\int_0^{t-\varepsilon}(\partial_sv+Lv)(s,X_s)\dd s+\int_0^{t-\varepsilon}\partial_yv(s,X_s)\sigma(X_s)\dd B_s.
\]
中间的积分为零. 由于 $\partial_yv$ 有界而 $\E\int_0^t\sigma(X_s)^2\dd s<\infty$ (由 \eqref{eq:H2} 与 \eqref{eq:sde-moment}), 随机积分的期望为零, 所以 $\E u(\varepsilon,X_{t-\varepsilon})=u(t,x)$. 令 $\varepsilon\downarrow0$, 由 $u$ 在 $[0,\infty)\times\R$ 上的连续性与轨道连续性, $u(\varepsilon,X_{t-\varepsilon})\to u(0,X_t)=f(X_t)$; 由有界收敛定理, $u(t,x)=\E_xf(X_t)$.
\end{proof}

定理说: 向后方程的 (足够正则的) 解必定由扩散过程表示, 特别地解是唯一的. 反过来, 若 $b,\sigma$ 与 $f$ 足够光滑 (例如具有有界的各阶导数), 则 $u(t,x)=\E_xf(X_t)$ 确实是 $C^{1,2}$ 的, 从而满足向后方程 (见~\cite{karatzas} 5.7 节). 它描述期望 $\E_xf(X_t)$ 如何随\emph{起点} $x$ 变化. 与之对偶, 向前方程描述 $X_t$ 的\emph{分布}如何随时间演化.

\begin{theorem}[Kolmogorov 向前方程 (Fokker-Planck 方程)]\label{thm:forward}
设 $b\in C^1$, $\sigma\in C^2$, $X_t$ ($X_0=x$) 对 $t>0$ 有密度 $p(t,y)$, $p\in C^{1,2}((0,\infty)\times\R)$, 并且可以在积分号下对 $t$ 求导 (即 $\frac{\dd}{\dd t}\int f(y)p(t,y)\dd y=\int f(y)\partial_tp(t,y)\dd y$ 对 $f\in C^2_c$ 成立). 则
\[
\partial_tp=L^*p:=-\partial_y\big(b\,p\big)+\frac{1}{2}\partial_{yy}\big(\sigma^2p\big).
\]
\end{theorem}

\begin{proof}
设 $f\in C^2_c$. 由 Dynkin 公式 (取 $\tau=t$), $\E_xf(X_t)=f(x)+\int_0^t\E_xLf(X_s)\dd s$, 被积函数关于 $s$ 连续, 所以
\[
\int f(y)\partial_tp(t,y)\dd y=\frac{\dd}{\dd t}\E_xf(X_t)=\E_xLf(X_t)=\int\big(b(y)f'(y)+\tfrac{1}{2}\sigma(y)^2f''(y)\big)p(t,y)\dd y.
\]
由于 $f$ 支撑紧, 分部积分没有边界项: $\int bf'p\dd y=-\int f\,\partial_y(bp)\dd y$, $\int\sigma^2f''p\dd y=\int f\,\partial_{yy}(\sigma^2p)\dd y$. 于是 $\int f\,(\partial_tp-L^*p)\dd y=0$ 对一切 $f\in C^2_c$ 成立. 连续函数若与一切 $C^2_c$ 函数的积分都为零, 则恒为零, 所以 $\partial_tp=L^*p$.
\end{proof}

\begin{figure}[htbp]
\centering
\begin{tikzpicture}
\begin{axis}[SAaxis, axis x line=bottom, axis y line=left, width=12cm, height=5.6cm, xmin=-2.6, xmax=4.4, ymin=0, ymax=1.42,
  xtick={-2,-1,0,1,2,3,4}, ytick={0,0.5,1}, xlabel={$y$}, samples=200, domain=-2.6:4.4,
  legend style={font=\small,draw=none,fill=none,at={(0.5,-0.22)},anchor=north}, legend columns=5, legend cell align=left]
  \addplot[SARed,thick] {exp(-(x-3*exp(-0.1))^2/(1-exp(-0.2)))/sqrt(pi*(1-exp(-0.2)))};
  \addplot[SAYellowText,thick] {exp(-(x-3*exp(-0.5))^2/(1-exp(-1)))/sqrt(pi*(1-exp(-1)))};
  \addplot[SAGreenText,thick] {exp(-(x-3*exp(-1))^2/(1-exp(-2)))/sqrt(pi*(1-exp(-2)))};
  \addplot[SABlue,thick] {exp(-(x-3*exp(-3))^2/(1-exp(-6)))/sqrt(pi*(1-exp(-6)))};
  \addplot[SADark,thick,dashed] {exp(-x^2)/sqrt(pi)};
  \legend{$t=0.1$, $t=0.5$, $t=1$, $t=3$, 平稳分布}
\end{axis}
\end{tikzpicture}
\caption{Fokker-Planck 方程的解: OU 过程 $\dd X_t=-X_t\dd t+\dd B_t$ 从 $x=3$ 出发时 $X_t$ 的密度 $N\big(3e^{-t},\frac12(1-e^{-2t})\big)$. 密度从点质量出发, 一边扩散一边向原点漂移, 最终收敛到平稳分布 $N(0,\frac12)$ (虚线)}
\label{fig:fokker-planck}
\end{figure}

若概率分布 $\pi$ 满足: 当 $X_0\sim\pi$ (与 $B$ 独立) 时, 对一切 $t$ 都有 $X_t\sim\pi$, 则称 $\pi$ 为\textbf{平稳分布}. 由向前方程, 平稳分布的密度应满足 $L^*\pi=0$. 在一维情形, 若进一步要求 ``概率流'' $b\pi-\frac{1}{2}(\sigma^2\pi)'$ 为零 (这在 $\pi$ 及其导数在 $\pm\infty$ 处趋于零时是自然的), 则得到一阶方程 $(\frac{1}{2}\sigma^2\pi)'=b\pi$, 其解为
\begin{equation}\label{eq:stationary-density}
\pi(y)=\frac{C}{\sigma(y)^2}\exp\Big(\int_0^y\frac{2b(z)}{\sigma(z)^2}\dd z\Big),
\end{equation}
其中 $C$ 是归一化常数 (要求右端可积). 直接验证: 令 $m=\frac{1}{2}\sigma^2\pi=\frac{C}{2}\exp(\int_0^y\frac{2b}{\sigma^2})$, 则 $m'=\frac{2b}{\sigma^2}m=b\pi$, 从而 $L^*\pi=-(b\pi)'+m''=0$. 可以证明, 在适当的条件下 (例如 $\sigma$ 有正下界且 \eqref{eq:stationary-density} 可积), 这就是唯一的平稳分布, 且对任意初值, $X_t$ 的分布收敛到 $\pi$.

\begin{example}[OU 过程与 Langevin 方程]\label{ex:stationary}
(1) 对 OU 过程 $\dd X=\kappa(\theta-X)\dd t+\sigma\dd B$, $\frac{2b}{\sigma^2}=\frac{2\kappa(\theta-y)}{\sigma^2}$, 由 \eqref{eq:stationary-density}, $\pi\propto e^{-\kappa(y-\theta)^2/\sigma^2}$, 即 $\pi=N(\theta,\frac{\sigma^2}{2\kappa})$. 这可以直接验证: 若 $X_0\sim N(\theta,\frac{\sigma^2}{2\kappa})$ 与 $B$ 独立, 由例~\ref{ex:linear-examples} (2), $X_t$ 是正态的, 均值为 $\theta$, 方差为 $e^{-2\kappa t}\frac{\sigma^2}{2\kappa}+\frac{\sigma^2}{2\kappa}(1-e^{-2\kappa t})=\frac{\sigma^2}{2\kappa}$.

(2) 设 $V$ 是光滑的 ``势能'' 函数, 考虑 \textbf{Langevin 方程} $\dd X_t=-V'(X_t)\dd t+\sqrt{2}\dd B_t$: 粒子沿势能下降的方向运动, 同时受到热噪声的扰动. 由 \eqref{eq:stationary-density}, 平稳分布为 \textbf{Gibbs 分布}
\[
\pi(y)=\frac{1}{Z}e^{-V(y)},\qquad Z=\int_{\R}e^{-V(y)}\dd y.
\]
当 $V$ 有多个极小点时 (例如双势阱 $V(y)=(y^2-1)^2$), 过程大部分时间停留在某个势阱附近, 偶尔借助噪声越过势垒跳到另一个势阱 (图~\ref{fig:langevin}). 反过来, 若想从一个已知 (至多相差一个常数因子) 的密度 $\pi$ 中抽样, 可以令 $V=-\log\pi$, 模拟 Langevin 方程足够长时间: 这是 Langevin Monte Carlo 方法的基础, 也是现代扩散生成模型的出发点之一.
\end{example}

\begin{figure}[htbp]
\centering
\begin{tikzpicture}
\begin{axis}[SAaxis, axis x line=bottom, axis y line=left, width=6.2cm, height=5.2cm, xmin=-2.1, xmax=2.1, ymin=0, ymax=1.6,
  xtick={-2,-1,0,1,2}, ytick={0,0.5,1,1.5}, xlabel={$y$}, samples=161, domain=-2:2,
  legend style={font=\scriptsize,draw=none,fill=none,at={(0.5,1.02)},anchor=south}, legend columns=2, legend cell align=left]
  \addplot[SADark,thick] {(x^2-1)^2};
  \addplot[SARed,thick,fill=SARed!12] {exp(-(x^2-1)^2)/1.97373} \closedcycle;
  \legend{$V(y)$, $\pi(y)$}
\end{axis}
\end{tikzpicture}\hfill
\begin{tikzpicture}
\begin{axis}[SAaxis, axis x line=bottom, axis y line=left, width=9cm, height=5.2cm, xmin=0, xmax=203, ymin=-2.1, ymax=2.1,
  xtick={0,50,100,150,200}, ytick={-1,0,1}, xlabel={$t$}]
  \addplot[SADark!25] coordinates {(0,0) (200,0)};
  \addplot[SABlue,very thin] coordinates {(0.000,-1.000) (0.250,-0.826) (0.500,-0.306) (0.750,-0.724) (1.000,-0.592) (1.250,0.757) (1.500,0.760) (1.750,1.278) (2.000,1.055) (2.250,1.144) (2.500,0.969) (2.750,1.230) (3.000,1.443) (3.250,0.577) (3.500,1.312) (3.750,1.028) (4.000,1.450) (4.250,1.280) (4.500,1.055) (4.750,0.932) (5.000,1.130) (5.250,0.572) (5.500,0.476) (5.750,0.522) (6.000,0.880) (6.250,0.493) (6.500,0.833) (6.750,0.934) (7.000,0.082) (7.250,0.841) (7.500,1.437) (7.750,1.173) (8.000,1.676) (8.250,1.006) (8.500,1.286) (8.750,0.432) (9.000,1.208) (9.250,1.375) (9.500,1.048) (9.750,1.058) (10.000,1.309) (10.250,1.232) (10.500,0.820) (10.750,0.475) (11.000,0.436) (11.250,-0.645) (11.500,-0.221) (11.750,-0.540) (12.000,-0.922) (12.250,-0.592) (12.500,-0.266) (12.750,-1.062) (13.000,-1.002) (13.250,-0.123) (13.500,1.088) (13.750,1.199) (14.000,1.007) (14.250,1.008) (14.500,0.669) (14.750,0.023) (15.000,-1.223) (15.250,-0.546) (15.500,-0.997) (15.750,-1.043) (16.000,-0.863) (16.250,0.192) (16.500,1.342) (16.750,0.603) (17.000,1.413) (17.250,1.052) (17.500,1.398) (17.750,1.075) (18.000,0.685) (18.250,0.378) (18.500,0.967) (18.750,1.613) (19.000,-0.181) (19.250,0.541) (19.500,0.809) (19.750,0.219) (20.000,1.127) (20.250,1.072) (20.500,1.143) (20.750,0.456) (21.000,1.080) (21.250,0.704) (21.500,1.485) (21.750,1.230) (22.000,0.260) (22.250,0.638) (22.500,0.042) (22.750,-0.679) (23.000,-0.801) (23.250,-1.043) (23.500,-1.240) (23.750,-0.575) (24.000,-0.246) (24.250,0.597) (24.500,1.048) (24.750,-0.201) (25.000,-0.872) (25.250,-1.169) (25.500,-1.195) (25.750,-0.787) (26.000,-0.386) (26.250,0.009) (26.500,-0.008) (26.750,-0.427) (27.000,-1.057) (27.250,-0.341) (27.500,-0.404) (27.750,-0.231) (28.000,-0.223) (28.250,-0.462) (28.500,-0.492) (28.750,-0.631) (29.000,-1.015) (29.250,-0.106) (29.500,1.487) (29.750,0.890) (30.000,-0.475) (30.250,-1.051) (30.500,-1.070) (30.750,-0.869) (31.000,-0.807) (31.250,-1.230) (31.500,-0.941) (31.750,-0.477) (32.000,-1.088) (32.250,-0.438) (32.500,-0.131) (32.750,-0.808) (33.000,-1.263) (33.250,-1.121) (33.500,-1.393) (33.750,-0.964) (34.000,-1.259) (34.250,-0.170) (34.500,-0.372) (34.750,0.928) (35.000,1.486) (35.250,1.121) (35.500,0.754) (35.750,1.147) (36.000,0.501) (36.250,0.747) (36.500,1.293) (36.750,0.701) (37.000,1.452) (37.250,1.053) (37.500,0.529) (37.750,-0.187) (38.000,-0.882) (38.250,-1.117) (38.500,0.258) (38.750,0.462) (39.000,0.447) (39.250,-0.426) (39.500,0.566) (39.750,1.057) (40.000,0.445) (40.250,0.290) (40.500,0.021) (40.750,0.641) (41.000,0.581) (41.250,0.015) (41.500,-0.335) (41.750,-0.593) (42.000,-1.044) (42.250,-1.113) (42.500,-0.867) (42.750,-0.472) (43.000,0.438) (43.250,0.529) (43.500,1.298) (43.750,0.903) (44.000,0.283) (44.250,-0.894) (44.500,-0.833) (44.750,-0.153) (45.000,0.463) (45.250,0.249) (45.500,0.563) (45.750,1.131) (46.000,0.940) (46.250,0.124) (46.500,1.223) (46.750,1.343) (47.000,0.930) (47.250,0.117) (47.500,0.001) (47.750,0.772) (48.000,0.616) (48.250,0.049) (48.500,0.146) (48.750,0.552) (49.000,1.326) (49.250,1.222) (49.500,0.855) (49.750,1.052) (50.000,1.249) (50.250,1.428) (50.500,1.196) (50.750,0.945) (51.000,1.231) (51.250,0.987) (51.500,1.018) (51.750,1.164) (52.000,0.557) (52.250,0.743) (52.500,1.042) (52.750,0.857) (53.000,0.260) (53.250,0.825) (53.500,0.897) (53.750,0.736) (54.000,0.952) (54.250,1.023) (54.500,0.528) (54.750,-0.409) (55.000,0.133) (55.250,-0.847) (55.500,-1.132) (55.750,-1.050) (56.000,-1.269) (56.250,-1.066) (56.500,-0.643) (56.750,-1.187) (57.000,-1.142) (57.250,-0.972) (57.500,-1.108) (57.750,-0.882) (58.000,-0.588) (58.250,-0.825) (58.500,-0.765) (58.750,-1.162) (59.000,-0.529) (59.250,-0.338) (59.500,-0.252) (59.750,-0.681) (60.000,-1.126) (60.250,-1.270) (60.500,-1.242) (60.750,-1.266) (61.000,-0.725) (61.250,-0.870) (61.500,-1.439) (61.750,-0.757) (62.000,-1.204) (62.250,-0.571) (62.500,-1.320) (62.750,-0.738) (63.000,-0.636) (63.250,-1.291) (63.500,-0.114) (63.750,0.070) (64.000,-0.401) (64.250,-0.579) (64.500,-0.574) (64.750,-0.358) (65.000,-1.066) (65.250,-0.943) (65.500,-1.187) (65.750,-0.633) (66.000,0.242) (66.250,0.625) (66.500,0.826) (66.750,0.804) (67.000,0.414) (67.250,0.725) (67.500,0.694) (67.750,0.872) (68.000,0.778) (68.250,0.857) (68.500,0.198) (68.750,0.230) (69.000,-0.803) (69.250,-0.920) (69.500,-1.430) (69.750,-0.845) (70.000,-0.369) (70.250,-0.805) (70.500,-1.132) (70.750,-0.809) (71.000,-1.188) (71.250,-0.631) (71.500,-0.937) (71.750,-1.040) (72.000,-0.660) (72.250,-1.013) (72.500,-0.742) (72.750,-1.494) (73.000,-1.100) (73.250,-1.104) (73.500,-0.480) (73.750,-0.492) (74.000,0.990) (74.250,0.559) (74.500,0.988) (74.750,1.117) (75.000,1.198) (75.250,0.969) (75.500,1.012) (75.750,0.702) (76.000,0.847) (76.250,-0.043) (76.500,0.489) (76.750,1.074) (77.000,-0.204) (77.250,-0.738) (77.500,-0.677) (77.750,-1.286) (78.000,-0.937) (78.250,-0.998) (78.500,-0.856) (78.750,-1.052) (79.000,-1.179) (79.250,-0.986) (79.500,-0.986) (79.750,-0.538) (80.000,-1.035) (80.250,-0.744) (80.500,-0.553) (80.750,-0.750) (81.000,-0.280) (81.250,-1.239) (81.500,-0.733) (81.750,-1.470) (82.000,-0.556) (82.250,1.106) (82.500,1.097) (82.750,1.467) (83.000,1.129) (83.250,-0.061) (83.500,-0.539) (83.750,-0.345) (84.000,-1.061) (84.250,-0.665) (84.500,-1.530) (84.750,-1.165) (85.000,-1.112) (85.250,-1.248) (85.500,-1.297) (85.750,-0.294) (86.000,0.431) (86.250,0.983) (86.500,1.304) (86.750,0.508) (87.000,0.588) (87.250,0.447) (87.500,1.218) (87.750,1.338) (88.000,1.269) (88.250,0.655) (88.500,1.013) (88.750,0.223) (89.000,-0.139) (89.250,0.044) (89.500,-0.356) (89.750,-1.023) (90.000,-0.926) (90.250,-0.640) (90.500,-1.205) (90.750,0.079) (91.000,-0.886) (91.250,-0.968) (91.500,-0.503) (91.750,0.391) (92.000,0.834) (92.250,-0.166) (92.500,-0.776) (92.750,-0.544) (93.000,-0.690) (93.250,-0.350) (93.500,-0.442) (93.750,0.985) (94.000,1.024) (94.250,0.242) (94.500,0.920) (94.750,0.825) (95.000,-0.064) (95.250,-0.656) (95.500,-0.340) (95.750,-0.441) (96.000,-1.212) (96.250,-0.531) (96.500,-0.765) (96.750,-0.899) (97.000,-1.271) (97.250,-1.157) (97.500,-1.097) (97.750,-0.856) (98.000,-0.096) (98.250,0.035) (98.500,1.029) (98.750,0.270) (99.000,0.541) (99.250,-0.367) (99.500,0.931) (99.750,1.191) (100.000,1.087) (100.250,1.095) (100.500,1.186) (100.750,1.262) (101.000,1.083) (101.250,0.964) (101.500,1.198) (101.750,0.967) (102.000,1.367) (102.250,1.193) (102.500,1.061) (102.750,1.102) (103.000,1.371) (103.250,0.265) (103.500,1.175) (103.750,0.904) (104.000,1.001) (104.250,0.027) (104.500,-0.937) (104.750,-0.663) (105.000,-1.074) (105.250,-0.941) (105.500,-0.870) (105.750,-0.784) (106.000,-0.438) (106.250,0.504) (106.500,0.507) (106.750,-0.136) (107.000,0.213) (107.250,1.320) (107.500,0.898) (107.750,1.546) (108.000,1.708) (108.250,0.151) (108.500,1.062) (108.750,1.176) (109.000,0.311) (109.250,0.386) (109.500,0.412) (109.750,0.468) (110.000,0.753) (110.250,0.972) (110.500,1.179) (110.750,1.235) (111.000,1.205) (111.250,0.896) (111.500,1.384) (111.750,0.954) (112.000,0.964) (112.250,1.438) (112.500,1.384) (112.750,1.321) (113.000,1.780) (113.250,-0.488) (113.500,-1.144) (113.750,-1.004) (114.000,-1.284) (114.250,-0.934) (114.500,-1.375) (114.750,-0.528) (115.000,-0.439) (115.250,-0.775) (115.500,-0.603) (115.750,-0.730) (116.000,-0.799) (116.250,-0.888) (116.500,-0.767) (116.750,-1.222) (117.000,-1.189) (117.250,-0.945) (117.500,-0.672) (117.750,-0.791) (118.000,-0.133) (118.250,-0.429) (118.500,-1.160) (118.750,-1.513) (119.000,-1.035) (119.250,-0.982) (119.500,-0.791) (119.750,-0.948) (120.000,-0.990) (120.250,-0.807) (120.500,-0.963) (120.750,-0.271) (121.000,-0.000) (121.250,-0.123) (121.500,-0.630) (121.750,-1.317) (122.000,-0.806) (122.250,-0.086) (122.500,-0.199) (122.750,0.659) (123.000,1.572) (123.250,1.210) (123.500,0.631) (123.750,0.243) (124.000,0.970) (124.250,1.092) (124.500,0.729) (124.750,1.071) (125.000,0.996) (125.250,0.765) (125.500,1.294) (125.750,1.250) (126.000,1.249) (126.250,0.888) (126.500,1.355) (126.750,1.153) (127.000,1.106) (127.250,0.771) (127.500,0.556) (127.750,0.363) (128.000,-0.304) (128.250,0.014) (128.500,-0.870) (128.750,-1.243) (129.000,-0.182) (129.250,-0.285) (129.500,-0.111) (129.750,-0.710) (130.000,-0.832) (130.250,-0.671) (130.500,-1.137) (130.750,-1.433) (131.000,-0.965) (131.250,-1.279) (131.500,-0.594) (131.750,-1.151) (132.000,-0.989) (132.250,-0.717) (132.500,-1.300) (132.750,-1.213) (133.000,-0.544) (133.250,-1.095) (133.500,-1.122) (133.750,-1.370) (134.000,-1.400) (134.250,-0.480) (134.500,-1.130) (134.750,-1.164) (135.000,-0.377) (135.250,-0.957) (135.500,-0.411) (135.750,-1.202) (136.000,-0.616) (136.250,-0.911) (136.500,-1.032) (136.750,-1.196) (137.000,-1.034) (137.250,-0.860) (137.500,-1.046) (137.750,-1.205) (138.000,-1.201) (138.250,-1.133) (138.500,-0.824) (138.750,-0.979) (139.000,-1.154) (139.250,-0.795) (139.500,-1.365) (139.750,-0.772) (140.000,-0.549) (140.250,-0.939) (140.500,-1.183) (140.750,-1.212) (141.000,0.018) (141.250,-0.318) (141.500,0.916) (141.750,0.737) (142.000,0.308) (142.250,-0.519) (142.500,-1.082) (142.750,-1.521) (143.000,-1.305) (143.250,-0.045) (143.500,-0.883) (143.750,-0.553) (144.000,-1.021) (144.250,-0.432) (144.500,-1.126) (144.750,-0.693) (145.000,-0.691) (145.250,0.131) (145.500,-0.432) (145.750,-0.933) (146.000,-0.717) (146.250,-0.738) (146.500,-0.852) (146.750,-0.652) (147.000,-0.930) (147.250,-1.362) (147.500,-1.461) (147.750,-0.800) (148.000,-1.299) (148.250,-0.954) (148.500,-1.239) (148.750,-1.019) (149.000,-1.111) (149.250,-1.124) (149.500,-0.987) (149.750,-0.873) (150.000,-0.993) (150.250,-1.179) (150.500,-0.895) (150.750,-0.620) (151.000,0.109) (151.250,0.241) (151.500,0.775) (151.750,-0.666) (152.000,-0.372) (152.250,-0.803) (152.500,0.226) (152.750,-0.464) (153.000,-0.716) (153.250,-0.860) (153.500,-0.938) (153.750,-0.474) (154.000,-1.468) (154.250,-0.974) (154.500,-1.169) (154.750,-1.107) (155.000,-1.034) (155.250,-1.186) (155.500,-0.668) (155.750,-1.109) (156.000,-0.509) (156.250,-0.884) (156.500,-0.893) (156.750,-0.717) (157.000,-1.054) (157.250,-0.948) (157.500,1.305) (157.750,1.708) (158.000,1.290) (158.250,1.639) (158.500,0.914) (158.750,1.102) (159.000,1.417) (159.250,1.135) (159.500,0.082) (159.750,0.168) (160.000,1.037) (160.250,0.934) (160.500,-0.190) (160.750,-0.645) (161.000,-1.358) (161.250,0.480) (161.500,-1.176) (161.750,-0.924) (162.000,-1.240) (162.250,-0.002) (162.500,-0.438) (162.750,-0.877) (163.000,-1.308) (163.250,-1.393) (163.500,-0.930) (163.750,-1.099) (164.000,-0.562) (164.250,-0.191) (164.500,0.886) (164.750,1.057) (165.000,0.896) (165.250,1.376) (165.500,1.009) (165.750,0.804) (166.000,0.807) (166.250,0.464) (166.500,1.135) (166.750,0.517) (167.000,1.370) (167.250,0.670) (167.500,1.440) (167.750,1.110) (168.000,1.290) (168.250,1.108) (168.500,0.711) (168.750,0.513) (169.000,0.538) (169.250,1.295) (169.500,0.967) (169.750,0.779) (170.000,0.477) (170.250,0.667) (170.500,0.870) (170.750,1.300) (171.000,1.082) (171.250,1.064) (171.500,1.275) (171.750,0.640) (172.000,1.252) (172.250,1.262) (172.500,1.160) (172.750,0.768) (173.000,-0.051) (173.250,0.092) (173.500,1.408) (173.750,0.969) (174.000,0.767) (174.250,0.866) (174.500,0.294) (174.750,0.980) (175.000,0.608) (175.250,-0.007) (175.500,-0.385) (175.750,-0.227) (176.000,-0.307) (176.250,-0.708) (176.500,-0.433) (176.750,0.061) (177.000,-0.000) (177.250,-0.278) (177.500,-0.630) (177.750,-1.054) (178.000,-0.960) (178.250,-0.951) (178.500,-0.373) (178.750,0.018) (179.000,-1.348) (179.250,0.122) (179.500,0.932) (179.750,1.389) (180.000,1.273) (180.250,1.264) (180.500,1.100) (180.750,1.159) (181.000,1.231) (181.250,0.641) (181.500,1.081) (181.750,0.650) (182.000,0.887) (182.250,0.709) (182.500,0.714) (182.750,0.939) (183.000,0.574) (183.250,0.371) (183.500,0.944) (183.750,1.439) (184.000,0.860) (184.250,-1.050) (184.500,-1.079) (184.750,-1.108) (185.000,-0.598) (185.250,-0.630) (185.500,-0.341) (185.750,-1.349) (186.000,-0.870) (186.250,-0.937) (186.500,-0.979) (186.750,-0.701) (187.000,-0.177) (187.250,-0.834) (187.500,-1.209) (187.750,-0.919) (188.000,0.432) (188.250,0.844) (188.500,0.972) (188.750,1.237) (189.000,1.112) (189.250,1.216) (189.500,0.395) (189.750,0.091) (190.000,-0.417) (190.250,-0.755) (190.500,-1.255) (190.750,-1.779) (191.000,-0.950) (191.250,-1.144) (191.500,-0.849) (191.750,-1.019) (192.000,-1.359) (192.250,-0.609) (192.500,-0.042) (192.750,-0.079) (193.000,-1.617) (193.250,-0.808) (193.500,-0.838) (193.750,-0.872) (194.000,-1.064) (194.250,-0.172) (194.500,-0.913) (194.750,-1.338) (195.000,-1.186) (195.250,-0.986) (195.500,0.866) (195.750,0.975) (196.000,1.606) (196.250,1.077) (196.500,0.458) (196.750,-0.172) (197.000,-0.625) (197.250,-0.149) (197.500,0.710) (197.750,0.071) (198.000,0.944) (198.250,0.698) (198.500,0.554) (198.750,0.516) (199.000,0.574) (199.250,0.600) (199.500,-1.106) (199.750,1.208) (200.000,0.978)};
\end{axis}
\end{tikzpicture}
\caption{双势阱 $V(y)=(y^2-1)^2$ 的 Langevin 方程. 左: 势能与 Gibbs 平稳密度 $\pi\propto e^{-V}$, 密度集中在两个势阱 $y=\pm1$ 附近. 右: 一条轨道, 它在两个势阱之间随机地来回跳跃}
\label{fig:langevin}
\end{figure}

% ---------------------------------------------------------------------
\subsection{一维扩散的出口问题}

第~4~章用鞅方法计算了布朗运动的出口概率与出口时间 (定理~\ref{thm:exit}). 对一般的一维扩散, 关键是找到使 $\phi(X_t)$ 成为鞅的函数 $\phi$, 即 $L\phi=0$ 的解.

设 $\ell<r$, 且 $\sigma>0$ 在 $[\ell,r]$ 上成立. 定义\textbf{尺度函数}
\begin{equation}\label{eq:scale}
s(x)=\int_{c}^x\exp\Big(-\int_{c}^y\frac{2b(z)}{\sigma(z)^2}\dd z\Big)\dd y,
\end{equation}
其中 $c\in[\ell,r]$ 是任取的参考点. 则 $s'>0$, 且 $s''=-\frac{2b}{\sigma^2}s'$, 即
\[
Ls=bs'+\frac{1}{2}\sigma^2s''=0.
\]
所以 $s(X_t)$ (在离开 $[\ell,r]$ 之前) 没有漂移: 尺度函数把扩散过程变成了 ``时间变换后的布朗运动'', 它的作用相当于布朗运动的恒等函数. 令 $\tau=\inf\{t: X_t\notin(\ell,r)\}$.

\begin{lemma}\label{lem:exit-finite}
对 $x\in(\ell,r)$, $\E_x\tau<\infty$.
\end{lemma}

\begin{proof}
记 $\sigma_0^2=\min_{[\ell,r]}\sigma^2>0$, $\beta=\max_{[\ell,r]}\abs{b}$, 取 $\lambda=\frac{2(\beta+1)}{\sigma_0^2}$. 令 $f(y)=e^{\lambda y}$ (在 $[\ell-1,r+1]$ 之外修改为 $C^2_c$ 函数, 不影响下面的论证). 在 $[\ell,r]$ 上,
\[
Lf=\Big(\lambda b+\frac{1}{2}\lambda^2\sigma^2\Big)e^{\lambda y}\ge\lambda\Big(\frac{1}{2}\lambda\sigma_0^2-\beta\Big)e^{\lambda\ell}=\lambda e^{\lambda\ell}=:c_0>0.
\]
对有界停时 $\tau\wedge n$ 应用 Dynkin 公式, $\E_xf(X_{\tau\wedge n})-f(x)=\E_x\int_0^{\tau\wedge n}Lf(X_s)\dd s\ge c_0\E_x(\tau\wedge n)$, 而左端不超过 $2e^{\lambda r}$. 令 $n\to\infty$, 由单调收敛定理 $\E_x\tau\le2e^{\lambda r}/c_0$.
\end{proof}

\begin{theorem}[一维扩散的出口问题]\label{thm:exit-diffusion}
设 $x\in(\ell,r)$. 则
\begin{enumerate}
\item $\P_x(X_\tau=r)=\dfrac{s(x)-s(\ell)}{s(r)-s(\ell)}$;
\item $\E_x\tau=v(x)$, 其中 $v$ 是边值问题 $Lv=-1$ ($\ell<x<r$), $v(\ell)=v(r)=0$ 的解.
\end{enumerate}
\end{theorem}

\begin{proof}
(1) 把 $s$ 在 $[\ell,r]$ 的一个邻域之外修改为 $C^2_c$ 函数. 由引理~\ref{lem:exit-finite} 与 Dynkin 公式 ($Ls=0$ 于 $[\ell,r]$), $\E_xs(X_\tau)=s(x)$. 而 $X_\tau\in\{\ell,r\}$, 记 $p=\P_x(X_\tau=r)$, 则 $p\,s(r)+(1-p)s(\ell)=s(x)$.

(2) 边值问题是二阶线性常微分方程的两点边值问题 (首项系数 $\frac{1}{2}\sigma^2>0$), 有唯一的 $C^2$ 解 (可用 $s$ 显式写出). 对 $v$ 应用 Dynkin 公式, $\E_xv(X_\tau)-v(x)=\E_x\int_0^\tau(-1)\dd s$, 而 $v(X_\tau)=0$, 所以 $\E_x\tau=v(x)$.
\end{proof}

\begin{example}\label{ex:exit-examples}
(1) 布朗运动: $s(x)=x$, $Lv=\frac{1}{2}v''=-1$ 给出 $v(x)=(x-\ell)(r-x)$, 与定理~\ref{thm:exit} 一致.

(2) 带漂移 $\mu\neq0$ 的布朗运动: $s(x)=\frac{1-e^{-2\mu x}}{2\mu}$ (取 $c=0$). 于是
\[
\P_x(X_\tau=r)=\frac{e^{-2\mu\ell}-e^{-2\mu x}}{e^{-2\mu\ell}-e^{-2\mu r}}.
\]
设 $\mu<0$, $x=0$, $r=a$, 令 $\ell\to-\infty$, 则 $e^{-2\mu\ell}\to0$, 右端趋于 $e^{2\mu a}$: 这就是定理~\ref{thm:drift-hitting} 中 $\P(\tau^\mu_a<\infty)=e^{-2\abs{\mu}a}$.

(3) 几何布朗运动 $\dd S=\mu S\dd t+\sigma S\dd B$ (在 $(0,\infty)$ 上): $\frac{2b}{\sigma^2}=\frac{2\mu}{\sigma^2y}$, $s'(y)=y^{-2\mu/\sigma^2}$ (取 $c=1$). 若 $\mu<\frac{1}{2}\sigma^2$, 则 $1-\frac{2\mu}{\sigma^2}>0$, $s(y)=\frac{y^{1-2\mu/\sigma^2}-1}{1-2\mu/\sigma^2}$, $s(0+)$ 有限. 令 $\ell\downarrow0$, 从 $x$ 出发在到达 $r$ 之前 ``跌到 $0$ 附近'' 的概率趋于 $1-\big(\frac{x}{r}\big)^{1-2\mu/\sigma^2}$. 当 $r\to\infty$ 时它趋于 $1$, 这与命题~\ref{prop:gbm} (4) 中 $\mu<\frac{1}{2}\sigma^2$ 时 $S_t\to0$ 一致.
\end{example}

\begin{remark}[边界分类与常返性]
尺度函数也决定了一维扩散的长期行为. 若 $\sigma>0$ 处处成立, 则扩散 $X$ (在 $\R$ 上) 常返 (即以概率 $1$ 无穷多次到达每一点) 当且仅当 $s(-\infty)=-\infty$ 且 $s(+\infty)=+\infty$; 若 $s(+\infty)<\infty$, 则以正概率 $X_t\to+\infty$. 更精细的 Feller 边界分类还需要 \textbf{速度测度} $m(\dd x)=\frac{2\dd x}{\sigma(x)^2s'(x)}$, 它满足 $\E_x\tau=\int_\ell^rG(x,y)\,m(\dd y)$, 其中 $G$ 是 Green 函数 (见~\cite{karatzas} 5.5 节). 例如 OU 过程的尺度函数 $s'(y)=e^{\kappa y^2/\sigma^2}$ (取 $\theta=0$) 在 $\pm\infty$ 处的积分发散, 所以 OU 过程常返.
\end{remark}

% =====================================================================
\section{随机微分方程与偏微分方程}
\label{sec:6-4}

向后方程 (定理~\ref{thm:backward}) 说明: 抛物型偏微分方程 $\partial_tu=Lu$ 的解可以表示为扩散过程的期望. 本节把这种 ``概率表示'' 推广到椭圆型方程的边值问题 (Dirichlet 问题) 与带位势和源项的抛物型方程 (Feynman-Kac 公式), 并以金融数学中的期权定价作为综合应用.

% ---------------------------------------------------------------------
\subsection{Dirichlet 问题}

设 $D\subseteq\R^d$ 是有界开集, $g$ 是 $\partial D$ 上的连续函数. \textbf{Dirichlet 问题}是求 $u\in C^2(D)\cap C(\overline{D})$, 使得
\[
\Delta u=0\quad\text{于 }D,\qquad u=g\quad\text{于 }\partial D.
\]
设 $\rvec{W}$ 是从 $\rvec{x}\in D$ 出发的 $d$ 维布朗运动 (概率与期望记为 $\P_{\rvec{x}}$, $\E_{\rvec{x}}$), $\tau_D=\inf\{t:\rvec{W}_t\notin D\}$.

\begin{theorem}[Dirichlet 问题的概率表示]\label{thm:dirichlet}
设 $k$ 是 $\overline{D}$ 上的有界连续函数. 若 $u\in C^2(D)\cap C(\overline{D})$ 满足
\[
\frac{1}{2}\Delta u=-k\quad\text{于 }D,\qquad u=g\quad\text{于 }\partial D,
\]
则 $u(\rvec{x})=\E_{\rvec{x}}\big[g(\rvec{W}_{\tau_D})+\int_0^{\tau_D}k(\rvec{W}_s)\dd s\big]$. 特别地 ($k=0$), 调和函数 $u$ 满足 $u(\rvec{x})=\E_{\rvec{x}}g(\rvec{W}_{\tau_D})$.
\end{theorem}

\begin{proof}
设 $D$ 包含于半径为 $R$ 的球中. 与定理~\ref{thm:exit} (1) 相同 (考察第一个分量离开 $(-R,R)$ 的时刻), $\E_{\rvec{x}}\tau_D<\infty$. 令 $D_n=\{\rvec{y}\in D:\operatorname{dist}(\rvec{y},\partial D)>\frac{1}{n}\}$, $\tau_n=\inf\{t:\rvec{W}_t\notin D_n\}$. $u$ 在 $\overline{D_n}$ 的一个邻域上是 $C^2$ 的 (把它在此邻域之外修改为 $\R^d$ 上的 $C^2$ 函数, 不影响下面的论证). 由多维 It\^o 公式~\ref{thm:ito-multi},
\[
u(\rvec{W}_{t\wedge\tau_n})=u(\rvec{x})+\int_0^{t\wedge\tau_n}\frac{1}{2}\Delta u(\rvec{W}_s)\dd s+\sum_{i=1}^d\int_0^t\1_{\{s\le\tau_n\}}\partial_iu(\rvec{W}_s)\dd W^i_s.
\]
被积过程 $\1_{\{s\le\tau_n\}}\partial_iu(\rvec{W}_s)$ 有界, 随机积分的期望为零, 所以
\[
u(\rvec{x})=\E_{\rvec{x}}\Big[u(\rvec{W}_{t\wedge\tau_n})+\int_0^{t\wedge\tau_n}k(\rvec{W}_s)\dd s\Big].
\]
先令 $t\to\infty$, 再令 $n\to\infty$. 由轨道连续性, $\tau_n\uparrow\tau_D$, $\rvec{W}_{\tau_n}\to\rvec{W}_{\tau_D}\in\partial D$. 由于 $u$ 在 $\overline{D}$ 上连续有界, $k$ 有界且 $\E_{\rvec{x}}\tau_D<\infty$, 由控制收敛定理即得结论.
\end{proof}

\begin{corollary}\label{cor:dirichlet}
\begin{enumerate}
\item (唯一性与最大值原理) Dirichlet 问题至多有一个解, 且解满足 $\min_{\partial D}g\le u\le\max_{\partial D}g$.
\item (平均值性质) 若 $u$ 在包含闭球 $\overline{B(\rvec{x},\rho)}$ 的开集上调和, 则 $u(\rvec{x})$ 等于 $u$ 在球面 $\partial B(\rvec{x},\rho)$ 上 (关于面积测度) 的平均值.
\item (球内的期望出口时间) 对 $D=B(\rvec{0},R)$, $\E_{\rvec{x}}\tau_D=\frac{R^2-\abs{\rvec{x}}^2}{d}$.
\end{enumerate}
\end{corollary}

\begin{proof}
(1) 由表示公式立即得到. (2) 由布朗运动的旋转不变性 (命题~\ref{prop:bm-rotation}), $\rvec{W}_{\tau}$ ($\tau$ 为离开球的时刻) 的分布在绕 $\rvec{x}$ 的旋转下不变, 所以是球面上的均匀分布. (3) 取 $u(\rvec{y})=\frac{R^2-\abs{\rvec{y}}^2}{d}$, 则 $\frac{1}{2}\Delta u=\frac{1}{2}\cdot\frac{-2d}{d}=-1$, 在球面上 $u=0$; 对 $k\equiv1$, $g\equiv0$ 应用定理.
\end{proof}

定理~\ref{thm:dirichlet} 假设了解的存在性. 反过来, $u(\rvec{x}):=\E_{\rvec{x}}g(\rvec{W}_{\tau_D})$ 总是 $D$ 内的调和函数 (由强马尔可夫性与平均值性质), 但它在边界上不一定连续. 若边界满足适当的正则性 (例如每个边界点处都有位于 $D$ 外部的小圆锥), 则 $u$ 在边界上连续, 从而是 Dirichlet 问题的解 (见~\cite{karatzas} 4.2 节). 正则性不能去掉: 设 $D=\{\rvec{y}\in\R^2:0<\abs{\rvec{y}}<1\}$ 是去心圆盘, 在单位圆上 $g=0$, 在原点 $g(\rvec{0})=1$. 由例~\ref{ex:bm-annulus}, 平面布朗运动不击中原点, 所以 $\E_{\rvec{x}}g(\rvec{W}_{\tau_D})=0$; 而 $u\equiv0$ 在原点处不等于 $g(\rvec{0})=1$, 由唯一性, 这个 Dirichlet 问题无解. 原点是一个 ``非正则'' 边界点: 布朗运动从它附近出发, 几乎必然不会从它离开.

% ---------------------------------------------------------------------
\subsection{Feynman-Kac 公式}

\begin{theorem}[Feynman-Kac 公式]\label{thm:feynman-kac}
设 $b,\sigma$ 满足 \eqref{eq:H1}\eqref{eq:H2}, $L$ 如 \eqref{eq:L-operator}, $c\ge0$ 与 $g$ 是 $\R$ 上的连续函数, $f$ 连续. 设 $v\in C^{1,2}([0,T)\times\R)\cap C([0,T]\times\R)$ 满足
\begin{equation}\label{eq:feynman-kac-pde}
\partial_tv+Lv-cv+g=0\quad\text{于 }[0,T)\times\R,\qquad v(T,x)=f(x),
\end{equation}
并设存在常数 $C$, 使 $\abs{v(t,x)}+\abs{g(x)}\le C(1+x^2)$. 则对 $(t,x)\in[0,T]\times\R$,
\[
v(t,x)=\E\Big[f(X^{t,x}_T)e^{-\int_t^Tc(X^{t,x}_u)\dd u}+\int_t^Tg(X^{t,x}_s)e^{-\int_t^sc(X^{t,x}_u)\dd u}\dd s\Big].
\]
\end{theorem}

\begin{proof}
记 $X=X^{t,x}$, $\Gamma_s=e^{-\int_t^sc(X_u)\dd u}\in(0,1]$, 并令
\[
Y_s=v(s,X_s)\Gamma_s+\int_t^sg(X_u)\Gamma_u\dd u,\qquad s\in[t,T].
\]
由 It\^o 公式与乘积法则 ($\dd\Gamma_s=-c(X_s)\Gamma_s\dd s$ 是有限变差的),
\[
\dd Y_s=\Gamma_s\big(\partial_tv+Lv-cv+g\big)(s,X_s)\dd s+\Gamma_s\,\partial_xv(s,X_s)\sigma(X_s)\dd B_s=\Gamma_s\,\partial_xv(s,X_s)\sigma(X_s)\dd B_s.
\]
令 $\tau_n=\inf\{s\ge t:\abs{X_s}\ge n\}\wedge(T-\frac{1}{n})$. 在 $[t,\tau_n]$ 上被积过程有界, 所以停止后的随机积分是鞅, $v(t,x)=Y_t=\E Y_{\tau_n}$. 令 $n\to\infty$: 由 $v$ 在 $[0,T]\times\R$ 上的连续性, $Y_{\tau_n}\to Y_T$; 又由增长条件,
\[
\abs{Y_{\tau_n}}\le C\Big(1+\sup_{s\le T}X_s^2\Big)+TC\Big(1+\sup_{s\le T}X_s^2\Big),
\]
由矩估计 \eqref{eq:sde-moment} 右端可积. 由控制收敛定理, $v(t,x)=\E Y_T$, 这就是所述公式 ($v(T,X_T)=f(X_T)$).
\end{proof}

\begin{intuition}
把 $c(x)$ 理解为粒子在位置 $x$ 处被 ``杀死'' (吸收) 的速率, $g(x)$ 理解为单位时间的 ``收益'', $f$ 理解为终端 ``奖励''. 则 $v(t,x)$ 是从 $(t,x)$ 出发、在被杀死之前累积的期望总收益. 在金融中, $c=r$ 是贴现率, $f$ 是期权的收益, $v$ 就是期权的价格. Feynman 最初用路径积分 (形式上的 ``Wiener 测度上的积分'') 表示 Schr\"odinger 方程的解; Kac 意识到对热方程这可以严格化, 并用它计算了布朗运动的占位时分布 (例如由 Feynman-Kac 公式可以推出 4.5.4 小节提到的第二个反正弦律).
\end{intuition}

\begin{remark}
(1) 定理同样假设了解的存在性. 若 $b,\sigma,c,g,f$ 足够光滑, 则右端的期望确实是 $C^{1,2}$ 的, 从而是 \eqref{eq:feynman-kac-pde} 的解 (见~\cite{karatzas} 定理 5.7.6). (2) 增长条件可以放宽为多项式增长, 此时需要解的高阶矩估计 $\E\sup_t\abs{X_t}^{2m}<\infty$, 它可以由 BDG 不等式~\ref{thm:bdg} 与 Gronwall 不等式得到.
\end{remark}

% ---------------------------------------------------------------------
\subsection{Black-Scholes 理论}

考虑一个由两种资产构成的市场: 无风险债券 $\beta_t=e^{rt}$ ($r\ge0$ 为利率) 与股票, 后者的价格满足
\[
\dd S_t=\mu S_t\dd t+\sigma S_t\dd B_t,\qquad S_0>0,
\]
$\{\mathcal{F}_t\}$ 是 $B$ 的扩充信息流. 一个\textbf{交易策略}是一对循序可测过程 $(\varphi_t,\psi_t)$, 分别表示持有的股票与债券的份数, 其价值为 $V_t=\varphi_tS_t+\psi_t\beta_t$. 若
\begin{equation}\label{eq:self-financing}
\dd V_t=\varphi_t\dd S_t+\psi_t\dd\beta_t,
\end{equation}
即价值的变化完全来自资产价格的变化 (没有外部资金的注入或抽出), 称策略是\textbf{自融资}的. 记贴现价格 $\widetilde{S}_t=e^{-rt}S_t$, $\widetilde{V}_t=e^{-rt}V_t$. 由乘积法则, 自融资条件等价于
\begin{equation}\label{eq:self-financing-discounted}
\dd\widetilde{V}_t=e^{-rt}\big(\dd V_t-rV_t\dd t\big)=e^{-rt}\varphi_t\big(\dd S_t-rS_t\dd t\big)=\varphi_t\dd\widetilde{S}_t.
\end{equation}

\emph{风险中性测度.} 令 $\theta=-\frac{\mu-r}{\sigma}$ (常数), $\widetilde{\P}=e^{\theta B_T-\frac{1}{2}\theta^2T}\cdot\P$. 由 Girsanov 定理~\ref{thm:girsanov}, $\widetilde{B}_t=B_t+\frac{\mu-r}{\sigma}t$ 是 $\widetilde{\P}$ 下的布朗运动, 且
\[
\dd S_t=rS_t\dd t+\sigma S_t\dd\widetilde{B}_t,\qquad\dd\widetilde{S}_t=\sigma\widetilde{S}_t\dd\widetilde{B}_t.
\]
所以在 $\widetilde{\P}$ 下, 股票的平均收益率等于无风险利率, 贴现价格 $\widetilde{S}$ 是鞅 (它是指数鞅). 由 \eqref{eq:self-financing-discounted}, 若 $\widetilde{\E}\int_0^T\varphi_t^2\widetilde{S}_t^2\dd t<\infty$ (称这样的策略为\textbf{容许}的), 则贴现价值 $\widetilde{V}$ 是 $\widetilde{\P}$-鞅. 这推出\textbf{无套利}: 不存在容许的自融资策略使 $V_0=0$, $V_T\ge0$ 且 $\P(V_T>0)>0$, 因为 $\widetilde{\E}\widetilde{V}_T=\widetilde{V}_0=0$ 迫使 $V_T=0$ $\widetilde{\P}$-$\as$, 从而 $\P$-$\as$ ($\widetilde{\P}$ 与 $\P$ 等价).

\begin{theorem}[Black-Scholes]\label{thm:black-scholes}
设欧式衍生品在时刻 $T$ 的收益为 $F=h(S_T)$, $h\ge0$ 连续且 $h(s)\le C(1+s)$.
\begin{enumerate}
\item (复制) 存在容许的自融资策略, 其价值满足 $V_T=F$. 任何这样的策略在时刻 $t$ 的价值都是
\begin{align*}
V_t&=e^{-r(T-t)}\widetilde{\E}\big[F\,\big|\,\mathcal{F}_t\big]=v(t,S_t),\\
v(t,s)&=e^{-r(T-t)}\E\,h\Big(s\,e^{(r-\frac{1}{2}\sigma^2)(T-t)+\sigma\sqrt{T-t}\,Z}\Big),
\end{align*}
其中 $Z\sim N(0,1)$. $V_0=v(0,S_0)$ 称为衍生品的 (无套利) 价格. 它与 $\mu$ 无关.
\item (看涨期权) 对 $h(s)=(s-K)^+$,
\[
v(t,s)=s\,\Phi(d_+)-Ke^{-r(T-t)}\Phi(d_-),\quad d_\pm=\frac{\log(s/K)+(r\pm\frac{1}{2}\sigma^2)(T-t)}{\sigma\sqrt{T-t}}.
\]
\item (偏微分方程与对冲) 若 $v\in C^{1,2}([0,T)\times(0,\infty))$ (对看涨期权由 (2) 显然), 则 $v$ 满足 \textbf{Black-Scholes 方程}
\[
\partial_tv+rs\,\partial_sv+\frac{1}{2}\sigma^2s^2\,\partial_{ss}v-rv=0,\qquad v(T,s)=h(s),
\]
且复制策略的股票头寸为 $\varphi_t=\partial_sv(t,S_t)$ (\textbf{Delta 对冲}).
\end{enumerate}
\end{theorem}

\begin{proof}
(1) 由于 $S_T$ 是对数正态的, $F\in L^2(\widetilde{\P})$. $B$ 与 $\widetilde{B}$ 只相差一个确定的函数, 所以它们生成相同的扩充信息流. 对 $\widetilde{\P}$ 与 $\widetilde{\P}$-布朗运动 $\widetilde{B}$ 应用鞅表示定理 (推论~\ref{cor:mart-rep}), 鞅 $M_t=\widetilde{\E}[e^{-rT}F\,|\,\mathcal{F}_t]$ 可以写成 $M_t=M_0+\int_0^tH_s\dd\widetilde{B}_s$, $H\in\mathscr{L}^2_T$ (关于 $\widetilde{\P}$). 令
\[
\varphi_t=\frac{H_t}{\sigma\widetilde{S}_t},\qquad\psi_t=M_t-\varphi_t\widetilde{S}_t.
\]
则 $\widetilde{V}_t=e^{-rt}(\varphi_tS_t+\psi_te^{rt})=\varphi_t\widetilde{S}_t+\psi_t=M_t$, 且 $\dd\widetilde{V}_t=H_t\dd\widetilde{B}_t=\varphi_t\sigma\widetilde{S}_t\dd\widetilde{B}_t=\varphi_t\dd\widetilde{S}_t$, 即策略自融资、容许, 且 $V_T=e^{rT}M_T=F$. 若另一个容许的自融资策略也复制 $F$, 其贴现价值是 $\widetilde{\P}$-鞅, 终值为 $e^{-rT}F$, 所以它等于 $M_t$. 最后, $S_T=S_t\exp\big((r-\frac{1}{2}\sigma^2)(T-t)+\sigma(\widetilde{B}_T-\widetilde{B}_t)\big)$, 其中 $\widetilde{B}_T-\widetilde{B}_t\sim N(0,T-t)$ 与 $\mathcal{F}_t$ 独立, 而 $S_t$ 是 $\mathcal{F}_t$-可测的, 由代入引理~\ref{lem:freezing}, $V_t=e^{rt}M_t=v(t,S_t)$.

(2) 记 $\tau=T-t$, $X=s\,e^{(r-\frac{1}{2}\sigma^2)\tau+\sigma\sqrt{\tau}Z}$. $X>K$ 当且仅当 $Z>-d_-$. 由配方 $\E[e^{aZ};Z>-d]=e^{a^2/2}\P(Z>-d-a)=e^{a^2/2}\Phi(d+a)$ (取 $a=\sigma\sqrt{\tau}$, 并注意 $d_-+\sigma\sqrt{\tau}=d_+$),
\[
\E\big[(X-K)^+\big]=s\,e^{(r-\frac{1}{2}\sigma^2)\tau}e^{\frac{1}{2}\sigma^2\tau}\Phi(d_+)-K\,\Phi(d_-)=se^{r\tau}\Phi(d_+)-K\Phi(d_-).
\]
乘以 $e^{-r\tau}$ 即得.

(3) $e^{-rt}v(t,S_t)=M_t$ 是 $\widetilde{\P}$-鞅. 对 $t<T$ 应用 It\^o 公式 (在 $\widetilde{\P}$ 下, $\dd S=rS\dd t+\sigma S\dd\widetilde{B}$),
\[
\dd\big(e^{-rt}v(t,S_t)\big)=e^{-rt}G(t,S_t)\dd t+e^{-rt}\partial_sv(t,S_t)\sigma S_t\dd\widetilde{B}_t,
\]
其中 $G=\partial_tv+rs\,\partial_sv+\frac{1}{2}\sigma^2s^2\partial_{ss}v-rv$.
于是 $\int_0^te^{-ru}G(u,S_u)\dd u$ 是连续的有限变差过程, 又等于两个连续局部鞅之差, 由 It\^o 过程分解的唯一性 (命题~\ref{prop:ito-process} (1)), 它恒为零, 且 $H_t=e^{-rt}\partial_sv(t,S_t)\sigma S_t$, 即 $\varphi_t=\partial_sv(t,S_t)$. 若在某点 $(t_0,s_0)$ 处 $G\ne0$, 由连续性, $G$ 在某个邻域 $U$ 内不为零 (不妨设为正); 而 $S_u$ 对每个 $u>0$ 有处处为正的密度, 所以 $\widetilde{\E}\int_0^{T}\1_U(u,S_u)\dd u>0$; 而 $\int_0^te^{-ru}G(u,S_u)\dd u\equiv0$ 推出 $G(u,S_u)=0$ 对几乎每个 $u$ 几乎必然成立, 矛盾. 所以 $G\equiv0$.
\end{proof}

\begin{figure}[htbp]
\centering
\begin{tikzpicture}
\begin{axis}[SAaxis, axis x line=bottom, axis y line=left, width=12cm, height=5.8cm, xmin=50, xmax=162, ymin=0, ymax=65,
  xtick={60,80,100,120,140,160}, ytick={0,20,40,60}, xlabel={$s$},
  legend style={font=\small,draw=none,fill=none,at={(0.02,0.98)},anchor=north west}, legend cell align=left]
  \addplot[SADark,thick,dashed] coordinates {(50,0) (100,0) (160,60)};
  \addplot[SABlue,thick] coordinates {(50.00,0.00) (51.00,0.00) (52.00,0.00) (53.00,0.00) (54.00,0.00) (55.00,0.00) (56.00,0.00) (57.00,0.00) (58.00,0.00) (59.00,0.00) (60.00,0.00) (61.00,0.00) (62.00,0.00) (63.00,0.00) (64.00,0.01) (65.00,0.01) (66.00,0.01) (67.00,0.02) (68.00,0.03) (69.00,0.03) (70.00,0.05) (71.00,0.06) (72.00,0.08) (73.00,0.10) (74.00,0.13) (75.00,0.17) (76.00,0.21) (77.00,0.26) (78.00,0.33) (79.00,0.40) (80.00,0.48) (81.00,0.58) (82.00,0.70) (83.00,0.83) (84.00,0.98) (85.00,1.14) (86.00,1.33) (87.00,1.54) (88.00,1.77) (89.00,2.03) (90.00,2.31) (91.00,2.61) (92.00,2.95) (93.00,3.31) (94.00,3.69) (95.00,4.11) (96.00,4.55) (97.00,5.01) (98.00,5.51) (99.00,6.03) (100.00,6.58) (101.00,7.16) (102.00,7.76) (103.00,8.39) (104.00,9.04) (105.00,9.71) (106.00,10.41) (107.00,11.13) (108.00,11.87) (109.00,12.63) (110.00,13.40) (111.00,14.20) (112.00,15.01) (113.00,15.84) (114.00,16.68) (115.00,17.53) (116.00,18.40) (117.00,19.28) (118.00,20.17) (119.00,21.08) (120.00,21.99) (121.00,22.91) (122.00,23.83) (123.00,24.77) (124.00,25.71) (125.00,26.66) (126.00,27.61) (127.00,28.57) (128.00,29.53) (129.00,30.50) (130.00,31.47) (131.00,32.44) (132.00,33.41) (133.00,34.39) (134.00,35.38) (135.00,36.36) (136.00,37.34) (137.00,38.33) (138.00,39.32) (139.00,40.31) (140.00,41.30) (141.00,42.29) (142.00,43.29) (143.00,44.28) (144.00,45.28) (145.00,46.27) (146.00,47.27) (147.00,48.26) (148.00,49.26) (149.00,50.26) (150.00,51.26) (151.00,52.25) (152.00,53.25) (153.00,54.25) (154.00,55.25) (155.00,56.25) (156.00,57.25) (157.00,58.25) (158.00,59.25) (159.00,60.25) (160.00,61.25)};
  \addplot[SARed,thick] coordinates {(50.00,0.12) (51.00,0.14) (52.00,0.17) (53.00,0.20) (54.00,0.24) (55.00,0.29) (56.00,0.34) (57.00,0.39) (58.00,0.46) (59.00,0.53) (60.00,0.61) (61.00,0.69) (62.00,0.79) (63.00,0.90) (64.00,1.01) (65.00,1.14) (66.00,1.27) (67.00,1.42) (68.00,1.58) (69.00,1.76) (70.00,1.94) (71.00,2.14) (72.00,2.35) (73.00,2.58) (74.00,2.81) (75.00,3.07) (76.00,3.33) (77.00,3.62) (78.00,3.91) (79.00,4.23) (80.00,4.55) (81.00,4.90) (82.00,5.25) (83.00,5.63) (84.00,6.01) (85.00,6.42) (86.00,6.84) (87.00,7.27) (88.00,7.72) (89.00,8.18) (90.00,8.66) (91.00,9.15) (92.00,9.66) (93.00,10.19) (94.00,10.72) (95.00,11.27) (96.00,11.84) (97.00,12.42) (98.00,13.01) (99.00,13.61) (100.00,14.23) (101.00,14.86) (102.00,15.50) (103.00,16.16) (104.00,16.83) (105.00,17.51) (106.00,18.19) (107.00,18.90) (108.00,19.61) (109.00,20.33) (110.00,21.06) (111.00,21.80) (112.00,22.55) (113.00,23.32) (114.00,24.09) (115.00,24.86) (116.00,25.65) (117.00,26.45) (118.00,27.25) (119.00,28.06) (120.00,28.88) (121.00,29.71) (122.00,30.54) (123.00,31.38) (124.00,32.23) (125.00,33.08) (126.00,33.94) (127.00,34.80) (128.00,35.67) (129.00,36.55) (130.00,37.43) (131.00,38.31) (132.00,39.20) (133.00,40.10) (134.00,41.00) (135.00,41.90) (136.00,42.81) (137.00,43.72) (138.00,44.64) (139.00,45.56) (140.00,46.48) (141.00,47.41) (142.00,48.34) (143.00,49.27) (144.00,50.21) (145.00,51.14) (146.00,52.09) (147.00,53.03) (148.00,53.98) (149.00,54.92) (150.00,55.88) (151.00,56.83) (152.00,57.79) (153.00,58.74) (154.00,59.70) (155.00,60.66) (156.00,61.63) (157.00,62.59) (158.00,63.56) (159.00,64.53) (160.00,65.50)};
  \legend{收益 $(s-K)^+$, $T-t=0.25$, $T-t=1$}
\end{axis}
\end{tikzpicture}
\caption{看涨期权的 Black-Scholes 价格 $v(t,s)$ ($K=100$, $r=0.05$, $\sigma=0.3$). 离到期越远, 价格越高于内在价值 $(s-K)^+$; 当 $t\to T$ 时价格曲线收敛到收益函数. 曲线的斜率 $\partial_sv=\Phi(d_+)\in(0,1)$ 就是对冲所需持有的股票份数}
\label{fig:black-scholes}
\end{figure}

\begin{intuition}
Black-Scholes 理论最令人意外的结论, 是期权价格与股票的期望收益率 $\mu$ 无关. 原因是期权可以用股票和债券完全复制, 复制成本只取决于股价的\emph{波动} (二次变差 $\sigma^2S^2\dd t$ 是逐条轨道的性质, 与概率测度无关), 而与股价的\emph{趋势}无关. Girsanov 定理恰好只改变趋势而不改变波动, 所以可以在一个 ``假想'' 的测度 $\widetilde{\P}$ 下计算, 在那里所有资产的平均收益率都是 $r$, 期权价格就是贴现收益的期望. 鞅表示定理则保证复制总是可能的 (市场完备).
\end{intuition}

由 Feynman-Kac 公式 (取 $c=r$, $g=0$, $L$ 为 $\widetilde{\P}$ 下 $S$ 的生成元 $rs\partial_s+\frac{1}{2}\sigma^2s^2\partial_{ss}$), Black-Scholes 方程与期望表示是同一件事的两个侧面: 前者是 ``对冲'' 的观点 (局部无风险组合的收益率等于 $r$), 后者是 ``风险中性定价'' 的观点.

% ---------------------------------------------------------------------
\subsection{障碍期权}

\textbf{障碍期权}的收益依赖于股价是否触及某个水平. 例如\textbf{向上敲出看涨期权} (up-and-out call) 的收益为
\[
F=(S_T-K)^+\1_{\{\max_{t\le T}S_t<H\}},\qquad S_0<H,\ K<H:
\]
若股价在到期前曾经涨到障碍 $H$, 期权作废. 它比普通看涨期权便宜, 适合预期股价温和上涨的投资者. 由于 $F$ 依赖于整条轨道, 定价需要 $(\max S,S_T)$ 的联合分布. 在 $\widetilde{\P}$ 下,
\[
S_t=S_0e^{\sigma\widehat{W}_t},\qquad\widehat{W}_t=\widetilde{B}_t+\nu t,\qquad\nu=\frac{r-\frac{1}{2}\sigma^2}{\sigma},
\]
所以问题归结为带漂移布朗运动的最大值与终点的联合分布. 第~4~章的反射原理处理了无漂移的情形, Girsanov 定理可以把漂移加上.

\begin{proposition}[带漂移布朗运动的镜像原理]\label{prop:image}
设 $W$ 是标准布朗运动, $\widehat{W}_t=W_t+\nu t$, $\widehat{M}_T=\max_{t\le T}\widehat{W}_t$, $h>0$. 则对任意满足 $f(x)=0$ ($x\ge h$) 的有界 Borel 函数 $f$,
\[
\E\big[f(\widehat{W}_T);\widehat{M}_T<h\big]=\E f(\widehat{W}_T)-e^{2\nu h}\,\E f(\widehat{W}_T+2h).
\]
\end{proposition}

\begin{proof}
由 Girsanov 定理 (取 $\theta\equiv\nu$), 在测度 $Z_T\cdot\P$ ($Z_T=e^{\nu W_T-\frac{1}{2}\nu^2T}$) 下 $W$ 是漂移为 $\nu$ 的布朗运动, 所以对轨道的有界可测泛函 $\Psi$, $\E\Psi(\widehat{W})=\E[\Psi(W)Z_T]$. 记 $M_T=\max_{t\le T}W_t$. 由反射原理 \eqref{eq:reflection-joint}, 对 $x\le h$, $\P(M_T\ge h,W_T\le x)=\P(W_T\ge2h-x)$, 即 $\{M_T\ge h\}$ 上 $W_T$ 在 $(-\infty,h]$ 中的 (次) 密度为 $\varphi_T(2h-x)$, $\varphi_T(y)=\frac{1}{\sqrt{2\pi T}}e^{-y^2/(2T)}$. 由于 $f$ 在 $[h,\infty)$ 上为零,
\begin{align*}
\E\big[f(\widehat{W}_T);\widehat{M}_T\ge h\big]&=\E\big[f(W_T)e^{\nu W_T-\frac{1}{2}\nu^2T};M_T\ge h\big]\\
&=\int_{-\infty}^hf(x)e^{\nu x-\frac{1}{2}\nu^2T}\varphi_T(2h-x)\dd x.
\end{align*}
展开平方可以验证恒等式
\[
e^{\nu x-\frac{1}{2}\nu^2T}\varphi_T(2h-x)=e^{2\nu h}\varphi_T(x-2h-\nu T),
\]
(两端的指数都等于 $-\frac{(2h-x)^2}{2T}+\nu x-\frac{1}{2}\nu^2T$), 而 $\varphi_T(y-\nu T)$ 正是 $\widehat{W}_T\sim N(\nu T,T)$ 的密度. 所以上式等于 $e^{2\nu h}\int f(x)\varphi_T(x-2h-\nu T)\dd x=e^{2\nu h}\E f(\widehat{W}_T+2h)$ (积分可以扩展到全体 $x$, 因为 $f$ 在 $[h,\infty)$ 上为零). 用 $\E f(\widehat{W}_T)$ 减去它即得结论.
\end{proof}

\begin{theorem}[向上敲出看涨期权]\label{thm:barrier}
设看涨期权的初始股价为 $s$, 执行价为 $k$, 到期时间为 $T$, 记其 Black-Scholes 价格为 $C(s,k)$ (定理~\ref{thm:black-scholes} (2), $t=0$), 并令
\[
U(s)=C(s,K)-C(s,H)-(H-K)e^{-rT}\Phi\big(d_-(s,H)\big),
\]
其中 $d_-(s,H)=\frac{\log(s/H)+(r-\frac{1}{2}\sigma^2)T}{\sigma\sqrt{T}}$. 则向上敲出看涨期权在时刻 $0$ 的价格为
\[
V_0=U(S_0)-\Big(\frac{H}{S_0}\Big)^{\frac{2r}{\sigma^2}-1}U\Big(\frac{H^2}{S_0}\Big).
\]
\end{theorem}

\begin{proof}
与定理~\ref{thm:black-scholes} (1) 相同 ($F$ 有界, 鞅表示定理仍适用), 价格为 $V_0=e^{-rT}\widetilde{\E}F$. 令 $g(s)=(s-K)^+\1_{\{s<H\}}$, $h=\frac{1}{\sigma}\log\frac{H}{S_0}>0$, $f(x)=g(S_0e^{\sigma x})$, 则 $f(x)=0$ 当 $x\ge h$, 且 $\{\max S<H\}=\{\widehat{M}_T<h\}$. 由命题~\ref{prop:image},
\[
\widetilde{\E}F=\widetilde{\E}g(S_T)-e^{2\nu h}\,\widetilde{\E}\,g\big(S_0e^{\sigma\widehat{W}_T+2\sigma h}\big)=\widetilde{\E}g(S_T)-\Big(\frac{H}{S_0}\Big)^{2\nu/\sigma}\widetilde{\E}\,g\Big(\frac{H^2}{S_0}e^{\sigma\widehat{W}_T}\Big),
\]
因为 $e^{\sigma h}=\frac{H}{S_0}$, $S_0e^{2\sigma h}=\frac{H^2}{S_0}$. 这里 $\frac{2\nu}{\sigma}=\frac{2r}{\sigma^2}-1$, 而 $\frac{H^2}{S_0}e^{\sigma\widehat{W}_T}$ 正是从 $\frac{H^2}{S_0}$ 出发的 (风险中性) 股价在 $T$ 时刻的值. 最后, 由 $g(s)=(s-K)^+-(s-H)^+-(H-K)\1_{\{s\ge H\}}$ 与 $\widetilde{\P}(S_T\ge H\,|\,S_0=s)=\Phi(d_-(s,H))$ (定理~\ref{thm:black-scholes} 证明中的计算), 从 $s$ 出发时 $e^{-rT}\widetilde{\E}g(S_T)=U(s)$.
\end{proof}

\begin{figure}[htbp]
\centering
\begin{tikzpicture}
\begin{axis}[SAaxis, axis x line=bottom, axis y line=left, width=12cm, height=5.8cm, xmin=50, xmax=136, ymin=0, ymax=6,
  xtick={60,80,100,120}, ytick={0,1,2,3,4,5,6}, xlabel={$S_0$}, clip=false,
  legend style={font=\small,draw=none,fill=none,at={(0.02,0.98)},anchor=north west}, legend cell align=left]
  \addplot[SADark!60,thick,dashed,restrict y to domain=0:6] coordinates {(50.00,0.12) (50.50,0.13) (51.00,0.14) (51.50,0.16) (52.00,0.17) (52.50,0.19) (53.00,0.20) (53.50,0.22) (54.00,0.24) (54.50,0.26) (55.00,0.29) (55.50,0.31) (56.00,0.34) (56.50,0.36) (57.00,0.39) (57.50,0.42) (58.00,0.46) (58.50,0.49) (59.00,0.53) (59.50,0.57) (60.00,0.61) (60.50,0.65) (61.00,0.69) (61.50,0.74) (62.00,0.79) (62.50,0.84) (63.00,0.90) (63.50,0.95) (64.00,1.01) (64.50,1.07) (65.00,1.14) (65.50,1.20) (66.00,1.27) (66.50,1.35) (67.00,1.42) (67.50,1.50) (68.00,1.58) (68.50,1.67) (69.00,1.76) (69.50,1.85) (70.00,1.94) (70.50,2.04) (71.00,2.14) (71.50,2.24) (72.00,2.35) (72.50,2.46) (73.00,2.58) (73.50,2.69) (74.00,2.81) (74.50,2.94) (75.00,3.07) (75.50,3.20) (76.00,3.33) (76.50,3.47) (77.00,3.62) (77.50,3.76) (78.00,3.91) (78.50,4.07) (79.00,4.23) (79.50,4.39) (80.00,4.55) (80.50,4.72) (81.00,4.90) (81.50,5.07) (82.00,5.25) (82.50,5.44) (83.00,5.63) (83.50,5.82) (84.00,6.01) (84.50,6.21) (85.00,6.42) (85.50,6.62) (86.00,6.84) (86.50,7.05) (87.00,7.27) (87.50,7.49) (88.00,7.72) (88.50,7.95) (89.00,8.18) (89.50,8.42) (90.00,8.66) (90.50,8.91) (91.00,9.15) (91.50,9.41) (92.00,9.66) (92.50,9.92) (93.00,10.19) (93.50,10.45) (94.00,10.72) (94.50,11.00) (95.00,11.27) (95.50,11.55) (96.00,11.84) (96.50,12.13) (97.00,12.42) (97.50,12.71) (98.00,13.01) (98.50,13.31) (99.00,13.61) (99.50,13.92) (100.00,14.23) (100.50,14.54) (101.00,14.86) (101.50,15.18) (102.00,15.50) (102.50,15.83) (103.00,16.16) (103.50,16.49) (104.00,16.83) (104.50,17.16) (105.00,17.51) (105.50,17.85) (106.00,18.19) (106.50,18.54) (107.00,18.90) (107.50,19.25) (108.00,19.61) (108.50,19.97) (109.00,20.33) (109.50,20.69) (110.00,21.06) (110.50,21.43) (111.00,21.80) (111.50,22.18) (112.00,22.55) (112.50,22.93) (113.00,23.32) (113.50,23.70) (114.00,24.09) (114.50,24.47) (115.00,24.86) (115.50,25.26) (116.00,25.65) (116.50,26.05) (117.00,26.45) (117.50,26.85) (118.00,27.25) (118.50,27.66) (119.00,28.06) (119.50,28.47) (120.00,28.88) (120.50,29.29) (121.00,29.71) (121.50,30.12) (122.00,30.54) (122.50,30.96) (123.00,31.38) (123.50,31.80) (124.00,32.23) (124.50,32.65) (125.00,33.08) (125.50,33.51) (126.00,33.94) (126.50,34.37) (127.00,34.80) (127.50,35.23) (128.00,35.67) (128.50,36.11) (129.00,36.55) (129.50,36.99) (130.00,37.43)};
  \addplot[SARed,thick] coordinates {(50.00,0.07) (50.50,0.08) (51.00,0.09) (51.50,0.09) (52.00,0.10) (52.50,0.11) (53.00,0.12) (53.50,0.13) (54.00,0.14) (54.50,0.15) (55.00,0.16) (55.50,0.17) (56.00,0.18) (56.50,0.19) (57.00,0.21) (57.50,0.22) (58.00,0.23) (58.50,0.25) (59.00,0.26) (59.50,0.28) (60.00,0.30) (60.50,0.31) (61.00,0.33) (61.50,0.35) (62.00,0.37) (62.50,0.38) (63.00,0.40) (63.50,0.42) (64.00,0.44) (64.50,0.46) (65.00,0.48) (65.50,0.51) (66.00,0.53) (66.50,0.55) (67.00,0.57) (67.50,0.59) (68.00,0.62) (68.50,0.64) (69.00,0.66) (69.50,0.69) (70.00,0.71) (70.50,0.74) (71.00,0.76) (71.50,0.79) (72.00,0.81) (72.50,0.83) (73.00,0.86) (73.50,0.88) (74.00,0.91) (74.50,0.93) (75.00,0.96) (75.50,0.98) (76.00,1.01) (76.50,1.03) (77.00,1.05) (77.50,1.08) (78.00,1.10) (78.50,1.12) (79.00,1.14) (79.50,1.17) (80.00,1.19) (80.50,1.21) (81.00,1.23) (81.50,1.25) (82.00,1.27) (82.50,1.29) (83.00,1.31) (83.50,1.33) (84.00,1.35) (84.50,1.36) (85.00,1.38) (85.50,1.39) (86.00,1.41) (86.50,1.42) (87.00,1.44) (87.50,1.45) (88.00,1.46) (88.50,1.47) (89.00,1.49) (89.50,1.49) (90.00,1.50) (90.50,1.51) (91.00,1.52) (91.50,1.53) (92.00,1.53) (92.50,1.54) (93.00,1.54) (93.50,1.54) (94.00,1.55) (94.50,1.55) (95.00,1.55) (95.50,1.55) (96.00,1.55) (96.50,1.54) (97.00,1.54) (97.50,1.54) (98.00,1.53) (98.50,1.53) (99.00,1.52) (99.50,1.51) (100.00,1.50) (100.50,1.49) (101.00,1.48) (101.50,1.47) (102.00,1.46) (102.50,1.45) (103.00,1.44) (103.50,1.42) (104.00,1.41) (104.50,1.39) (105.00,1.38) (105.50,1.36) (106.00,1.34) (106.50,1.32) (107.00,1.30) (107.50,1.28) (108.00,1.26) (108.50,1.24) (109.00,1.22) (109.50,1.20) (110.00,1.18) (110.50,1.15) (111.00,1.13) (111.50,1.10) (112.00,1.08) (112.50,1.05) (113.00,1.03) (113.50,1.00) (114.00,0.98) (114.50,0.95) (115.00,0.92) (115.50,0.89) (116.00,0.86) (116.50,0.84) (117.00,0.81) (117.50,0.78) (118.00,0.75) (118.50,0.72) (119.00,0.69) (119.50,0.66) (120.00,0.63) (120.50,0.60) (121.00,0.57) (121.50,0.53) (122.00,0.50) (122.50,0.47) (123.00,0.44) (123.50,0.41) (124.00,0.38) (124.50,0.35) (125.00,0.31) (125.50,0.28) (126.00,0.25) (126.50,0.22) (127.00,0.19) (127.50,0.16) (128.00,0.12) (128.50,0.09) (129.00,0.06) (129.50,0.03) (130.00,0.00)};
  \addplot[SADark!30,forget plot] coordinates {(130,0) (130,6)};
  \node[font=\small,SADark!70,anchor=south] at (axis cs:130,6) {$H$};
  \legend{普通看涨期权, 向上敲出看涨期权}
\end{axis}
\end{tikzpicture}
\caption{向上敲出看涨期权 ($K=100$, $H=130$, $T=1$, $r=0.05$, $\sigma=0.3$) 与普通看涨期权 (虚线, 超出纵轴范围的部分未画出) 的价格. 敲出条款使 $S_0=100$ 处的期权价格从普通期权的 $14.23$ 降至 $1.503$, 而且价格关于 $S_0$ 不再单调, 股价越接近障碍, 期权越可能作废}
\label{fig:barrier}
\end{figure}

\begin{remark}
(1) 这里假设障碍是\emph{连续}监测的. 实际合约常在离散时刻 (例如每个交易日收盘) 检查障碍, 此时价格略高; 6.5 节将看到, 用 Monte Carlo 方法逐步模拟时, 离散监测带来的偏差只以 $O(\sqrt{h})$ 的速度消失. (2) 镜像原理的名字来自偏微分方程: $V$ 满足 Black-Scholes 方程, 边界条件为 $V(t,H)=0$; 公式中的第二项恰是在 ``镜像点'' $H^2/S_0$ 处放置的一个反号的源, 使边界上的值相互抵消. 这与静电学中用镜像电荷求解边值问题的方法完全相同.
\end{remark}

% ---------------------------------------------------------------------
\subsection{美式期权与最优停时}

\textbf{美式期权}可以在到期日 $T$ 之前的任何时刻行权. 持有者选择行权时刻 $\tau$, 它必须是停时 (行权决定只能依赖于已有的信息), 行权时获得 $g(S_\tau)$, 例如看跌期权 $g(s)=(K-s)^+$. 可以证明 (见~\cite{shreve} 第 8 章), 美式期权的无套利价格是一个\textbf{最优停时问题}的值:
\begin{equation}\label{eq:american}
V_0=\sup_{\tau\le T}\widetilde{\E}\big[e^{-r\tau}g(S_\tau)\big],
\end{equation}
上确界取遍取值于 $[0,T]$ 的停时. 在离散时间中, 这个问题可以用逆向归纳求解: 在每个时刻比较 ``立即行权的收益'' 与 ``继续持有的价值'', 二者的较大者就是该时刻的价值 (第~3~章可选停时定理的一个应用, 称为 Snell 包络). 连续时间的情形更为微妙. 下面两个结果分别展示了鞅方法与偏微分方程方法.

\begin{proposition}[不付红利的美式看涨期权]\label{prop:american-call}
对 $g(s)=(s-K)^+$, $\sup_{\tau\le T}\widetilde{\E}[e^{-r\tau}(S_\tau-K)^+]=\widetilde{\E}[e^{-rT}(S_T-K)^+]$. 即提前行权永远不是最优的, 美式看涨期权与欧式看涨期权价格相同.
\end{proposition}

\begin{proof}
我们证明 $Y_t=e^{-rt}(S_t-K)^+=(\widetilde{S}_t-Ke^{-rt})^+$ 是 $\widetilde{\P}$-下鞅. 设 $t\le u\le T$. 由于 $\widetilde{S}$ 是鞅, 函数 $x\mapsto(x-Ke^{-rt})^+$ 是凸的, 由条件 Jensen 不等式~\ref{thm:cjensen},
\begin{align*}
Y_t&=\big(\widetilde{\E}[\widetilde{S}_u\,|\,\mathcal{F}_t]-Ke^{-rt}\big)^+\le\widetilde{\E}\big[(\widetilde{S}_u-Ke^{-rt})^+\,\big|\,\mathcal{F}_t\big]\\
&\le\widetilde{\E}\big[(\widetilde{S}_u-Ke^{-ru})^+\,\big|\,\mathcal{F}_t\big]=\widetilde{\E}[Y_u\,|\,\mathcal{F}_t],
\end{align*}
其中第二个不等号用到 $Ke^{-ru}\le Ke^{-rt}$ ($r\ge0$). 由可选停时定理 (对下鞅, 有界停时 $\tau\le T$; 与推论~\ref{cor:ct-ost-bounded} 的证明相同, 离散化后由定理~\ref{thm:ost-bounded} 得到), $\widetilde{\E}Y_\tau\le\widetilde{\E}Y_T$.
\end{proof}

直观上, 提前行权意味着提前支付执行价 $K$, 损失了 $K$ 的利息, 同时放弃了股价下跌时不行权的权利. 看跌期权的情形则不同, 股价很低时立即行权得到 $K-S_t$ 并获取利息可能更有利. 有限期限的美式看跌期权没有显式解, 但永久 ($T=\infty$) 的美式看跌期权可以完全求解.

\begin{theorem}[永久美式看跌期权]\label{thm:perpetual-put}
设 $r>0$, $g(s)=(K-s)^+$. 令 $\gamma=\frac{2r}{\sigma^2}$, $s^*=\frac{\gamma K}{1+\gamma}$,
\[
V^*(s)=\begin{cases}K-s,&0<s\le s^*,\\(K-s^*)\big(s/s^*\big)^{-\gamma},&s>s^*.\end{cases}
\]
则对一切 $s>0$,
\[
V^*(s)=\sup_\tau\widetilde{\E}_s\big[e^{-r\tau}(K-S_\tau)^+\1_{\{\tau<\infty\}}\big],
\]
上确界取遍一切停时, 且上确界在 $\tau^*=\inf\{t:S_t\le s^*\}$ 处达到. ($\widetilde{\E}_s$ 表示 $S_0=s$.)
\end{theorem}

\begin{proof}
\emph{第一步: 阈值策略.} 对 $0<b\le s$, 令 $\tau_b=\inf\{t:S_t\le b\}$. 由于 $\log\frac{S_t}{s}=\sigma\widehat{W}_t$, $\tau_b$ 是漂移为 $-\nu$ 的布朗运动 $-\widehat{W}$ 首次到达 $a=\frac{1}{\sigma}\log\frac{s}{b}$ 的时刻. 由定理~\ref{thm:drift-hitting} (取 $\mu=-\nu$, $\lambda=r$) 与 $\nu^2+2r=\frac{(r+\frac{1}{2}\sigma^2)^2}{\sigma^2}$,
\[
\widetilde{\E}_s\big[e^{-r\tau_b}\1_{\{\tau_b<\infty\}}\big]=\exp\Big(-a\Big(\nu+\frac{r+\frac{1}{2}\sigma^2}{\sigma}\Big)\Big)=\exp\Big(-\frac{2r}{\sigma}a\Big)=\Big(\frac{s}{b}\Big)^{-\gamma}.
\]
所以阈值策略 $\tau_b$ ($b<K$) 的价值为 $(K-b)(s/b)^{-\gamma}$. 对 $b$ 求最大值: $\frac{\dd}{\dd b}\big[(K-b)b^\gamma\big]=b^{\gamma-1}\big(\gamma(K-b)-b\big)$, 在 $b=s^*$ 处取得最大值. 于是对 $s>s^*$, 策略 $\tau^*$ 的价值为 $V^*(s)$; 对 $s\le s^*$, 立即行权 ($\tau^*=0$) 的价值为 $K-s=V^*(s)$. 所以上确界 $\ge V^*(s)$.

\emph{第二步: 验证 $V^*$ 是上界.} (a) $V^*$ 在 $(0,\infty)$ 上连续可微: 在 $s^*$ 处两段的值都是 $K-s^*$, 右导数为 $-\gamma\frac{K-s^*}{s^*}=-1$, 等于左导数 (\textbf{光滑粘贴}条件, 这正是 $s^*$ 的最优性条件). $V^*$ 在 $s^*$ 之外二阶连续可微, 且 $V^{*\prime\prime}$ 在 $s^*$ 附近有界. (b) $V^*\ge g$: 在 $(0,s^*]$ 上相等; 在 $(s^*,\infty)$ 上 $V^*>0$, 且 $V^*$ 是凸函数, 其在 $s^*$ 处的切线就是 $K-s$, 所以 $V^*\ge K-s$. (c) 记 $\widetilde{L}=rs\partial_s+\frac{1}{2}\sigma^2s^2\partial_{ss}$ 为 $\widetilde{\P}$ 下 $S$ 的生成元. 对 $s>s^*$, 由于 $\gamma$ 满足 $\frac{1}{2}\sigma^2\gamma(\gamma+1)-r\gamma-r=0$, $\widetilde{L}V^*-rV^*=0$; 对 $s<s^*$, $\widetilde{L}V^*-rV^*=-rs-r(K-s)=-rK<0$. 所以 $\widetilde{L}V^*-rV^*\le0$ (除 $s^*$ 一点外).

It\^o 公式对 $C^1$ 且除有限个点外 $C^2$、二阶导数局部有界的函数仍然成立 (例外点处二阶导数任意取值), 因为 $S$ 在任何一点停留的总时间为零; 证明与 Tanaka 公式相同, 用光滑函数逼近 (见~\cite{karatzas} 3.6 节). 于是
\[
e^{-rt}V^*(S_t)=V^*(s)+\int_0^te^{-ru}\big(\widetilde{L}V^*-rV^*\big)(S_u)\dd u+\int_0^te^{-ru}V^{*\prime}(S_u)\sigma S_u\dd\widetilde{B}_u.
\]
由于 $\abs{V^{*\prime}}\le1$ 且 $\widetilde{\E}\int_0^tS_u^2\dd u<\infty$, 随机积分是鞅. 对任意停时 $\tau$ 与 $t>0$, 由停止性质与 (c),
\begin{align*}
\widetilde{\E}_s\big[e^{-r(\tau\wedge t)}g(S_{\tau\wedge t})\big]&\le\widetilde{\E}_s\big[e^{-r(\tau\wedge t)}V^*(S_{\tau\wedge t})\big]\\
&=V^*(s)+\widetilde{\E}_s\int_0^{\tau\wedge t}e^{-ru}\big(\widetilde{L}V^*-rV^*\big)(S_u)\dd u\le V^*(s).
\end{align*}
令 $t\to\infty$, 由 Fatou 引理 ($g\ge0$), $\widetilde{\E}_s[e^{-r\tau}g(S_\tau)\1_{\{\tau<\infty\}}]\le V^*(s)$.
\end{proof}

\begin{figure}[htbp]
\centering
\begin{tikzpicture}
\begin{axis}[SAaxis, axis x line=bottom, axis y line=left, width=11cm, height=5.6cm, xmin=0, xmax=205, ymin=0, ymax=105,
  xtick={0,52.63,100,150,200}, xticklabels={$0$,$s^*$,$K$,$150$,$200$}, ytick={0,25,50,75,100}, xlabel={$s$},
  legend style={font=\small,draw=none,fill=none,at={(0.98,0.98)},anchor=north east}, legend cell align=left]
  \addplot[SADark,thick,dashed] coordinates {(0,100) (100,0) (200,0)};
  \addplot[SARed,thick,domain=52.63:200,samples=100,forget plot] {47.368*(x/52.63)^(-1.1111)};
  \addplot[SARed,thick] coordinates {(0,100) (52.63,47.37)};
  \fill[SARed] (axis cs:52.63,47.37) circle (2pt);
  \addplot[SADark!40,densely dotted,forget plot] coordinates {(52.63,0) (52.63,47.37)};
  \legend{行权收益 $(K-s)^+$, 期权价值 $V^*(s)$}
\end{axis}
\end{tikzpicture}
\caption{永久美式看跌期权 ($K=100$, $r=0.05$, $\sigma=0.3$, $s^*\approx52.6$). 在行权区域 $s\le s^*$ 价值等于收益; 在继续持有区域 $s>s^*$ 价值严格大于收益. 两段在 $s^*$ 处光滑粘贴 (函数值与导数都连续)}
\label{fig:perpetual-put}
\end{figure}

\begin{remark}[有限期限的美式看跌期权]
对 $T<\infty$, 价值函数 $v(t,s)$ 与一条随时间变化的\textbf{最优行权边界} $s^*(t)$ (当 $t\uparrow T$ 时上升到 $K$) 由下面的\textbf{自由边界问题}决定: 在继续持有区域 $\{s>s^*(t)\}$ 中 $v$ 满足 Black-Scholes 方程, 在行权区域中 $v=K-s$, 在边界上 $v$ 与 $\partial_sv$ 连续 (光滑粘贴). 等价地, $v$ 满足变分不等式
\[
\max\Big(\partial_tv+rs\partial_sv+\frac{1}{2}\sigma^2s^2\partial_{ss}v-rv,\ (K-s)^+-v\Big)=0.
\]
这个问题没有显式解, 实际中用二叉树、有限差分或回归型 Monte Carlo 方法 (Longstaff-Schwartz) 数值求解 (见~\cite{shreve} 第 8 章, \cite{glasserman} 第 8 章). 定理~\ref{thm:perpetual-put} 的证明结构, 即 ``先猜测阈值策略并优化阈值, 再验证上鞅性质'', 是求解最优停时问题的一般方法.
\end{remark}

% =====================================================================
\section{数值方法}
\label{sec:6-5}

大多数随机微分方程没有显式解 (有限期限的美式期权就是一个例子), 需要数值求解. 数值方法有两类目标: 逼近解的\emph{轨道} (强逼近, 例如模拟一个具体的系统、做滤波), 或逼近解的\emph{分布}的泛函 $\E f(X_T)$ (弱逼近, 例如 Monte Carlo 定价). 本节介绍最基本的 Euler-Maruyama 格式, 证明它的强收敛阶, 并讨论 Monte Carlo 方法.

% ---------------------------------------------------------------------
\subsection{Euler-Maruyama 格式}

设 $b,\sigma$ 不依赖于时间, 满足 \eqref{eq:H1}\eqref{eq:H2}. 取步长 $h=T/N$, 格点 $t_k=kh$. 模仿常微分方程的 Euler 方法, 在每一步把系数 ``冻结'' 在左端点:
\begin{equation}\label{eq:euler-maruyama}
\widehat{X}_{t_{k+1}}=\widehat{X}_{t_k}+b(\widehat{X}_{t_k})\,h+\sigma(\widehat{X}_{t_k})\,\Delta B_k,\qquad\Delta B_k=B_{t_{k+1}}-B_{t_k},\qquad\widehat{X}_0=\xi.
\end{equation}
模拟时, $\Delta B_k$ 是独立的 $N(0,h)$ 随机数. 注意系数必须取在左端点: 这正是 It\^o 积分的约定 (5.1 节), 取在中点的格式会收敛到 Stratonovich 方程的解.

\begin{theorem}[Euler-Maruyama 格式的强收敛]\label{thm:em}
设 $\xi\in L^2(\mathcal{F}_0)$, $X$ 是方程的解, $\widehat{X}$ 由 \eqref{eq:euler-maruyama} 定义. 则存在只依赖于 $K,T,\E\xi^2$ 的常数 $C$, 使得对一切 $h\le1$,
\[
\E\Big[\max_{0\le k\le N}\big\lvert X_{t_k}-\widehat{X}_{t_k}\big\rvert^2\Big]\le C\,h.
\]
\end{theorem}

\begin{proof}
记 $\eta(s)=t_k$ ($s\in[t_k,t_{k+1})$), 并定义连续时间的插值
\[
\widehat{X}_t=\xi+\int_0^tb(\widehat{X}_{\eta(s)})\dd s+\int_0^t\sigma(\widehat{X}_{\eta(s)})\dd B_s,\qquad t\in[0,T].
\]
在 $[t_k,t_{k+1}]$ 上被积过程是常数, 所以在格点上它与 \eqref{eq:euler-maruyama} 一致, 且在格点之间 $\widehat{X}_t-\widehat{X}_{\eta(t)}=b(\widehat{X}_{\eta(t)})(t-\eta(t))+\sigma(\widehat{X}_{\eta(t)})(B_t-B_{\eta(t)})$.

\emph{第一步: 矩估计.} 由 \eqref{eq:euler-maruyama} 与 \eqref{eq:H2} 对 $k$ 归纳可知 $\E\sup_t\abs{\widehat{X}_t}^2<\infty$. 记 $\widehat{m}(t)=\E\sup_{u\le t}\abs{\widehat{X}_u}^2$. 与定理~\ref{thm:sde-existence} 中矩估计的推导相同, $\widehat{m}(t)\le3\E\xi^2+6K^2(T+4)\int_0^t(1+\widehat{m}(s))\dd s$, 由 Gronwall 不等式, $\widehat{m}(T)\le M$, 其中 $M$ 与定理~\ref{thm:sde-existence} 中的上界相同, 不依赖于 $h$.

\emph{第二步: 局部误差.} 由于 $B_s-B_{\eta(s)}$ 与 $\mathcal{F}_{\eta(s)}$ 独立且方差不超过 $h$, 由 \eqref{eq:H2},
\[
\E\abs{\widehat{X}_s-\widehat{X}_{\eta(s)}}^2\le2h^2\E b(\widehat{X}_{\eta(s)})^2+2h\,\E\sigma(\widehat{X}_{\eta(s)})^2\le2(h^2+h)\cdot2K^2(1+M)\le C_1h.
\]

\emph{第三步: 全局误差.} $X=\Psi(\xi,X)$, 而 $\widehat{X}=\Psi(\xi,\widehat{X}_{\eta(\cdot)})$, 其中 $\widehat{X}_{\eta(\cdot)}$ 表示过程 $s\mapsto\widehat{X}_{\eta(s)}$ (它右连续且适应). 由基本估计引理~\ref{lem:sde-estimate}, $e(t)=\E\sup_{u\le t}\abs{X_u-\widehat{X}_u}^2$ 满足
\begin{align*}
e(t)&\le L\int_0^t\E\abs{X_s-\widehat{X}_{\eta(s)}}^2\dd s\le2L\int_0^t\Big(\E\abs{X_s-\widehat{X}_s}^2+\E\abs{\widehat{X}_s-\widehat{X}_{\eta(s)}}^2\Big)\dd s\\
&\le2L\int_0^te(s)\dd s+2LTC_1h.
\end{align*}
由 Gronwall 不等式, $e(T)\le2LTC_1e^{2LT}h$.
\end{proof}

所以 Euler-Maruyama 格式的\textbf{强收敛阶}为 $\frac{1}{2}$: 误差 $(\E\abs{X_T-\widehat{X}_T}^2)^{1/2}=O(h^{1/2})$. 这比常微分方程的 Euler 方法 (一阶) 差, 原因是布朗增量的量级为 $h^{1/2}$ 而非 $h$. 在 It\^o-Taylor 展开中, 被略去的最低阶项是 $\frac{1}{2}\sigma\sigma'\big((\Delta B)^2-h\big)$, 其量级为 $O(h)$, 高于 $O(h^2)$. 把这一项加入格式, 就得到 \textbf{Milstein 格式}
\[
\widehat{X}_{t_{k+1}}=\widehat{X}_{t_k}+b\,h+\sigma\,\Delta B_k+\frac{1}{2}\sigma\sigma'\big((\Delta B_k)^2-h\big)\qquad(\text{系数都在 }\widehat{X}_{t_k}\text{ 处取值}),
\]
它的强收敛阶为 $1$ (见~\cite{kloeden} 第 10 章). 若 $\sigma$ 是常数 (加性噪声), 二者相同, Euler-Maruyama 格式本身就有一阶强收敛. 图~\ref{fig:em-error} 用几何布朗运动 (有显式解, 可以精确计算误差) 验证了这两个收敛阶.

\begin{figure}[htbp]
\centering
\begin{tikzpicture}
\begin{axis}[SAaxis, axis x line=bottom, axis y line=left, width=6.6cm, height=5.6cm, xmode=log, ymode=log, xmin=0.003, xmax=0.65, ymin=0.0015, ymax=0.6,
  xlabel={步长 $h$}, ylabel={$\E\abs{X_1-\widehat{X}_1}$}, log basis x=2,
  legend style={font=\scriptsize,draw=none,fill=none,at={(0.5,-0.24)},anchor=north}, legend columns=2, legend cell align=left]
  \addplot[SABlue,thick,mark=*,mark size=1.4pt] coordinates {(0.5,0.35689) (0.25,0.27216) (0.125,0.20235) (0.0625,0.14587) (0.03125,0.10395) (0.015625,0.07426) (0.0078125,0.05224) (0.00390625,0.03759)};
  \addplot[SARed,thick,mark=square*,mark size=1.4pt] coordinates {(0.5,0.21145) (0.25,0.12737) (0.125,0.07035) (0.0625,0.03726) (0.03125,0.01925) (0.015625,0.00973) (0.0078125,0.00486) (0.00390625,0.00244)};
  \addplot[SABlue!50,dashed,domain=0.0039:0.5] {0.25*x^0.5};
  \addplot[SARed!50,dashed,domain=0.0039:0.5] {0.2*x};
  \legend{Euler-Maruyama, Milstein, 斜率 $\frac12$, 斜率 $1$}
\end{axis}
\end{tikzpicture}\hfill
\begin{tikzpicture}
\begin{axis}[SAaxis, axis x line=bottom, axis y line=left, width=8cm, height=5.6cm, xmin=0, xmax=1.03, ymin=0.7, ymax=2.6,
  xtick={0,0.25,0.5,0.75,1}, ytick={1,1.5,2,2.5}, xlabel={$t$},
  legend style={font=\scriptsize,draw=none,fill=none,at={(0.5,-0.24)},anchor=north}, legend columns=2, legend cell align=left]
  \addplot[SADark,thin] coordinates {(0.000,1.000) (0.004,1.039) (0.008,0.986) (0.012,0.998) (0.016,1.039) (0.020,0.988) (0.023,0.963) (0.027,0.934) (0.031,0.901) (0.035,0.917) (0.039,0.895) (0.043,0.902) (0.047,0.864) (0.051,0.902) (0.055,0.891) (0.059,0.889) (0.062,0.952) (0.066,0.965) (0.070,0.963) (0.074,0.999) (0.078,0.889) (0.082,0.939) (0.086,0.922) (0.090,0.815) (0.094,0.830) (0.098,0.869) (0.102,0.906) (0.105,0.933) (0.109,0.938) (0.113,0.975) (0.117,0.927) (0.121,1.020) (0.125,1.127) (0.129,1.213) (0.133,1.229) (0.137,1.158) (0.141,1.190) (0.145,1.159) (0.148,1.178) (0.152,1.186) (0.156,1.174) (0.160,1.175) (0.164,1.139) (0.168,1.150) (0.172,1.254) (0.176,1.312) (0.180,1.341) (0.184,1.327) (0.188,1.291) (0.191,1.223) (0.195,1.316) (0.199,1.312) (0.203,1.282) (0.207,1.226) (0.211,1.284) (0.215,1.328) (0.219,1.484) (0.223,1.412) (0.227,1.382) (0.230,1.397) (0.234,1.348) (0.238,1.440) (0.242,1.489) (0.246,1.554) (0.250,1.657) (0.254,1.797) (0.258,1.813) (0.262,1.900) (0.266,1.844) (0.270,1.863) (0.273,1.900) (0.277,1.997) (0.281,2.010) (0.285,2.012) (0.289,2.297) (0.293,2.265) (0.297,2.203) (0.301,2.283) (0.305,2.162) (0.309,2.101) (0.312,2.001) (0.316,2.039) (0.320,2.039) (0.324,2.108) (0.328,2.022) (0.332,2.019) (0.336,2.097) (0.340,2.229) (0.344,2.386) (0.348,2.417) (0.352,2.329) (0.355,2.283) (0.359,2.295) (0.363,2.150) (0.367,2.019) (0.371,1.879) (0.375,1.845) (0.379,1.715) (0.383,1.545) (0.387,1.539) (0.391,1.509) (0.395,1.457) (0.398,1.426) (0.402,1.506) (0.406,1.473) (0.410,1.503) (0.414,1.434) (0.418,1.392) (0.422,1.384) (0.426,1.495) (0.430,1.440) (0.434,1.324) (0.438,1.241) (0.441,1.154) (0.445,1.027) (0.449,0.994) (0.453,1.019) (0.457,1.064) (0.461,1.059) (0.465,1.151) (0.469,1.243) (0.473,1.158) (0.477,1.213) (0.480,1.189) (0.484,1.127) (0.488,1.144) (0.492,1.092) (0.496,1.188) (0.500,1.215) (0.504,1.191) (0.508,1.357) (0.512,1.373) (0.516,1.374) (0.520,1.321) (0.523,1.334) (0.527,1.274) (0.531,1.211) (0.535,1.179) (0.539,1.190) (0.543,1.188) (0.547,1.198) (0.551,1.204) (0.555,1.191) (0.559,1.201) (0.562,1.276) (0.566,1.274) (0.570,1.219) (0.574,1.339) (0.578,1.351) (0.582,1.300) (0.586,1.316) (0.590,1.405) (0.594,1.434) (0.598,1.415) (0.602,1.402) (0.605,1.348) (0.609,1.358) (0.613,1.430) (0.617,1.443) (0.621,1.391) (0.625,1.315) (0.629,1.270) (0.633,1.189) (0.637,1.229) (0.641,1.171) (0.645,1.259) (0.648,1.285) (0.652,1.331) (0.656,1.203) (0.660,1.175) (0.664,1.102) (0.668,1.192) (0.672,1.219) (0.676,1.197) (0.680,1.241) (0.684,1.217) (0.688,1.090) (0.691,1.159) (0.695,1.266) (0.699,1.304) (0.703,1.290) (0.707,1.221) (0.711,1.204) (0.715,1.249) (0.719,1.274) (0.723,1.230) (0.727,1.235) (0.730,1.237) (0.734,1.146) (0.738,1.173) (0.742,1.174) (0.746,1.212) (0.750,1.157) (0.754,1.139) (0.758,1.194) (0.762,1.226) (0.766,1.318) (0.770,1.257) (0.773,1.242) (0.777,1.179) (0.781,1.150) (0.785,1.122) (0.789,1.115) (0.793,1.082) (0.797,1.129) (0.801,1.079) (0.805,1.159) (0.809,1.153) (0.812,1.179) (0.816,1.168) (0.820,1.150) (0.824,1.164) (0.828,1.255) (0.832,1.266) (0.836,1.313) (0.840,1.271) (0.844,1.233) (0.848,1.239) (0.852,1.235) (0.855,1.109) (0.859,1.123) (0.863,1.074) (0.867,1.082) (0.871,1.000) (0.875,0.925) (0.879,0.876) (0.883,0.965) (0.887,0.889) (0.891,0.871) (0.895,0.948) (0.898,0.905) (0.902,0.954) (0.906,0.926) (0.910,0.930) (0.914,0.902) (0.918,0.948) (0.922,1.011) (0.926,0.932) (0.930,0.908) (0.934,0.958) (0.938,0.938) (0.941,0.999) (0.945,0.957) (0.949,0.973) (0.953,0.938) (0.957,0.928) (0.961,0.993) (0.965,1.062) (0.969,0.984) (0.973,0.989) (0.977,0.933) (0.980,0.897) (0.984,0.904) (0.988,0.993) (0.992,0.963) (0.996,0.986) (1.000,0.963)};
  \addplot[SARed,thick,mark=*,mark size=1.3pt] coordinates {(0.000,1.000) (0.125,1.160) (0.250,1.653) (0.375,1.897) (0.500,1.181) (0.625,1.321) (0.750,1.204) (0.875,0.983) (1.000,1.062)};
  \legend{精确解, Euler-Maruyama ($h=\frac18$)}
\end{axis}
\end{tikzpicture}
\caption{左: 几何布朗运动 $\dd X=\frac12X\dd t+0.8X\dd B$ 在 $T=1$ 处的强误差 (双对数坐标, $2\times10^4$ 条轨道), Euler-Maruyama 格式与 Milstein 格式的误差分别以 $h^{1/2}$ 与 $h$ 的速度下降. 右: 同一条布朗轨道驱动下的精确解与粗步长的 Euler-Maruyama 近似}
\label{fig:em-error}
\end{figure}

% ---------------------------------------------------------------------
\subsection{Monte Carlo 方法}

为计算 $\E f(X_T)$ (例如 Feynman-Kac 表示中的期望, 或期权价格), 用 Euler-Maruyama 格式独立地模拟 $n$ 条轨道, 取样本平均 $\frac{1}{n}\sum_{j=1}^nf(\widehat{X}^{(j)}_T)$. 其误差由两部分组成:
\[
\frac{1}{n}\sum_{j=1}^nf\big(\widehat{X}^{(j)}_T\big)-\E f(X_T)=\underbrace{\Big(\frac{1}{n}\sum_{j=1}^nf\big(\widehat{X}^{(j)}_T\big)-\E f\big(\widehat{X}_T\big)\Big)}_{\text{统计误差}}+\underbrace{\Big(\E f\big(\widehat{X}_T\big)-\E f(X_T)\Big)}_{\text{离散化偏差 (弱误差)}}.
\]
由中心极限定理, 统计误差的量级为 $\frac{\sigma_f}{\sqrt{n}}$ (第~1~章), 可以用样本标准差估计并给出置信区间. 对光滑的 $f$, 弱误差是 $O(h)$, 比强误差好 (\textbf{弱收敛阶}为 $1$, 见~\cite{kloeden} 第 14 章): 它只要求分布接近, 无需轨道接近. 要使总误差为 $\varepsilon$, 需要 $h\sim\varepsilon$, $n\sim\varepsilon^{-2}$, 计算量为 $n/h\sim\varepsilon^{-3}$. Giles 的\textbf{多层 Monte Carlo} 方法把不同步长的模拟组合起来, 可以把计算量降到 $\varepsilon^{-2}(\log\varepsilon)^2$ (见~\cite{glasserman}).

对依赖于整条轨道的泛函, 离散化偏差可能更大. 以 6.4.4 小节的向上敲出看涨期权为例: 若只在格点上检查股价是否超过障碍, 就会漏掉格点之间的越界, 从而高估期权价格. 图~\ref{fig:barrier-mc} 显示, 这种偏差只以 $O(\sqrt{h})$ 的速度消失: 在步长为 $h$ 的两个格点之间, 布朗轨道的振幅是 $\sqrt{h}$ 的量级.

\begin{figure}[htbp]
\centering
\begin{tikzpicture}
\begin{axis}[SAaxis, axis x line=bottom, axis y line=left, width=11cm, height=5.4cm, xmode=log, log basis x=2, xmin=3, xmax=700, ymin=1.3, ymax=3.2,
  xtick={4,16,64,256}, xticklabels={$4$,$16$,$64$,$256$}, ytick={1.5,2,2.5,3}, xlabel={$N$},
  legend style={font=\small,draw=none,fill=none,at={(0.98,0.98)},anchor=north east}, legend cell align=left]
  \addplot[SABlue,thick,mark=*,mark size=1.5pt,error bars/.cd,y dir=both,y explicit] coordinates {(4,3.0184) +- (0,0.0203) (8,2.6183) +- (0,0.0187) (16,2.2899) +- (0,0.0174) (32,2.0749) +- (0,0.0164) (64,1.9025) +- (0,0.0156) (128,1.7839) +- (0,0.0150) (256,1.6978) +- (0,0.0146) (512,1.6394) +- (0,0.0143)};
  \addplot[SARed,thick,dashed] coordinates {(3,1.503) (700,1.503)};
  \legend{离散监测的 Monte Carlo 估计, 连续监测的精确价格}
\end{axis}
\end{tikzpicture}
\caption{向上敲出看涨期权 ($S_0=K=100$, $H=130$, 其余参数同图~\ref{fig:barrier}) 的 Monte Carlo 价格 (横轴为每年的监测次数 $N=1/h$, $4\times10^5$ 条轨道, 误差棒为 $95\%$ 置信区间). 只在 $N$ 个等距时刻检查障碍时, 价格随 $N$ 增大而缓慢地下降到连续监测的价格 (定理~\ref{thm:barrier}), 偏差约与 $N^{-1/2}$ 成正比}
\label{fig:barrier-mc}
\end{figure}

\begin{remark}[布朗桥修正]
反射原理给出了消除这一偏差的办法. 由推论~\ref{cor:max-joint}, 对从 $x$ 出发、在时刻 $\Delta$ 到达 $y$ 的布朗运动 ($x,y<h$), 中间曾经到达 $h$ 的条件概率为
\[
\P\Big(\max_{u\le\Delta}W_u\ge h\,\Big|\,W_0=x,W_\Delta=y\Big)=\frac{\varphi_\Delta(2h-x-y)}{\varphi_\Delta(y-x)}=\exp\Big(-\frac{2(h-x)(h-y)}{\Delta}\Big)
\]
(先对 $x=0$ 计算: 联合密度与 $W_\Delta$ 的密度之比为 $\exp\big(-\frac{(2h-y)^2-y^2}{2\Delta}\big)=\exp\big(-\frac{2h(h-y)}{\Delta}\big)$, 一般情形平移即可). 漂移不改变这个条件概率 (给定两端点后, 漂移项被吸收), 扩散系数为 $\sigma$ 时把 $\Delta$ 换成 $\sigma^2\Delta$. 模拟时, 在每一步以这个概率判定 ``格点之间越界'', 即可得到连续监测的无偏估计. 用这种方法, 即使每年只监测 $50$ 次, $2\times10^6$ 条轨道给出的估计 $1.5038\pm0.0060$ 也与精确值 $1.503$ 吻合.
\end{remark}

% ---------------------------------------------------------------------
\subsection*{本章小结}
\addcontentsline{toc}{subsection}{本章小结}

\begin{center}
\begin{minipage}{\textwidth}
\centering
\captionof{table}{随机微分方程与扩散过程的主要结论}
\label{tab:ch6-summary}
\renewcommand{\arraystretch}{1.5}
\begin{tabular}{>{\raggedright\arraybackslash}p{0.24\textwidth} >{\raggedright\arraybackslash}p{0.46\textwidth} >{\raggedright\arraybackslash}p{0.20\textwidth}}
\toprule
结论 & 内容 & 用到的工具 \\
\midrule
线性方程 & 常数变易法给出显式解 (GBM, OU, 布朗桥) & It\^o 公式, 乘积法则 \\
存在唯一性 & Lipschitz + 线性增长 $\Rightarrow$ 唯一强解 & Picard 迭代, Doob, Gronwall \\
弱解 & Girsanov 构造; Tanaka 方程无强解 & Girsanov, L\'evy 刻画, Tanaka \\
马尔可夫性 & $\E[f(X_{s+h})\,|\,\mathcal{F}_s]=P_hf(X_s)$ & 唯一性, 代入引理, 单调类 \\
生成元 & $L=b\partial_x+\frac12\sigma^2\partial_{xx}$, Dynkin 公式 & It\^o 公式, 可选停时 \\
Kolmogorov 方程 & $\partial_tu=Lu$, $\partial_tp=L^*p$, 平稳分布 & It\^o 公式, 分部积分 \\
出口问题 & 尺度函数 $Ls=0$, 期望出口时间 $Lv=-1$ & Dynkin 公式 \\
Dirichlet 问题 & $u(\rvec{x})=\E_{\rvec{x}}g(\rvec{W}_{\tau_D})$ & 多维 It\^o 公式 \\
Feynman-Kac & 带贴现与源项的抛物方程的概率表示 & It\^o 公式, 矩估计 \\
期权定价 & Black-Scholes 公式, 障碍期权, 美式期权 & Girsanov, 鞅表示, 反射原理 \\
Euler-Maruyama & 强收敛阶 $\frac12$, 弱收敛阶 $1$ & 基本估计, Gronwall \\
\bottomrule
\end{tabular}
\end{minipage}
\end{center}

\subsection*{结语}
\addcontentsline{toc}{subsection}{结语}

回顾全书, 第~1~章的测度论提供了严格的语言, 第~2~章用条件期望刻画信息; 第~3~章的鞅论描述了 ``公平'' 的随机演化并提供了最有力的不等式与收敛定理, 第~4~章的布朗运动则是连续随机性的普适模型. 在此基础上, It\^o 积分与 It\^o 公式建立了随机微积分 (第~5~章), 随机微分方程把它们用于描述和计算一般的连续随机系统 (第~6~章). 进一步的学习可以沿几个方向展开: 一般半鞅与带跳过程的随机积分 (\cite{revuz}), 随机控制与最优停时, 滤波理论, Malliavin 分析, 随机偏微分方程, 以及更深入的金融数学 (\cite{shreve}). 扩散过程与偏微分方程之间的联系 (\cite{karatzas}, \cite{oksendal}) 则贯穿了所有这些方向.
